% !TeX program = xelatex
% Standalone English volume. Use XeLaTeX until the references stabilize.
\documentclass[a4paper,11pt]{article}
\usepackage[T1]{fontenc}
\usepackage{lmodern}
\usepackage[english]{babel}
\usepackage[left=26mm,right=26mm,top=25mm,bottom=25mm,headheight=15pt,headsep=7mm,footskip=11mm]{geometry}
\usepackage{amsmath,amssymb,mathtools,amsthm}
\newtheorem{proposition}{Proposition}[section]
\newtheorem{lemma}[proposition]{Lemma}
\newtheorem{theorem}[proposition]{Theorem}
\newtheorem{corollary}[proposition]{Corollary}
\newtheorem{maintheorem}[proposition]{Main Theorem}
\theoremstyle{definition}
\newtheorem{definition}[proposition]{Definition}
\theoremstyle{remark}
\newtheorem{remark}[proposition]{Remark}
\makeatletter
\renewcommand*\l@section{\@dottedtocline{1}{0em}{2.4em}}
\renewcommand*\l@subsection{\@dottedtocline{2}{1.45em}{3.45em}}
\renewcommand*\l@subsubsection{\@dottedtocline{3}{3.8em}{4.55em}}
\makeatother
\usepackage{xcolor,fancyhdr,titlesec,needspace,xurl}
\usepackage{longtable,multicol}
\usepackage{hyperref}
\definecolor{ink}{HTML}{1A3542}
\definecolor{muted}{HTML}{4F606A}
\hypersetup{unicode=true,colorlinks=true,linkcolor=ink,filecolor=ink,urlcolor=ink,citecolor=ink,bookmarksnumbered=true,
 pdfdisplaydoctitle=true,pdflang={en},
 pdftitle={The generalized Riemann hypothesis for Dirichlet L-functions. Edition 930, Volume XV},
 pdfsubject={A natural-language reconstruction of the proposed proof},
 pdfauthor={A mathematical reconstruction}}
\titleformat{\section}{\Large\bfseries\color{ink}}{\thesection}{.65em}{}
\titlespacing*{\section}{0pt}{.8em}{.65em}
\setlength{\parindent}{1.2em}
\linespread{1.045}
\setlength{\parskip}{.22em plus .05em minus .03em}
\setlength{\emergencystretch}{2em}
\setcounter{topnumber}{2}
\setcounter{tocdepth}{3}\setcounter{secnumdepth}{3}
\widowpenalty=4500
\clubpenalty=4500
\pagestyle{fancy}\fancyhf{}
\fancyhead[L]{\small\color{muted}Dirichlet GRH \textperiodcentered\ Volume I}
\fancyhead[R]{\small\color{muted}Dirichlet zeros and power intervals}
\fancyfoot[C]{\small\thepage}
\renewcommand{\headrulewidth}{.3pt}
\newcommand{\lead}[1]{\par\smallskip\noindent{\bfseries\color{ink}#1.}\enspace}
\newcommand{\sour}[2]{\unskip\penalty150\hspace{.35em}{\small\hyperlink{ref#1}{[#2]}}}
\titleformat{\subsection}{\large\bfseries}{\thesubsection}{.6em}{}
\titlespacing*{\subsection}{0pt}{1.15em plus .15em}{.4em}
\setlength{\abovedisplayskip}{9pt plus 2pt minus 3pt}
\setlength{\belowdisplayskip}{9pt plus 2pt minus 3pt}
\newcommand{\RR}{\mathbb R}\newcommand{\CC}{\mathbb C}\newcommand{\QQ}{\mathbb Q}
\newcommand{\ZZ}{\mathbb Z}\newcommand{\ord}{\operatorname{ord}}
\newcommand{\Frob}{\operatorname{Frob}}\newcommand{\Gal}{\operatorname{Gal}}
\newcommand{\len}{\operatorname{length}}\newcommand{\Res}{\operatorname*{Res}}
\newcommand{\wmean}{\operatorname{av}}\newcommand{\Ker}{\operatorname{ker}}
\newcommand{\Lh}{\mathcal L}\newcommand{\Gg}{\mathcal G}
\newcommand{\real}{\operatorname{Re}}\newcommand{\imag}{\operatorname{Im}}
\newcommand{\abs}[1]{\left|#1\right|}\newcommand{\eps}{\varepsilon}
\newcommand{\anchoritem}[2]{\par\Needspace{2.8\baselineskip}\hypertarget{ref#1}{}\noindent{\bfseries\color{ink}#2.}\enspace}
\newcommand{\srcloc}[2]{{\small\nolinkurl{#1.lean}}, line~#2}
\numberwithin{equation}{section}
\raggedbottom

\newcommand{\Sel}{\mathcal S}
\newcommand{\Cl}{\operatorname{Cl}}
\newcommand{\Aut}{\operatorname{Aut}}
\newcommand{\GL}{\operatorname{GL}}
\newcommand{\OO}{\mathcal O}
\hypersetup{bookmarksdepth=3}


\addto\captionsenglish{\renewcommand{\partname}{Volume}}


\makeatletter
\expandafter\def\csname r@cor:fieldrelations712\endcsname{{12.4}{32}{}{proposition.12.4}{Dirichlet-GRH-v930-Volume-II.pdf}}
\expandafter\def\csname r@cor:gluedperiod811\endcsname{{73.25}{97}{}{proposition.73.25}{Dirichlet-GRH-v930-Volume-XIV.pdf}}
\expandafter\def\csname r@lem:basis59\endcsname{{65.8}{27}{}{proposition.65.8}{Dirichlet-GRH-v930-Volume-XIII.pdf}}
\expandafter\def\csname r@lem:dualcartan811\endcsname{{73.26}{97}{}{proposition.73.26}{Dirichlet-GRH-v930-Volume-XIV.pdf}}
\expandafter\def\csname r@lem:mass32\endcsname{{121.17}{48}{}{proposition.121.17}{Dirichlet-GRH-v930-Volume-XXIII.pdf}}
\expandafter\def\csname r@lem:oscillatorcartan710\endcsname{{142.12}{45}{}{proposition.142.12}{Dirichlet-GRH-v930-Volume-XXVII.pdf}}
\expandafter\def\csname r@lem:split61\endcsname{{65.12}{32}{}{proposition.65.12}{Dirichlet-GRH-v930-Volume-XIII.pdf}}
\expandafter\def\csname r@lem:vacuumtail709\endcsname{{142.8}{42}{}{proposition.142.8}{Dirichlet-GRH-v930-Volume-XXVII.pdf}}
\expandafter\def\csname r@lem:zak242\endcsname{{67.7}{80}{}{proposition.67.7}{Dirichlet-GRH-v930-Volume-XIII.pdf}}
\expandafter\def\csname r@prop:doubledkernel711\endcsname{{142.14}{47}{}{proposition.142.14}{Dirichlet-GRH-v930-Volume-XXVII.pdf}}
\expandafter\def\csname r@sec:integralgluing811\endcsname{{73.10}{96}{}{subsection.73.10}{Dirichlet-GRH-v930-Volume-XIV.pdf}}
\expandafter\def\csname r@seriesguide\endcsname{{1.1}{4}{}{subsection.1.1}{Dirichlet-GRH-v930-Volume-I.pdf}}
\expandafter\def\csname r@thm:class61\endcsname{{65.11}{31}{}{proposition.65.11}{Dirichlet-GRH-v930-Volume-XIII.pdf}}
\expandafter\def\csname r@thm:elliptic312\endcsname{{67.12}{86}{}{proposition.67.12}{Dirichlet-GRH-v930-Volume-XIII.pdf}}
\expandafter\def\csname r@thm:gluedcost811\endcsname{{73.27}{98}{}{proposition.73.27}{Dirichlet-GRH-v930-Volume-XIV.pdf}}
\expandafter\def\csname r@thm:inert175\endcsname{{72.16}{59}{}{proposition.72.16}{Dirichlet-GRH-v930-Volume-XIV.pdf}}
\expandafter\def\csname r@thm:negativeinterval740\endcsname{{15.3}{4}{}{proposition.15.3}{Dirichlet-GRH-v930-Volume-III.pdf}}
\expandafter\def\csname r@thm:one61\endcsname{{65.13}{33}{}{proposition.65.13}{Dirichlet-GRH-v930-Volume-XIII.pdf}}
\expandafter\def\csname r@thm:torsion312\endcsname{{66.35}{74}{}{proposition.66.35}{Dirichlet-GRH-v930-Volume-XIII.pdf}}
\expandafter\def\csname r@thm:transport59\endcsname{{65.9}{28}{}{proposition.65.9}{Dirichlet-GRH-v930-Volume-XIII.pdf}}
\expandafter\def\csname r@thm:vac242\endcsname{{67.8}{81}{}{proposition.67.8}{Dirichlet-GRH-v930-Volume-XIII.pdf}}
\makeatother
\begin{document}
\setcounter{part}{14}
\fancyhead[L]{\small\color{muted}Dirichlet GRH \textperiodcentered\ Volume XV}
\fancyhead[R]{\small\color{muted}Definite kernels and theta exchanges}
\thispagestyle{plain}\begin{center}
{\LARGE\bfseries The generalized Riemann hypothesis\\[2pt]for Dirichlet \(L\)-functions}\\[6pt]
{\small 27 November 2026 \quad\textperiodcentered\quad Edition 930}\end{center}
\part{Definite theta kernels and the three global exchanges}\label{part:XV}
\begin{abstract}\noindent We construct the definite class kernels and their periods. For each of the three theta exchanges we compare the local test factors, the relevant automorphic lines and the normalization of the global integral.\end{abstract}
\setcounter{section}{76}
\noindent The \hyperref[seriesguide]{guide in Volume~I} lists all twenty-seven volumes.
\begingroup\small\setlength{\parskip}{0pt}\tableofcontents\endgroup
\clearpage

\section{Definite periods and the carrier integral}\label{nav:carrier-integral}\label{nav:volume-XV}

\subsection{The periods on the complementary definite unitary group}
\label{sec:complementperiod812}
We determine which functions on \(U(W_j)\) are detected by the
second factor of Corollary~\ref{cor:gluedperiod811}. Fix a node
of the same interval and \(1\le j\le m\). Put
\[
 G=U(W_j),\qquad H=U(V^\circ),\qquad
 a=k+j-1,\quad b=m-j,\quad N=\dim_E V^\circ=5s-9m-1 .
\]
Thus \(a+b=s\), \(G(F_\infty)=\prod_{\tau}U(a,b)\) and
\(H(F_\infty)=\prod_\tau U(N)\). Write \([G]=G(F)\backslash
G(\mathbb A_F)\), and likewise for \(H\). Let \(J\) be the
product of the source principal subgroups modulo \(\pi_v^2\)
at the old marks and the source lattice stabilizers elsewhere.
Define \(J_H\) by the same rule on the target. Normalize each
finite lattice stabilizer to have mass one. We use the bilinear
lift \(\Theta u(y)=\int_{[G]}u(g)\theta(g,y)\,dg\), and take
Hermitian inner products linear in their second argument.

In the negative source columns order the coordinates by the
numbering used for the minors in Section~\ref{sec:integralgluing811}.
At \(\tau\) write \(c_{h,\tau}\) for the number of negative
columns of height \(h\) there. Only \(h\le b\) can occur.
Indeed, in the notation of Theorem~\ref{thm:gluedcost811},
when \(b<h<s\) one has \(\sigma_{s-h}\ge H_0-w_j\), where
\(H_0=s^8\), and the other endpoint is at most
\(-w_j-(t-1)/2<H_0-w_j\). At heights \(h\ge s\) the count is
zero. After choosing an ordering of the target rows take the
Fock polynomial \(Q_\tau=\prod_h\Delta_h(Z_-)^{c_{h,\tau}}\),
where \(\Delta_h\) is the leading \(h\)-rowed square minor and
\(Z_-\) consists of the \(N\)-by-\(b\) negative variables.
It is nonzero at a matrix with an identity in these columns.
Write \(v\) for the tensor of the corresponding Fock vectors,
normalized to length one. Take the source and target circle
characters of this pair to be respectively \(\upsilon^{m-j}\)
and \(\upsilon^{m/2-j}\), with the rank normalization of the
Weil action used above. They have trivial tame part. At each
old mark take the indicator of \(X\equiv z\bmod\pi_v^2\)
on the integral \(N\)-by-\(s\) matrices, and elsewhere take
the unit lattice vacuum. These are precisely the data of
the complementary factor for both choices of twist.

\begin{lemma}[The real projection selected by the negative columns]
\leavevmode\label{lem:negativecoefficient812}
Put \(\delta=j-m/2\). Write \(g_\tau\in U(a,b)\) in blocks
\(\left(\begin{smallmatrix}A&B\\ C&D\end{smallmatrix}\right)\)
in positive and negative coordinates, and set
\(S_\tau=D-CA^{-1}B=(D^*)^{-1}\).
The normalized oscillator coefficient on the source is
\[
 K_\infty(g)=\langle v,\omega(g,1)v\rangle
 =\prod_\tau (\det g_\tau)^\delta(\det S_\tau)^N
           \prod_{h=1}^b\Delta_h(S_\tau)^{c_{h,\tau}} .
\]
An empty determinant is one. This continuous function is in
\(L^1(G(F_\infty))\cap L^2(G(F_\infty))\). Put
\(d_\infty=\int|K_\infty(h)|^2\,dh>0\). Then
\[
 P_\infty u(x)=d_\infty^{-1}\int_{G(F_\infty)}
                       K_\infty(h)u(xh)\,dh
\]
is an orthogonal projection on \(L^2([G])\).
\end{lemma}
\begin{proof}
Here \(A,D\) are invertible by the defining unitary relations.
Block inversion of \(g_\tau\), using
\(g_\tau^{-1}=\operatorname{diag}(1,-1)g_\tau^*
\operatorname{diag}(1,-1)\), gives the formula for \(S_\tau\).
For \(k=\operatorname{diag}(k_+,k_-)\) the Fock action on the
negative variables substitutes \(Z_-\mapsto Z_- k_-\). Its
scalar on the vacuum is \((\det k)^\delta(\det k_-)^N\).
In fact the source acts on a Hom row by the inverse source
matrix. The inverse determinant of its negative block in the
Fock vacuum gives one \(\det k_-\) per target row. The source
action also includes the inverse power \(m/2-j\) of the
determinant from the target circle character.

In a Cartan decomposition insert one real boost of size \(u_i\)
in each paired positive and negative coordinate. In a single
target row its lowering operator differentiates once in each
coordinate of the pair, so it kills a polynomial in negative
letters. Its creation operator introduces a positive letter
in every nonconstant term, orthogonal to any polynomial in
negative letters. The middle, diagonal oscillator operator
gives \((\cosh u_i)^{-1}\) and replaces the paired negative
variable by \((\cosh u_i)^{-1}\) times itself, as in
Lemma~\ref{lem:oscillatorcartan710}. After the two compact
rotations the resulting substitution is \(Z_-\mapsto Z_- S_\tau\).
Indeed on the boost \(S_\tau\) has diagonal entries \(\operatorname{sech}
u_i\) and ones, and under left and right compact rotations it
acquires the corresponding negative blocks on the two sides.
The scalar transforms in the same way as
\((\det g_\tau)^\delta(\det S_\tau)^N\).
For any matrix \(L\) on these variables the normalized
coefficient of \(Q_\tau(Z_- L)\) against \(Q_\tau(Z_-)\) is
\(\prod_h\Delta_h(L)^{c_{h,\tau}}\). To see this, factor a
matrix with nonzero leading minors as lower unitriangular,
diagonal, and upper unitriangular. The upper unitriangular
factor fixes the leading-column polynomial. The lower one
fixes its dual functional, by unitary adjunction and the same
upper triangular calculation. The diagonal entries give the
stated product by telescoping. Polynomial continuation gives
it for all \(L\). This proves the formula.

All singular values of \(S_\tau\) are at most one. Its
determinant has modulus \(\prod_i(\cosh u_i)^{-1}\), so the
coefficient has modulus at most \(\prod_i(\cosh u_i)^{-N}\)
at this place. The Cartan density is bounded by a constant
times \(\exp(2(s-1)\sum_i u_i)\), by the same lemma.
Here \(N>2s\) at sufficiently large nodes. Integrability
follows, also when the group is compact. Its square is
integrable and its value at the identity is one.

The vector is a joint highest weight harmonic. Indeed every
mixed contraction differentiates in a positive letter and
thus gives zero, while triangular changes of the ordered
rows and columns give the two highest lines of the minors.
In the harmonic decomposition for the compact dual member
\(U(N)\), its fixed target highest line generates an
irreducible source representation. This is the compact
unitary dual-pair decomposition of
\href{https://msp.org/pjm/1987/129-1/pjm-v129-n1-p02-s.pdf}
{Davidson, \S2, Theorem~2.7}. Tensoring with the determinant
characters preserves irreducibility. For \(b=0\) the vector
is the vacuum and its source representation is the indicated
character. The tensor over real places is irreducible and
has a square-integrable coefficient. Schur orthogonality
therefore gives \(K_\infty*K_\infty=d_\infty K_\infty\).
Also \(K_\infty(h^{-1})=\overline{K_\infty(h)}\). Integrating
against the unitary right action proves the projection claim,
with the sign of the kernel fixed by the bilinear lift.
\end{proof}

\begin{lemma}[The finite part at the complementary rank]
\leavevmode\label{lem:complementfinite812}
Let \(f=[F:\QQ]\). On \(L^2([G])^J\) convolution with the
finite oscillator coefficient is \(c_{\mathcal M}U\), where
\[
 c_{\mathcal M}=\prod_{v\in\mathcal M}
       \frac{q_v^{-2Ns}}{|\GL_s(\mathcal O_v/\pi_v^2)|},
 \qquad
 U=U^*,\qquad \|U-I\|\ll f\,2^{-s/3}.
\]
The constant is absolute for all sufficiently large nodes.
In particular \(U\) is positive and invertible there, with
\(\|U^{-1}\|\le2\), and it commutes with \(P_\infty\).
\end{lemma}
\begin{proof}
The \(z\)-rows on the point coset are the identity modulo
\(\pi_v^2\). If both \(X\) and \(Xg\) are in its support,
the two square blocks force \(g\equiv1\bmod\pi_v^2\).
Conversely this subgroup fixes the test. Thus its coefficient
is its squared norm \(q_v^{-2Ns}\) on that subgroup and zero
elsewhere. Integrating gives the corresponding factor in
\(c_{\mathcal M}\).

At every other place \(v\) let \(K_v\) be the lattice
stabilizer and \(A_v\) the vacuum coefficient. It is one on
\(K_v\) and bi-invariant. Write \(L_v\) for the integral
source lattice on the quadratic algebra. The Heisenberg
lattice calculation in Lemma~\ref{lem:vacuumtail709} gives
\(|A_v(g)|\le[L_v:L_v\cap gL_v]^{-N/2}\): the Hom lattice
here consists of \(N\) copies of the source dual. The same
lemma's sublattice count bounds the number of cosets at an
inert place with intersection index \(Q^i\) by
\(16Q^{2(s-1)i}\), where \(Q=q_v^2\). At a split place
the two indices in the chosen component are \(q_v^u,q_v^w\)
and there are at most \(16q_v^{(s-1)(u+w)}\) cosets with
these indices. Its inverse dual supplies the other component.
These are all cases since \(E/F\) is unramified at finite
places by Corollary~\ref{cor:fieldrelations712}.

Now \(N/2-2(s-1)=(s-9m+3)/2\ge s/3\), using \(s\ge x^5\)
and \(m\le2\sqrt x\). Hence the sum of absolute values over
cosets other than \(K_v\) is bounded at inert places by
\(16\sum_{i\ge1}Q^{-si/3}\), and at split places by
\(16\sum_{i\ge1}(i+1)q_v^{-si/3}\).
Both are at most \(128q_v^{-s/3}\). In the space of
\(K_v\)-fixed functions the \(K_v\) term of convolution is
the identity. Each other coset contributes a translate
with operator norm bounded by its absolute coefficient.
The resulting bi-invariant operator is self-adjoint since
\(A_v(g^{-1})=\overline{A_v(g)}\). The operators at
distinct places commute. Moreover
\(\sum_v q_v^{-s/3}\le f\sum_{l=2}^\infty l^{-s/3}
\ll f\,2^{-s/3}\).
Taking the norm limit of the products bounds \(\|U-I\|\)
by \(\exp(128\sum_v q_v^{-s/3})-1\), with the sum away
from the marks. The degree \(f\ll_{q_0,d}x^{12}\) on
these intervals makes this less than \(1/2\) eventually.
The Neumann series and commutation of finite and real
translations give the conclusions.
\end{proof}

\begin{theorem}[The complementary Gram operator on the prescribed level]
\leavevmode\label{thm:complementgram812}
Let \(\Theta\) use the data at the beginning of this subsection,
with the unit vector \(v\) at infinity. There is a positive
constant \(c_H\), fixed by the Haar normalizations and the
two spaces, such that
\[
 \Theta^*\Theta=c_H\,c_{\mathcal M}\,d_\infty\, U P_\infty
                 \quad\text{on }L^2([G])^J .
\]
Here the lift on \(L^2\) is the continuous extension of its
integral on smooth rapidly decreasing automorphic functions
of this level. In particular its squared norm on \(u\), divided
by \(c_Hc_{\mathcal M}d_\infty\), is between
\((1-\|U-I\|)\|P_\infty u\|^2\) and
\((1+\|U-I\|)\|P_\infty u\|^2\).
\end{theorem}
\begin{proof}
Apply the convergent Siegel--Weil formula to the doubled
source \(W_j\oplus(-W_j)\) and the definite space of rank
\(N\). Its Siegel parabolic has Levi \(\GL_s(E)\).
Since the Siegel point \((N-s)/2\) is greater than \(s/2\),
the Eisenstein sum is absolutely convergent and the
definite-group integral is its value, up to the positive
Haar constant. See
\href{https://msp.org/pjm/2013/264-1/pjm-v264-n1-p11-s.pdf}
{Yamana, Theorem~2 and the convergence range preceding it}.
Use the doubled tensor \(\bar\Phi\otimes\Phi\) of the
specified datum. Its theta integral on the doubled group,
restricted to a pair \(g',g\), is the integral over \([H]\)
of \(\overline{\theta(g',y)}\theta(g,y)\). This follows
by using the direct sum polarization and summing over the
two rational position vectors. The conjugate splitting
on the first copy gives precisely the complex conjugate.

We verify the terms of the Eisenstein sum because the
function to which the resulting identity is applied can
be a product of automorphic functions. Each term corresponds
to a rational maximal isotropic plane. Write a basis of it
in blocks \(\binom{K}{K'}\) with \(s\) columns. At one old
split mark write \(X,X'\) for the two position matrices of
the doubled tensor. The Siegel functional for this plane is
supported where \(XK-X'K'=0\), choosing the sign of the
second coordinates accordingly. Indeed, in the polarization
adapted to the plane it is evaluation at zero. In the split
matrix polarization the partial Fourier transform giving
this evaluation integrates over matrices annihilating the
plane. Equivalently the functional is invariant under all
multipliers \(\psi_v(\operatorname{tr}(Z(XK-X'K')))\),
whose characters separate points. In the
translated point tests the matrices still have column rank
\(s\), since right translation at a split place acts by
an invertible change of source coordinates. If \(K'c=0\)
then \(XKc=0\), hence \(Kc=0\) and consequently \(c=0\).
Thus \(K'\) is invertible at this place, and over \(E\).
Total isotropy says the plane is the graph of an isometry.
Every rational isometry gives one plane, so the surviving
sum is indexed by \(G(F)\).

For the identity graph the partial Fourier transform at
zero pairs the two data on \(X=X'\). With self-dual additive
measure this is their oscillator inner product. Equivariance
takes the graph of \(\gamma\in G(F)\) to the coefficient at
\(g'^{-1}\gamma g\), with trivial phase on the rational
splitting. Thus the kernel is \(c_H\) times
the sum of translates of the adelic coefficient
\(A(h)=\langle\Phi,\omega(h,1)\Phi\rangle\).
Lemmas~\ref{lem:negativecoefficient812} and
\ref{lem:complementfinite812}, including the absolute
coset bounds in its proof, give \(A\in L^1(G(\mathbb A_F))\).
Unfolding against \(\overline{u'(g')}u(g)\) therefore gives
\[
 \langle\Theta u',\Theta u\rangle
       =c_H\int_{G(\mathbb A_F)}
                         A(h)\langle u',R(h)u\rangle\,dh ,
 \qquad (R(h)u)(x)=u(xh).
\]
Indeed the unfolded absolute integral is bounded by
\(\|A\|_1\|u'\|_2\|u\|_2\). For smooth rapidly decreasing
inputs the original lift integrals also converge absolutely,
since theta kernels have moderate growth uniformly on the
compact target quotient. Factoring the last integral gives
the stated identity on these inputs. Compactly supported
smooth functions on the source quotient, averaged over
\(J\), form a dense subspace of the level space. The upper
bound gives a bounded extension. The lower bound follows
from self-adjointness of \(U\) and commutation with
\(P_\infty\).
\end{proof}

\begin{corollary}[What the complementary class evaluations span]
\leavevmode\label{cor:complementreproduce812}
Choose representatives \(y\) of
\(H(F)\backslash H(\mathbb A_{F,f})/J_H\). At each such class
test against the polynomials in the span of the compact
target translates of \(\prod_\tau Q_\tau\). These finitely
many classes, with a basis of that polynomial space, give
linear period functionals whose common kernel is exactly
\(\ker P_\infty\) on the level space. Their Petersson dual
vectors span \(P_\infty L^2([G])^J\).

In particular, if \(F_0\) is a good source form and \(B\) is
a first-factor theta function of Corollary~\ref{cor:gluedperiod811}
with the compatible splitting, the complementary periods of
\(F_0B\) detect precisely \(P_\infty(F_0B)\). When \(j=m\)
this is the ordinary compact-group integral of
\(F_0(xh)B(xh)\prod_\tau(\det h_\tau)^{m/2}\), using
probability Haar measure at infinity.
\end{corollary}
\begin{proof}
The finite class set is finite by compactness of \([H]\).
The finite test is right \(J_H\)-invariant and the span of
the compact translates of \(v\) is finite dimensional.
Their polynomial factors belong to the allowed negative-letter
tests because rotating the \(N\) target rows expands each minor into a
linear combination of minors on the same source columns
by Cauchy--Binet. The determinant character on the target
acts by scalars. Rational invariance and right translation
at infinity express every value of \(\Theta u\) as a
combination of these class evaluations. Conversely every
evaluation in this span is a linear combination of values
of the translated lift. Finite stabilizers merely impose
relations between values. The resulting finite-dimensional
space of target functions is closed in \(L^2([H])\) and all
its evaluation maps are continuous. The equivalence thus
holds on the completed source. Theorem~\ref{thm:complementgram812}
identifies the common kernel. Taking orthogonal complements
proves the spanning statement. Equivalently, the image of
\(\Theta^*\Theta\) is exactly the range of \(P_\infty\)
since \(U\) is invertible on that range.

The good form is cuspidal and the first factor has moderate
growth, so their product is rapidly decreasing. At the old
marks it is fixed modulo \(\pi_v^2\). At a new mark the
character of the determinant of the marked rows cancels
the source action on the first finite test, and the second
factor is spherical. The product has full source level there
and at every other unmarked place. It is therefore an
admissible input. For \(j=m\) the negative block is empty
and Lemma~\ref{lem:negativecoefficient812} gives the displayed
character integral.
\end{proof}

For the comparison of individual determinants, the remaining
spanning assertion concerns the first factors. Their prescribed
arithmetic family must give \(P_\infty(F_0B)\ne0\) for at least
one member whenever \(F_0\ne0\). The complex span of the
complementary class evaluations by itself supplies no bound
on the integral index of these product periods.
\clearpage

\subsection{The carrier form and shortest vectors at its peaks}
\label{sec:carrierpeaks813}
The first factors paired with \(P_\infty\) in
Corollary~\ref{cor:complementreproduce812} use a single form on
\(U(2m,m)\). We construct it on each of the fields already in
use and locate the points where it and a good source form
have large absolute value. The ranks here are smaller than
that of the complementary group in the preceding subsection.

\begin{lemma}[Counting self-dual neighbors of an unramified lattice]
\leavevmode\label{lem:selfdualtail813}
Let \(L\) be a self-dual Hermitian lattice of rank \(r\) over
an unramified quadratic extension of a local field, with
residue cardinality \(Q\) in the extension. For \(i\ge1\),
the number of self-dual lattices \(L'\) for which
\([L:L\cap L']=Q^i\) is at most \(8^iQ^{ri}\).
Consequently at target rank \(3r\), for the unramified unit
vacuum on the fields of Corollary~\ref{cor:fieldrelations712},
the sum of absolute source coefficients over cosets outside
the lattice stabilizer is \(O(q^{-r/2})\) at every finite
place of residue cardinality \(q\) in \(F\), with an absolute
constant for sufficiently large \(r\).
\end{lemma}
\begin{proof}
Fix the intersection \(I=L\cap L'\) and the isomorphism type
\(A=L/I\) over the valuation ring \(O\). Duality gives
\(I^\#=L+L'\). In \(I^\#/I\) the modules \(A\) and \(L'/I\)
are complementary. The perfect pairing with values in
\(\operatorname{Frac}(O)/O\) identifies \(L'/I\) with the
conjugate dual of \(A\), of the same elementary divisor type.
If there is one complement, the others are graphs of maps
from it to \(A\). There are thus at most
\(|\operatorname{End}_O(A)|\) choices for \(L'\) for this \(I\).
On the other hand the number of possible \(I\) of this type
is bounded by \(Q^{ri}/|\operatorname{Aut}_O(A)|\), by
counting surjections \(L\to A\).

Write \(A=\bigoplus_{h\ge1}(O/\pi^h)^{e_h}\).
The fraction of its endomorphisms that are invertible equals
\(\prod_{h:e_h>0}\prod_{l=1}^{e_h}(1-Q^{-l})\).
Indeed reduction modulo \(\pi\) gives a block triangular
matrix when one orders the block lengths increasingly.
Its diagonal blocks are arbitrary and independent. By
Nakayama's lemma invertibility is equivalent to the
invertibility of all those blocks. Since
\(\prod_{l\ge1}(1-2^{-l})>1/4\) and at most \(i\) block
lengths occur, each type costs at most \(4^i Q^{ri}\).
The types are partitions of \(i\), of which there are at
most \(2^{i-1}\) by comparison with ordered compositions.

At an inert place the coset determines \(L'\), \(Q=q^2\),
and the vacuum bound of Lemma~\ref{lem:vacuumtail709} is
\(Q^{-3ri/2}\). Thus the tail is bounded by
\(\sum_{i\ge1}(8Q^{-r/2})^i\).
At a split place use the same lemma's count of
\(16(i+1)q^{(r-1)i}\) for total index exponent \(i\) on
the two components, and the coefficient bound
\(q^{-3ri/2}\). Summation proves the assertion, including
at places of residue characteristic two.
\end{proof}

\begin{theorem}[The prescribed full-level cuspidal carrier]
\leavevmode\label{thm:carrier813}
Let \(D^-=-I_m\) be the negative definite standard space,
and \(S=I_{2m}\oplus(-I_m)\). Take the theta lift of the
constant function on \([U(D^-)]\) to \(U(S)\), with circle
characters \((\upsilon^m,1)\), the full finite lattice
vacua, and the archimedean polynomial
\[
          \prod_\tau \det(Z_{\tau,-})^{2m}.
\]
Here \(Z_{\tau,-}\) uses the \(m\) negative target rows and
all \(m\) source columns in the Fock coordinates for this
pair. At every sufficiently large node this lift, denoted
by \(\phi\), is a nonzero smooth cuspidal form fixed by the
full finite lattice stabilizer. It is holomorphic of scalar
weight \((0,\ldots,0;3m,\ldots,3m)\) in the convention that
positive coordinates come first.
\end{theorem}
\begin{proof}
At a real place let \(u\in U(D^-)\). On the Fock vacuum it
acts by \((\det u)^{2m}\). The \(2m\) positive target rows
have negative Hom sign and each contributes one determinant
through the vacuum factor. The target circle character is
trivial. The negative target rows have the opposite sign,
so the indicated determinant polynomial acquires the
character \((\det u)^{-2m}\). The test is nonzero and fixed
by the compact source.

We verify first occurrence to control the constant terms.
The lower members of the rational Witt tower are
\(I_m\oplus\mathbb H^l\), \(0\le l<m\), with \(\mathbb H\)
a hyperbolic plane. Over \(\RR\) their signatures are
\((m+l,l)\). Use the splittings compatible with taking
constant terms in this tower. The compact source vacuum
character in the smaller oscillator is still
\((\det u)^{2m}\). Indeed adjoining one hyperbolic plane
tensors with the oscillator on its isotropic half in the
Schr\"odinger polarization. On that space \(U(D^-)\) acts
by linear substitution with unit Jacobian and fixes the
standard Gaussian. This Gaussian is the vacuum in its
compatible Fock polarization. This fixes the source
character on descent at each step.

Apart from that vacuum character, the polynomial algebra
has \(m+l\) copies of the standard source letters and \(l\)
copies of the dual letters. The latter are the rows whose
determinant had negative weight above. A summand, as a
compact source representation, has the form
\(S_\lambda(\CC^m)\otimes S_\mu(\CC^m)^*\) with
\(\ell(\mu)\le l\). This is the polynomial Cauchy
decomposition on each set of letters, obtained by taking
simultaneous permutation invariants in the tensor powers
of row and column spaces. For this summand to contain
\(\det^{-2m}\) one must have
\(\mu=\lambda+(2m,\ldots,2m)\), by Schur's lemma applied
after moving \(S_\mu\) across the Hom pairing.
This requires \(m\) nonzero rows of \(\mu\), impossible
for \(l<m\). Averaging over the compact source on each
homogeneous degree proves vanishing on the dense polynomial
vectors and hence on the full oscillator. Every smaller
global lift of the trivial source vanishes, already by
averaging at one real place. The source is cuspidal because
it is anisotropic. Rallis's tower property therefore gives
zero for every proper parabolic constant term of the lift
to \(S\). We use the constant-term tower for the same
splittings as in \href{https://arxiv.org/pdf/1207.4709}
{Gan, Qiu and Takeda, \S11.2}.

To verify nonvanishing for the chosen tensor, normalize its
real vector to length one. Its real source coefficient is
identically one. At finite places the unit vacua are fixed
by the hyperspecial lattice stabilizers and have value one
there in their source coefficients. All characters are
unramified in the lattice splitting. By
Lemma~\ref{lem:selfdualtail813} the unmarked convolution on
the constant source equals the identity up to an operator
of norm \(O(f2^{-m/2})\). Indeed multiply the local
operators and use
\(\sum_v q_v^{-m/2}\ll f2^{-m/2}\), as in
Lemma~\ref{lem:complementfinite812}. Here \(f\ll_{q_0,d}x^{12}\)
and \(m\ge\sqrt{x}\).

The doubled Siegel--Weil formula applies at rank \(3m>2m\).
The target has Witt index \(m\), so the difference \(3m-m\)
also exceeds the Siegel rank \(m\). In particular both
the theta integral and the Eisenstein sum converge
absolutely by \href{https://msp.org/pjm/2013/264-1/pjm-v264-n1-p11-s.pdf}
{Yamana, Theorem~2}. Every rational maximal isotropic plane
of \(D^-\oplus(-D^-)\) is a graph. Its intersection with
either summand is isotropic in a definite space and hence
zero. As in the proof of Theorem~\ref{thm:complementgram812},
unfolding these graphs expresses the squared norm of
the lift as a positive Haar constant times the integral
of the oscillator coefficient against the constant source
pairing. The source quotient is compact, and the integral over the
target of the diagonal squared theta kernel is the continuous
convergent Eisenstein value on that quotient. Cauchy--Schwarz
bounds the mixed integrals there. The real integral is the
positive volume of the compact real source and the finite
factor differs from one by \(O(f2^{-m/2})\). This proves
strict positivity for large \(x\).

On the target compact blocks \(k_+,k_-\) the vacuum contributes
\((\det k_+)^{-m}\), from the negative Hom coordinates. The
source circle character contributes \((\det(k_+\oplus
k_-))^m\), and the polynomial contributes \((\det k_-)^{2m}\).
This gives the asserted weight. All noncompact lowering
operators differentiate once in a positive and once in a
negative target row, so they kill the polynomial. These are
the holomorphic highest weight equations, and the theta
functional intertwines the real group action. The specified
finite vacua give full level. Smooth compact-finite
cuspidal automorphic forms are rapidly decreasing and
bounded on the quotient.
\end{proof}

For a signed space with matrix \(J_0=\operatorname{diag}(1,-1)\)
in blocks, put \(\Lambda_g=E^n\cap g_f\widehat{\mathcal O}_E^n\).
Use the Euclidean norm at the embeddings of the chosen CM
type and set
\[
 \lambda(g)=\min\left(1,\min_{0\ne y\in\Lambda_g}
      \Big(\sum_{\tau}\|g_\tau^{-1}\tau(y)\|^2\Big)^{1/2}\right).
\]
A peak of a nonzero bounded form \(L\) means a point \(g\)
where \(|L(g)|>\tfrac12\sup |L|\).

\begin{lemma}[A shortest-vector bound from a holomorphic transvection]
\leavevmode\label{lem:peakminimum813}
Let \(L\) be a bounded automorphic scalar coefficient,
holomorphic of weight \(u_{\tau,i}\) for an order placing
positive coordinates before negative ones. This means
that each translate at infinity extends holomorphically
to \(G(F_\infty)B\) with right character
\(\prod_{\tau,i}b_{\tau,ii}^{u_{\tau,i}}\) on the complex
upper triangular group \(B\). Suppose \(|u_{\tau,i}|\le A\),
\(A\ge2\), and every negative weight exceeds every positive
weight at the same embedding by at least one. Let
\(\nu\in\mathcal O_E\) satisfy \(\bar\nu=-\nu\) and
\(\tau(\nu)=i t_\tau\) with \(0<t_\tau\le T,\ T\ge1\).
Suppose the finite level contains the principal congruence
subgroup modulo \(\nu\). Then any peak satisfies
\(\lambda(g)\ge(AT)^{-1/2}\).
\end{lemma}
\begin{proof}
If a nonzero \(y\in\Lambda_g\) violates the bound, put
\(v_\tau=g_\tau^{-1}\tau(y)\). Its Hermitian length
\(y^*J_0y\) lies in \(\mathcal O_F\) because the finite
coordinates at \(g_f\) are integral. Its absolute value
at every real place is less than one. The norm of a
nonzero algebraic integer cannot have this property,
so \(y\) is isotropic. In a definite space this already
gives a contradiction.

Now \(\gamma=1-\nu yy^*J_0\) is a rational isometry,
and \(g_f^{-1}\gamma g_f\) belongs to the finite principal
level. Define
\[
       n_\tau(z)=1-i t_\tau z\,v_\tau v_\tau^*J_0.
\]
The rank-one nilpotent matrices give \(n(z+1)=n(1)n(z)\).
Automorphy and the finite level imply
\(L(g n(1)h)=L(gh)\) for every real \(h\), where the
right multiplication is in the archimedean component.

Write \(Y=\Im z\) and \(\epsilon_i=(J_0)_{ii}\). Set
\(a_l=\sum_{i\le l}\epsilon_i |v_{\tau,i}|^2\) and
\(B_l=1+2t_\tau Y a_l\), suppressing \(\tau\). The
positive-first order and isotropy imply
\(0\le a_l\le\|v_\tau\|^2/2\), with equality on the right
after the last positive coordinate. Also \(a_0=a_n=0\).
The leading Gram determinants of \(n_\tau(z)^*J_0 n_\tau(z)\)
are the corresponding determinants of \(J_0\) times \(B_l\).
Indeed the Gram matrix is \(J_0+2t_\tau Y
J_0v_\tau v_\tau^*J_0\), to which the rank-one determinant
formula applies. Since \(A\sum_\tau t_\tau\|v_\tau\|^2<1\),
all \(B_l\) are positive on \(Y>-1\). Indefinite
Gram--Schmidt in this order now gives \(n_\tau(z)=h_\tau b_\tau\)
with \(h_\tau\) unitary and \(b_\tau\) upper triangular,
\(|b_{\tau,ll}|^2=B_l/B_{l-1}\). Therefore the holomorphic
extension evaluated at \(g n(z)\) defines a holomorphic
function on this half-plane, periodic by the preceding
translation and triangular equivariance.

Let \(C=\sup|L|\). For \(Y\ge0\) the diagonal moduli are
at least one on the positive block and at most one on the
negative block. Subtracting the largest positive weight
from every weight does not affect the modulus of the
character since \(|\det b_\tau|=1\). The weight gap bounds
its modulus by the product of the negative diagonal moduli,
namely \((1+t_\tau Y\|v_\tau\|^2)^{-1/2}\).
Thus the periodic function tends to zero as \(Y\to\infty\).
For \(-1<Y<0\) the two inequalities for the diagonal
moduli reverse. The absolute weight bound gives a character
modulus at most \((1-t_\tau\|v_\tau\|^2)^{-A}\).
The product of these bounds over the real places is at
most \(e^2\), since \(-\log(1-z)\le2z\) for
\(0\le z\le1/2\). The periodic function is consequently
bounded by \(Ce^2\). In the disk coordinate
\(\exp(2\pi i(z+i))\) it extends with value zero at the
origin. The Schwarz lemma at \(z=0\) gives
\(|L(g)|\le C e^{2-2\pi}<C/2\), contradicting the peak.
\end{proof}

\begin{theorem}[Simultaneous minima for the two forms]
\leavevmode\label{thm:jointminima813}
Let \(M=p_0p_1\), and let \(F_0\ne0\) be any good form on
\(U(W_j)\) on its specified real line. Choose any peak
\(g_1\) of the carrier and any peak \(g\) of \(F_0\).
At every sufficiently distant node of the fixed-\((q_0,d)\)
family,
\[
 \lambda_S(g_1)\ge(3m\sqrt d)^{-1/2},\qquad
 \lambda_{W_j}(g)\ge(5b_x M\sqrt d)^{-1/2},\qquad
 \lambda_S(g_1)\lambda_{W_j}(g)\ge b_x^{-12}.
\]
The last bound also holds after adjoining the standard
padding \((-I_j)\oplus I_{3m-j}\) to \(W_j\) and extending
\(g\) by the identity.
\end{theorem}
\begin{proof}
Use the sign of \(\sqrt{-d}\) whose imaginary part is
positive at the chosen embeddings. The CM type is induced
from a single embedding of \(\QQ(\sqrt{-d})\), as in
the proof of Theorem~\ref{thm:negativeinterval740}.
Theorem~\ref{thm:carrier813} supplies a bounded nonzero
holomorphic form of full level. Apply
Lemma~\ref{lem:peakminimum813} with
\(\nu=\sqrt{-d}\) and \(A=3m\).

For the good form apply the lemma to \(\overline{F_0}\),
which has the same peaks. We spell out the holomorphic
convention. The ordered Harish--Chandra entries for
\(W_j\) are the \(k-1\) centered consecutive integers
followed by \(H_0-w_1,\ldots,H_0-w_m\), with the largest
\(k+j-1\) entries positive in signature. Negate the
entries for the conjugate representation and sort each
compact block decreasingly, calling the result \(\eta\).
Every negative-block entry then exceeds every positive-block
entry by at least \(2H_0\) when both blocks occur. The
separation is the adjacent gap between \(H_0-w_j\) and
\(H_0-w_{j+1}\). The fixed line becomes the compact highest
line in this reversed order. The compact-Cartan description
of the holomorphic discrete series gives its weight
\(u=\eta+\rho_{\rm n}-\rho_{\rm c}\), with the half-sums
of positive noncompact and compact roots in the chamber
of \(\eta\). In this chamber all positive noncompact roots
point from the smaller positive-sign block entries to
the larger negative-sign block entries. The lowest compact
type is killed by the opposite noncompact root spaces.
Its compact highest line is also killed by the compact
upper root spaces. Together these are exactly the upper
triangular root spaces for the positive-first order.
To obtain the holomorphic extension in the lemma, extend
to the open set \(G(F_\infty)B\) by the triangular character.
This is well-defined
since \(G(F_\infty)\cap B\) is the compact diagonal torus.
The root equations and the diagonal weight are the
Cauchy--Riemann equations for this extension. This also
describes the compact case \(j=m\) by the usual algebraic
highest weight representation. For the holomorphic model
of the noncompact groups see
\href{https://msp.org/pjm/1987/129-1/pjm-v129-n1-p02-s.pdf}
{Davidson, \S\S3--4}.

Each entry of \(\rho_{\rm n}-\rho_{\rm c}\) has modulus
at most \(s\), since at most \(s-1\) roots involving that
coordinate enter the half-sums. Thus the gap is at least
\(2H_0-2s\ge1\). The weight bound \(w_j\le2b_x\) from
Theorem~\ref{thm:gluedcost811} holds also for \(w_m\),
and \(k,s,H_0\ll x^{40}\), so \(|u_i|\le5b_x\) everywhere
for large \(x\). The good form is cuspidal and bounded.
Its finite level contains the subgroup principal modulo
\(M\sqrt{-d}\), since divisibility by \(M\) gives the
prescribed congruence at every old mark. Apply the lemma
with that element and \(T=M\sqrt d\).

Peaks exist for both forms by the definition of supremum.
Multiplying the two estimates gives
\((15m b_x M d)^{-1/2}\). Since \(M\le b_x^{22}\) on
these intervals and \(15md\le b_x\) eventually, this is
at least \(b_x^{-12}\). On the padded summands at the
identity a nonzero integral coordinate has the sum of
its squared archimedean moduli at least one, by the norm
of an algebraic integer. The squared Euclidean lengths add
in the direct sum, proving the last assertion.
\end{proof}

These estimates locate the points where the first-factor
theta kernels must approximate the prescribed real
coefficients. Such an approximation still has to use the
finite residue data and level of the integral periods.
\clearpage

\subsection{Approximating the real kernel at the prescribed finite level}
\label{sec:sparseperiod814}
Keep the carrier space \(S\), the good form \(F_0\), and the
two peaks of Theorem~\ref{thm:jointminima813}. Put
\(T=W_j\oplus(-I_j)\oplus I_{3m-j}\), and let \(g_2\)
denote the extension of the peak of \(F_0\) by the identity
on the padding. We work at \((g_1,g_2)\).
Let \(r=3m,\ r'=s+3m,\ f=[F:\QQ]\), and set
\(n_*=2ft^3\). Both ranks are at most \(t=5s\).
The finite residue data will be functions on
\(\operatorname{Mat}_{r',r}(\mathcal O_E/\mathfrak a^2)\),
where \(\mathfrak a\) is the product of the designated old
places of \(E\). Thus \(\mathfrak a\bar{\mathfrak a}=(M)\)
with \(M=p_0p_1\). At the other places we keep the lattice
vacuum. Write \(b_x\) for the large weight of the node.

\begin{lemma}[Separation of the rational Hom vectors]
\leavevmode\label{lem:sparsedistance814}
For a rational \(E\)-linear map \(y:S\to T\) write
\(y(g_1,g_2)=g_2^{-1}y g_1\) for its adelic coordinates.
Let \(\alpha\) be the finite scale of the self-dual oscillator
lattice, so the local vacuum lattice on Hom in these
coordinates is \(\alpha\operatorname{Mat}_{r',r}
(\widehat{\mathcal O}_E)\). Suppose
\[
  \alpha^{-1}y(g_1,g_2)_f\ \in\
        \mathfrak a^2\operatorname{Mat}_{r',r}
             (\widehat{\mathcal O}_E).
\]
At sufficiently large nodes the real Fock coordinates of
these rational maps are \(b_x^6\)-separated in Euclidean
distance. In particular every nonzero such vector has
length at least \(b_x^6\).
\end{lemma}
\begin{proof}
Let \(K_0=q_1 d\). The field \(E\) is contained in the
\(K_0\)-th cyclotomic field, so \(K_0\mathfrak D_E^{-1}
\subset\mathcal O_E\), with \(\mathfrak D_E\) the absolute
different. Indeed the different of the full cyclotomic
field is generated by \(\Phi_{K_0}'(\zeta_{K_0})\), which
divides \(K_0\) by differentiation of \(X^{K_0}-1\).
Transitivity of the different proves the assertion for
the subfield. In particular \(|D_E|^{1/[E:\QQ]}\le K_0\).
In the prescribed oscillator normalization the imaginary
scalar is \(-1/\sqrt{-d}\), and \(q_1\alpha\) is integral.
For the latter fact, the finite scale is one or the
reciprocal of the local generator of the cyclotomic relative
different, with the inverse placed in the specified component.
Since \(d\) is a prime congruent to \(7\bmod 8\) and \(d\nmid q_1\), the discriminants of
\(\QQ(\sqrt{-d})\) and \(\QQ(\mu_{q_1})\) are coprime.
Their integer rings tensor to the integer ring of \(E\),
so the relative different over the quadratic field is
generated by \(\Phi_{q_1}'(\zeta_{q_1})\), a divisor of \(q_1\).

Use the Euclidean realization of the source lattice
\(\Lambda_{g_1}\) in Lemma~\ref{lem:peakminimum813}, of
real dimension \(2fr\). Its Euclidean dual is obtained
by applying \(2J_0\) to the lattice in these coordinates
corresponding to \(\mathfrak D_E^{-1}\Lambda_{g_1}\).
Indeed the absolute trace of its Hermitian pairing is
twice the real Euclidean pairing after inserting \(J_0\),
and \(\Lambda_{g_1}\) is Hermitian self-dual at every
finite place. Thus the shortest dual length is at least
\(2\lambda_S(g_1)/K_0\).
The Euclidean transference inequality
\(\lambda_D(L)\lambda_1(L^\vee)\le D\) in dimension \(D\)
gives \(2fr\) independent source vectors of length at
most \(2frK_0/\lambda_S(g_1)\).
See \href{https://doi.org/10.1007/BF01445125}
{Banaszczyk, \emph{New bounds in some transference theorems
in the geometry of numbers}}. These vectors span over
\(\QQ\) and hence over \(E\). A nonzero integral Hom map
sends at least one of them to a nonzero vector of the target
lattice. The sum-of-squares Frobenius norm of its real
matrices is therefore at least
\(\lambda_S(g_1)\lambda_T(g_2)/(2frK_0)\).

Minkowski's convex body theorem on the fractional ideal
\(\mathfrak a^{-2}\) gives \(c_0\ne0\) in this ideal with
\(\max_\tau|\tau(c_0)|\le 2\sqrt{K_0}/M\).
In fact the real covolume is
\(2^{-f}\sqrt{|D_E|}\,M^{-2f}\).
The product of the \(f\) disks of the indicated radius
has greater than \(2^{2f}\) times this volume. This choice
works for any designated member of each old split pair.
Multiplying \(y\) in the statement by \(q_1c_0\) makes it
an integral Hom map between the moved lattices. Its
archimedean coordinates grow by a factor of at most
\(2q_1\sqrt{K_0}/M\). In Fock coordinates Euclidean
matrix length is multiplied by
\((2\pi/\sqrt d)^{1/2}\), from the imaginary scalar above.
The lower bound for a nonzero displacement is consequently
\[
       (2\pi/\sqrt d)^{1/2}
       \frac{M\,\lambda_S(g_1)\lambda_T(g_2)}
            {4 f r q_1 K_0^{3/2}} .
\]
Here \(M\ge b_x^{20}\), and the product of minima is
at least \(b_x^{-12}\). All factors in the denominator
are polynomial in \(x\) of total order at most \(77\),
with constant depending only on \(q_0,d\). Indeed
\(f\ll x^{12}\), \(r\le t\ll x^5\), and
\(q_1\le2\varphi(q_1)^2\le 8 f^2\), the first inequality
following prime power by prime power in Euler's formula.
Since \(b_x\ge x^{400}\), the bound exceeds \(b_x^6\)
for large \(x\). Subtracting two admissible vectors
preserves the finite congruence, giving the separation.
\end{proof}

\begin{lemma}[Decay under a Heisenberg displacement]
\leavevmode\label{lem:oscillatordecay814}
In real Fock coordinates let \(\rho(z)\) be displacement
by a real phase-space vector \(z\), and \(\omega(U)\) a
unitary metaplectic operator covering a real linear
symplectic map \(U\). For a polynomial vector \(v\) of
degree at most \(D\), an integer \(L\ge1\), and \(z\ne0\),
\[
 |\langle v,\rho(z)\omega(U)v\rangle|
 \ \le\ \|v\|^2
 \left(\frac{C(1+\|U^{-1}\|)\sqrt{D+L+1}}{|z|}\right)^L
\]
with an absolute constant in the Euclidean Fock normalization,
independent of the number of variables.
\end{lemma}
\begin{proof}
Let \(X_e\) be the self-adjoint generator of displacement in
the unit direction \(e\). Choose \(e\) symplectically dual
to \(z/|z|\). The Euclidean metric and the standard
symplectic form are compatible, so the commutator
\([X_e,\rho(z)]\) is a scalar of modulus \(|z|\) times
\(\rho(z)\), up to the fixed normalization of the central
character. Covariance gives
\(\omega(U)^{-1}X_e\omega(U)=X_{U^{-1}e}\).
The field operator \(X_w\) is a sum of one creation and
one annihilation operator in the complex direction of \(w\).
On total degrees through \(d_0\) their norms are bounded
by \(|w|\sqrt{d_0+1}\) times a fixed constant, by using
an orthonormal basis with that direction first. A field
operator raises the maximum degree by at most one.
The \(L\)-fold commutator with \(X_e\), paired between
\(v\) and \(\omega(U)v\), gives a binomial sum. In its
term with \(l\) operators on the right, move them through
\(\omega(U)\) and bound the two norms by the preceding
degree estimate. Summation gives
\(\|v\|^2(C\sqrt{D+L+1})^L(1+\|U^{-1}\|)^L\).
Division by \(|z|^L\) proves the result. All differentiated
vectors belong to Schwartz space.
\end{proof}

\begin{theorem}[Interpolation by the prescribed residue functions]
\leavevmode\label{thm:sparseapprox814}
Let \(v\) be any archimedean Fock polynomial vector on
\(\operatorname{Hom}_E(S,T)\) of degree at most \(b_x^{3/2}\).
For the unitary splitting characters of this pair there
exists one function \(\Psi\) on the residue matrices
modulo \(\mathfrak a^2\) such that, for all real \(u\in U(S)\)
and \(u'\in U(T)\),
\[
\begin{split}
 \big|\theta_{\Psi,v}(g_1u,g_2u')
       -\langle v,\omega_\infty(u,u')v\rangle\big|
 &\le b_x^{-20n_*}\|v\|^2
           \bigl(\|u\|_{\max}\|u'\|_{\max}\bigr)^{4n_*}.
\end{split}
\]
The maximum is the largest Euclidean operator norm over
the embeddings. The residue function depends on the two
peaks and \(v\) and works simultaneously for all their
real translates. The threshold on the interval is
independent of the peaks and \(j\).
\end{theorem}
\begin{proof}
Write \(\mathcal D\) for the additive group of residue
matrices, and use all its additive characters \(\chi\).
At the chosen split places let the finite position test be
the lattice indicator times \(\chi\) of the residue
matrix, using probability mass for the position lattice.
Use unit vacua at all other places. Call this test \(\phi_\chi\).
By linearity the function
\[
 \Psi=|\mathcal D|^{-1}\sum_{\chi\in\widehat{\mathcal D}}
             \overline{\theta_{\chi,v}(g_1,g_2)}\,\chi
\]
gives a kernel equal to the average of the products
\(\overline{\theta_{\chi,v}(g_1,g_2)}
                    \theta_{\chi,v}(g_1u,g_2u')\).
The characters form the full dual even when the modulus has
several prime factors.

We give the Poisson identity evaluating this average.
For any two adelic Schwartz vectors \(v_1,v_2\) in a rational
position polarization \(X\), the product of theta functionals
\(\overline{\theta(v_1)}\theta(v_2)\) equals
\(\sum_{y}\langle v_1,\rho(y)v_2\rangle\) over the full
rational phase space. To prove it, group the two rational
position variables by their difference \(a\in X(F)\).
Poisson summation of
\(\overline{v_1(x)}v_2(x+a)\) in \(x\) introduces all
dual rational frequencies. The half-pairing in the Weyl
operator has trivial character on rational points. Use
self-dual measures, for which the rational quotient has
volume one. The matrix coefficients on phase space are
Schwartz, being partial Fourier transforms after a linear
change of the two position variables. Hence these sums
converge absolutely. Local changes of polarization are
unitary and preserve the inner products.

Apply this identity to the two translates in the average.
Moving the common \((g_1,g_2)\) through the Heisenberg
operator replaces the rational vector \(y\) by
\(y(g_1,g_2)\). The common splitting phase cancels. In the
split position model the finite inner product of a character
test with its displacement by position \(a\) and frequency
\(b'\) is zero unless the normalized \(a,b'\) are integral.
Indeed the supports of the two lattice indicators must
meet, and the frequency must integrate trivially over the
lattice. On that support the coefficient is the character
of the position residue of \(a\), up to a fixed phase
independent of \(\chi\). Reversing the displacement convention
inverts all the residue characters and gives the same
average. At unmarked places the unit vacuum has coefficient
of modulus one on its self-dual phase-space lattice and
zero outside, by commutation with that lattice.
Orthogonality of \(\widehat{\mathcal D}\) thus leaves
exactly the finite condition of Lemma~\ref{lem:sparsedistance814},
with a phase of modulus one. In particular the \(y=0\)
term is precisely \(\langle v,\omega_\infty(u,u')v\rangle\).

For the other terms use \(L=4n_*\) in
Lemma~\ref{lem:oscillatordecay814}. The symplectic map
on Hom has inverse norm at most
\(V=\|u\|_{\max}\|u'\|_{\max}\), because each unitary
matrix preserving a sign form has inverse of the same
Euclidean norm. Thus \(V\ge1\). In real dimension
\(d_0=2frr'\) the remaining displacements have separation
\(R=b_x^6\) and are at distance at least \(R\) from zero.
Balls of radius \(R/2\) about these points have disjoint
interiors. The shell \(2^h R\le |z|<2^{h+1}R\) contains
at most \((1+2^{h+2})^{d_0}\) points by comparison of
volumes. Since \(d_0\le n_*\), summing over shells gives
\(\sum_{z\ne0}|z|^{-L}\le(6/R)^L\).
Finally \(n_*\ll_{q_0,d}x^{27}\) and
\(D\le b_x^{3/2}\). For sufficiently large nodes
\(12C\sqrt{D+L+1}/R\le b_x^{-5}\).
The asserted estimate follows on taking the \(L\)-th power.
\end{proof}

\begin{corollary}[The determinant polynomial of the first factor]
\leavevmode\label{cor:paddedminor814}
The interpolation theorem applies to the following actual
polynomial with the splitting characters
\((\upsilon^m,\chi^0\upsilon^{\,m+j+(t+k-1)/2-2js^8})\).
Here \(\chi^0\) is the fixed character of weight \(-w_0\)
and trivial tame part. Put \(a=k+j-1\), \(p=s+9m+1\),
\(\Delta'=s^8-(t+1)/2\), and let \(c^+_{h,\tau}\) be the
positive column counts of Theorem~\ref{thm:gluedcost811}.
For each \(1\le h<a\) form a \(2m\)-rowed positive square
minor \(D_{h,\tau}\) using the \(2m\) positive carrier
columns. For the \(i\)-th row, \(0\le i<2m\), take positive
row \(h+i\) of \(W_j\) if \(h+i<a\), and otherwise positive
padding row \(i\), numbering from zero. Form the \(m\)-rowed
square \(D_{-,\tau}\) on the negative carrier columns and
all negative padded target rows, with the \(j\) padding
rows first. The polynomial is
\[
 \prod_\tau\left(D_{-,\tau}^{\,p+m+\Delta'+\sum_{1\le h<a}c^+_{h,\tau}}
              \prod_{1\le h<a}D_{h,\tau}^{\,c^+_{h,\tau}}\right).
\]
In particular there is a linear combination of the elementary
residue-character tests whose value at the two peaks is
nonzero and differs from the squared Fock norm of this
polynomial by at most \(b_x^{-20n_*}\) times its squared norm.
\end{corollary}
\begin{proof}
Each square uses distinct rows, with at least \(2m\)
available positive padding rows. Each determinant is
therefore a nonzero polynomial and their product is nonzero.
At one embedding write \(A_\tau=\sum_{h<a}c^+_{h,\tau}\).
The degree there is \(2m A_\tau+m(p+m+\Delta'+A_\tau)\).
The telescoping count in Theorem~\ref{thm:gluedcost811}
gives \(A_\tau\le w_j+k\ll b_x\), and the height-zero count
vanishes. Also \(p+m+\Delta'\ll b_x\). The total degree
is \(O(fm b_x)\le b_x^{3/2}\) on the same intervals.
Apply the theorem and take \(u=u'=1\).
\end{proof}

The coefficients in the displayed residue average are
complex values of theta functions. The result supplies the
approximation used to test complex spanning. The arithmetic
comparison still requires pairing these tests with the
carrier and controlling the prescribed integral first factors.
\clearpage

\subsection{Integrating the theta kernel against the carrier at infinity}
\label{sec:carrierintegration815}
Fix the field, the two old marks, and the index \(j\) of
Section~\ref{sec:sparseperiod814}. Let \(\phi\) be the carrier,
\(g_1\) any peak of \(\phi\), and \(g_2\) the padded peak
of a nonzero good form. Write \(G_S=\prod_{\tau}U(2m,m)\)
and \(G_T=\prod_{\tau}U(s+2m,m)\) for the two real groups.
Let \(v\) be the determinant polynomial of
Corollary~\ref{cor:paddedminor814}, regarded as a Fock vector.
Define
\[
 c_v(h,u)=\langle v,\omega_\infty(h,u)v\rangle,\qquad
 c_v(1,1)=\|v\|^2>0 .
\]
Here and below real translates have identity finite component.
We shall integrate in \(h\) and compare the resulting value
to an explicit positive mass.

At one place write \(h=\left(\begin{smallmatrix}A&B\\ C&D\end{smallmatrix}\right)\)
in positive and negative coordinates. For integers \(a_0\le b_0\) put
\[
 H_{a_0,b_0}(h)=\frac{\det(h)^{b_0}}{\det(A)^{b_0-a_0}} .
\]
The relation \(A^*A-C^*C=I\) guarantees that the denominator
is nonzero. On \(K=U(2m)\times U(m)\) the character of this
function on either side is
\(\det(k_+)^{a_0}\det(k_-)^{b_0}\). For a list of weight
pairs \(u_0\) indexed by the real places, write \(H_{u_0}\)
for the product of these functions.

Use \(p=s+9m+1\), \(\Delta'=s^8-(t+1)/2\), and set
\[
 A_\tau=\sum_{1\le h<a}c^+_{h,\tau},\qquad
 \mu_\tau=A_\tau+\Delta',\qquad a=k+j-1 .
\]
Define four weights, giving their positive and negative
entries in that order,
\[
 \begin{aligned}
 u_\phi&=(0,3m),                 &u_1&=(-\Delta',6m+1),\\
 u_{2,\tau}&=(\Delta',\mu_\tau+s+2m),
                            &u_{3,\tau}&=(0,\mu_\tau+p+2m).
 \end{aligned}
\]
The identity \(u_\phi+u_1+u_2=u_3\) holds entry by entry.

\begin{lemma}[The carrier variable of the real coefficient]
\leavevmode\label{lem:carriercoefficient815}
For \(h\in G_S\),
\[
 c_v(h,1)=\|v\|^2 H_{u_2}(h)\overline{H_{u_3}(h)} .
\]
\end{lemma}
\begin{proof}
The target circle character acts on the source determinant
with exponent \(-e_\tau\), where \(e_\tau=m-\mu_\tau\).
To check the exponent, use the list \(\sigma_i\) in the proof
of Theorem~\ref{thm:gluedcost811}. Through index \(a-1\)
its entries are at least \(\sigma_{a-1}=s^8-w_{j,\tau}\),
which exceeds \((t-1)/2-w_{j,\tau}\). Its consecutive
entries differ by at least one there. The sum of the positive
column counts is therefore
\(\sigma_0-\sigma_{a-1}-(a-1)
  =w_{j,\tau}-s^8+k/2-a\).
It follows that \(\mu_\tau=w_{j,\tau}-j-(t+k-1)/2\).
Since \(w_{j,\tau}=w_{0,\tau}+2js^8\), the exponent of the
circle character in Corollary~\ref{cor:paddedminor814},
including the weight of \(\chi^0\), is \(m-\mu_\tau\).

We compute the normalized coefficient at each real place
by compact multiplication and on a real diagonal boost.
There are \(s+2m\) positive and \(m\) negative target rows.
On \(k=(k_+,k_-)\) a positive target row contributes
\(\det k_-\) to the vacuum and a negative target row
contributes \(\det k_+\). Indeed the vacuum multiplier is
the inverse determinant of the action on negative
phase-space coordinates. The action on a Hom row is by
the inverse source matrix, and exactly the source columns
of opposite sign are negative coordinates. The splitting
character contributes \(\det(k)^{-e_\tau}\). A positive minor
in \(v\) uses all positive source columns, and the negative
minor uses all negative ones. In the variables of like signs
the action on these determinants contributes
\(\det(k_+)^{-A_\tau}\det(k_-)^{-\kappa_\tau'}\), with
\(\kappa_\tau'=p+m+\mu_\tau\). Thus the two compact
weights are
\[
       m-e_\tau-A_\tau=\Delta',\qquad
       s+2m-e_\tau-\kappa_\tau'=s-p .
\]
They are the weights of \(H_{u_2}\overline{H_{u_3}}\).
As the vector transforms by a unitary character, its
coefficient has this character on both sides.

Now pair \(m\) positive and negative source coordinates in
a Cartan boost with parameters \(\xi_1,\ldots,\xi_m\).
In each target row the two oscillator variables for each
pair have opposite phase-space signs. Every letter in the
polynomial has positive phase-space sign. In the creation,
diagonal, annihilation factorization of the boost, the
annihilation terms kill it because they differentiate in
an opposite-sign pair. Every nonconstant creation term is
orthogonal to it because it contains a negative letter.
The diagonal term supplies \((\cosh\xi_i)^{-1}\) for the
vacuum in each row and rescales each paired coordinate by
\((\cosh\xi_i)^{-1}\), as in Lemma~\ref{lem:oscillatorcartan710}.
Each determinant uses every paired column of its own sign
exactly once. The normalized coefficient on the boost is
therefore
\[
       \prod_{i=1}^{m}(\cosh\xi_i)^{-(s+3m+A_\tau+\kappa_\tau')}.
\]
The determinant of the boost is one, and \(\det A\) is
\(\prod_i\cosh\xi_i\). The displayed exponent is the sum
of the two gaps \(u_{2,-}-u_{2,+}\) and \(u_{3,-}-u_{3,+}\).
The Cartan decomposition gives equality on the whole group.
Taking the tensor product over the real places proves the lemma.
\end{proof}

\begin{lemma}[A comparison of two masses on the real carrier group]
\leavevmode\label{lem:carrierweightmass815}
With the same Haar measure in both integrals, let
\(J=\int_{G_S}|H_{u_3}(h)|^2\,dh\). Uniformly for large
nodes on these intervals,
\[
       0<J<\infty,\qquad
       \int_{G_S}|H_{u_1}(h)|\,\|h\|_{\max}^{4n_*}\,dh
                   \ \le\ b_x^{n_*}J .
\]
\end{lemma}
\begin{proof}
Write \(r=3m\). The map \(hK\mapsto Z=CA^{-1}\) identifies
the symmetric space with the matrix ball
\(\mathcal B=\{Z\in M_{m,2m}(\CC):I-ZZ^*>0\}\).
Indeed \(Z\) gives the graph of a positive \(2m\)-plane.
Orthonormal bases for this plane and its negative orthogonal
complement recover \(h\) up to \(K\). Put
\(\delta(Z)=\det(I-ZZ^*)\). The identity
\(A^*(I-Z^*Z)A=I\) gives
\(\delta=|\det A|^{-2}\). In particular
\(|H_{a_0,b_0}|=\delta^{(b_0-a_0)/2}\).

The quotient measure in this chart is
\(c\,\delta^{-r}\,dZ\) for a constant \(c>0\), where \(dZ\)
is real Lebesgue measure. To verify it, under a block matrix
\(\left(\begin{smallmatrix}A&B\\ C&D\end{smallmatrix}\right)\)
the new coordinate is \(Z'=(C+DZ)(A+BZ)^{-1}\), and its
differential sends \(X\) to \((D-Z'B)X(A+BZ)^{-1}\).
The block determinant formula gives
\(\det(D-Z'B)=\det(h)/\det(A+BZ)\).
Hence the real Jacobian is \(|\det(A+BZ)|^{-2r}\).
Preservation of the Hermitian form gives
\(\delta(Z')=\delta(Z)/|\det(A+BZ)|^2\).
These formulas give the asserted invariant density and the
same constant for all integrals at this place.

The largest singular value of a boost of parameter
\(\xi\ge0\) is \(e^\xi\le2\cosh\xi\). All factors
\(\cosh\xi_i\) are at least one. Thus
\(\|h\|_{\max}\le2\prod_\tau|\det A(h_\tau)|\).
Since \(s\gg x^5\), \(n_*=2ft^3\ll x^{27}\), and
\(\Delta'=s^8-(t+1)/2\), one has
\((6m+1+\Delta'-4n_*)/2\ge r\) on this interval for
all sufficiently large \(x\). The weighted integral in the
statement is bounded above by \(2^{4n_*}\) times the
product of the Lebesgue volumes of the balls, each with
its measure constant. This follows because each remaining
exponent of \(\delta\) is nonnegative and \(\delta\le1\).

Put \(\ell_\tau=\mu_\tau+p+2m\) and
\(L=1+\max_\tau\ell_\tau\). We have
\(r\le\ell_\tau\) and \(L\ll b_x\) by the column count
and the bounds of Corollary~\ref{cor:paddedminor814}.
The local integrand for \(J\) against \(c\,dZ\) is
\(\delta^{\ell_\tau-r}\), with values at most one.
On the smaller ball \(\eta\mathcal B\),
where \(\eta=(2mL)^{-1/2}\), it is at least
\((1-\eta^2)^{mL}\ge e^{-1}\). This uses
\(\log(1-y)\ge-2y\) for \(0\le y\le1/2\).
As the complex dimension is \(2m^2\), this smaller ball
has \((2mL)^{-2m^2}\) times the volume of the full ball.
This proves positivity and finiteness and bounds the
ratio of the two integrals by
\[
          2^{4n_*} e^f (2mL)^{2m^2f}.
\]
Its logarithm divided by \(n_*\log b_x\) tends to zero.
In fact \(\log(2mL)=O(\log b_x)\) and \(2m^2 f/n_*=m^2/t^3\)
tends to zero on the original interval. This proves the
claimed uniform comparison.
\end{proof}

\begin{theorem}[Nonvanishing after integration against the carrier]
\leavevmode\label{thm:integratedpeak815}
Choose \(\Psi\) from Theorem~\ref{thm:sparseapprox814}
for the two peaks and \(v\). For \(u\in G_T\) put
\[
 \begin{aligned}
  B_0(u)&=\int_{G_S}\phi(g_1h)c_v(h,u)H_{u_1}(h)\,dh,\\
  B_\Psi(u)&=\int_{G_S}\phi(g_1h)
         \theta_{\Psi,v}(g_1h,g_2u)H_{u_1}(h)\,dh .
 \end{aligned}
\]
Both integrals converge absolutely and \(B_0\) is bounded.
At every sufficiently large node they satisfy
\[
 \begin{aligned}
    B_0(1)&=\|v\|^2\phi(g_1)J\ne0,\\
    |B_\Psi(u)-B_0(u)|
       &\le 2 b_x^{-19n_*}\|u\|_{\max}^{4n_*}|B_0(1)|.
 \end{aligned}
\]
In particular \(B_\Psi(1)\ne0\) at the padded peak of
every nonzero good form. At each such peak some elementary
additive character of the prescribed residue group also
has a nonzero integral of this kind.
\end{theorem}
\begin{proof}
The carrier satisfies \(\sup|\phi|<2|\phi(g_1)|\).
Unitarity gives \(|c_v(h,u)|\le\|v\|^2\) for every \(h,u\).
The preceding lemma gives absolute convergence and a bound
independent of \(u\) for the first integral.

At one real place a holomorphic scalar group function of
weight \((a_0,b_0)\), divided by \(H_{a_0,b_0}\), descends
to a holomorphic function on \(\mathcal B\). We recall
the trivialization to fix the orientation. The weight on the
compact blocks is scalar and extends on the complex block
upper parabolic to \(\det(P_{11})^{a_0}\det(P_{22})^{b_0}\).
The upper triangular transformation law and the compact
block law give this parabolic law by factoring each
diagonal block into a unitary and an upper triangular
matrix. For \(Z=CA^{-1}\), the graph matrix
\(\left(\begin{smallmatrix}I&0\\ Z&I\end{smallmatrix}\right)\)
is a holomorphic section. Its inverse times \(h\) is
block upper with diagonal \(A,D-ZB\).
The character of this block upper factor is precisely
\(\det(h)^{b_0}/\det(A)^{b_0-a_0}\). Evaluating the
holomorphic extension on the section proves the assertion.
The same argument works at all embeddings together.

Apply this to \(h\mapsto\phi(g_1h)\) of weight \(u_\phi\)
from Theorem~\ref{thm:jointminima813}, and write the quotient
as \(Q(Z)\). Since \(H_{u_\phi}H_{u_1}H_{u_2}=H_{u_3}\),
Lemma~\ref{lem:carriercoefficient815} identifies the
integrand for \(B_0(1)\) with
\(\|v\|^2 Q(Z)|H_{u_3}(h)|^2\).
Both its mass measure and the product of the balls are
invariant under simultaneous rotation \(Z_\tau\mapsto
e^{i\theta}Z_\tau\). The circle average of \(Q\) on
each orbit is \(Q(0)=\phi(g_1)\). Indeed the function of
one variable \(z\mapsto Q((zZ_\tau)_\tau)\) is holomorphic
on a disk of radius greater than one. Absolute integrability
and rotation invariance justify averaging under the integral.
This gives the first identity.

For the second, Theorem~\ref{thm:sparseapprox814} bounds
the pointwise difference of kernels by
\[
 b_x^{-20n_*}\|v\|^2\|u\|_{\max}^{4n_*}\|h\|_{\max}^{4n_*}.
\]
After multiplication by the carrier and \(H_{u_1}\),
Lemma~\ref{lem:carrierweightmass815} proves absolute
integrability of the difference and bounds its integral
in modulus by
\(2|\phi(g_1)|\|v\|^2 b_x^{-19n_*}J\|u\|_{\max}^{4n_*}\).
At \(u=1\) this error is strictly smaller than \(|B_0(1)|\).

For clarity, each individual residue character \(\chi\)
also defines an absolutely convergent integral at \(u=1\).
Apply the two-vector Poisson identity in the proof of
Theorem~\ref{thm:sparseapprox814} to the squared modulus
of its kernel, with both translates equal to \((g_1h,g_2)\).
Each finite coefficient has modulus at most one and vanishes
outside the moved integral phase-space lattice. Before the
real translation by \(h\), the corresponding rational
vectors form a fixed lattice \(\mathcal L\) in real Hom
coordinates. Indeed intersecting the local integral
lattices gives a full fractional integer lattice in the
rational Hom space, and the fixed real coordinates differ
by an invertible linear map. For \(z\ne0\) in \(\mathcal L\),
the further displacement has length at least
\(|z|/\|h\|_{\max}\). Lemma~\ref{lem:oscillatordecay814}
with \(U=1\) and \(L'=4n_*\) bounds the absolute sum of
the nonzero terms by a fixed multiple of
\(\|h\|_{\max}^{L'}\sum_{z\ne0}|z|^{-L'}\), which is
finite since \(L'\) exceeds the real dimension. Taking
square roots shows that the individual kernel is
\(O(\|h\|_{\max}^{2n_*})\), with a constant allowed to
depend on these fixed data. The mass bound gives the
required integrability. Expanding the finite sum defining
\(\Psi\) under the integral now proves the last assertion.
\end{proof}

The integral here is over \(G_S\) with \(g_1\) fixed. The
arithmetic comparison uses periods on the adelic quotient
of the carrier group. Relating these integrals requires
the definite auxiliary theta kernels and their prescribed
residue functions and row choices.
\clearpage

\subsection{The adelic periods of the first factor}
\label{sec:adelicfirst816}
Write \(r=3m\), \(N=6m+1=2r+1\), and \(V^\flat=I_N\).
We use \(U(V^\flat)\) to express the carrier integral of
Theorem~\ref{thm:integratedpeak815} in terms of the prescribed
periods on \([U(S)]\). The small excess of \(N\) over
\(2r\) requires a more precise calculation at the finite
places. In this subsection a positive lower bound is needed
on each fixed field. The real approximation will retain its
uniform constants.

\begin{lemma}[Absolute convergence at rank \(2r+1\)]
\leavevmode\label{lem:thresholdtail816}
At an unmarked finite place of the present fields let
\(q\) be the residue cardinality in \(F\), \(K_v\) the
stabilizer of the integral lattice in \(S\), and \(A_v\)
the normalized coefficient on \(U(S)\) of the unit lattice
vacuum for \(V^\flat\). Give \(K_v\) mass one. Then
\[
     \|A_v-\mathbf 1_{K_v}\|_1
     \ \ll_r\
     \begin{cases}
       q^{-2},&E_v/F_v\text{ inert},\\
       q^{-3/2},&E_v/F_v\text{ split}.
     \end{cases}
\]
\end{lemma}
\begin{proof}
These are all places by Corollary~\ref{cor:fieldrelations712}.
The coefficient is one on \(K_v\). At an inert place let
\(Q=q^2\) and let \(L,L'\) be the integral lattice and a
translate. Fix \(I=L\cap L'\) with \(|L/I|=Q^i\) and put
\(A=L/I\). For a fixed module type of \(A\) write
\(\lambda_1,\ldots,\lambda_b\) for its nonzero elementary
divisor exponents, and \(d_A=\sum_{u,v}\min(\lambda_u,\lambda_v)\).
There are at most \(Q^{ri}/|\operatorname{Aut}(A)|\)
possible kernels \(I\), by counting surjections out of \(L\).
The endomorphism ring has order \(Q^{d_A}\), and
\[
   |\operatorname{Aut}(A)|\ge c_r Q^{d_A},\qquad
   c_r=\left(\prod_{a=1}^{\infty}(1-2^{-a})\right)^r>0 .
\]
Indeed an endomorphism is invertible if and only if its
reduction on \(A/\pi A\) is invertible. In order of
elementary divisor lengths the reduction is block triangular.
The equal-length diagonal blocks vary independently over
all square matrices, and the proportion invertible in
a block of size \(b'\) is \(\prod_{a=1}^{b'}(1-Q^{-a})\).

In \(I^\#/I\), \(A\) is isotropic and \(L'/I\) is an
isotropic complement perfectly paired with it. Indeed
\(I^\#=L+L'\) by self-duality, and the intersection in
the quotient is zero. If there is one such complement \(B\),
every other is the graph of a map \(B\to A\).
The graph is isotropic precisely when the associated
pairing on \(B\), with values in \(E_v/\mathcal O_{E_v}\),
is skew Hermitian. There are \(Q^{d_A/2}\) such pairings.
For a pair of different cyclic factors one off-diagonal
entry ranges over \(Q^{\min(\lambda_u,\lambda_v)}\) choices
and determines the other. On a diagonal factor of length
\(\lambda_u\) there are \(q^{\lambda_u}\) choices, since
the trace onto \(\mathcal O_{F_v}/\pi^{\lambda_u}\) is onto
in the unramified extension. The uniformizer is taken in
\(F_v\). This count includes every \(L'\) with intersection
\(I\).

There are at most \((i+1)^r\) types, and \(d_A\ge i\).
The count of translates is consequently at most
\(c_r^{-1}(i+1)^r Q^{(r-1/2)i}\). Lemma~\ref{lem:vacuumtail709}
gives \(|A_v|\le Q^{-Ni/2}\). The sum for \(i\ge1\)
is \(O_r(q^{-2})\).
At a split place the same lemma bounds the number of
cosets of total index exponent \(i\) by
\(16(i+1)q^{(r-1)i}\) and the coefficient by \(q^{-Ni/2}\).
Summation gives \(O(q^{-3/2})\).
\end{proof}

\begin{lemma}[Positivity at each unmarked place]
\leavevmode\label{lem:thresholdpositive816}
For every unitary representation \(R\) of \(U(S)(F_v)\)
and every \(K_v\)-fixed vector \(w\),
\[
   \int A_v(g)\langle w,R(g)w\rangle\,dg
     \ \ge\ \beta_v\|w\|^2,\qquad
   \beta_v=\frac{|\operatorname{GL}_r(\mathcal O_{E_v}/\pi)|}
                  {|\operatorname{Mat}_r(\mathcal O_{E_v}/\pi)|}>0 .
\]
At a split place the numerator and denominator use both
components. The integrals are absolutely convergent.
\end{lemma}
\begin{proof}
Locally the target lattice is a sum of \(r\) hyperbolic
planes and a line. In an unramified quadratic algebra
one can split a pair of unit diagonal entries in \(I_N\)
by choosing a unit of norm \(-1\). The residue norm is
onto and the norm lifts through successive principal-unit
quotients by the surjective residue trace. Pair the
resulting primitive isotropic vector with a vector of
pairing one. Adding an appropriate multiple of the first
makes the second isotropic by trace surjectivity.
The same operations hold in a split quadratic algebra.

Choose one isotropic line of each plane for the position
space. The resulting part of the oscillator acts on
\(r\) rows of \(S^\vee\), hence on
\(\operatorname{Mat}_r(E_v)\), by \(X\mapsto Xg\).
Include the scalar splitting factors in the unitary
action \(T(g)\) of the remaining line. Indeed
the source preserves the position polarization and the
absolute norm of its determinant is one. In the
coordinates of the source and target integral lattices,
the unit vacuum is \(\mathbf 1_C\otimes\eta\) with
\(C=\operatorname{Mat}_r(\mathcal O_{E_v})\), \(C\) of
mass one, and \(\eta\) of norm one. At an odd inert
place the translation part of the self-dual Heisenberg
lattice fixes the indicator of \(C\), and its dual
multipliers force support in \(C\), which characterizes
this vacuum line. The finite different scale can be
absorbed into the position coordinates and their duals.
At a split place this also follows from the split
matrix model by Fourier transforming \(r\) of the
\(2r+1\) rows. A lattice indicator goes to the indicator
of its dual with the corresponding unit normalization.
This covers the dyadic places, which split since the
quadratic field of discriminant \(-d\) already splits
at two and complex conjugation acts nontrivially on it.
As the full vacuum is \(K_v\)-fixed,
\(T(k)\eta=\eta\) for \(k\in K_v\).

Put \(T_w=R\otimes T\) and \(a=w\otimes\eta\). The
integral in the statement is the integral over \(X,g\) of
\[
  \mathbf 1_C(X)\mathbf 1_C(Xg)
                \langle a,T_w(g)a\rangle .
\]
Its absolute integral is at most
\(\|w\|^2\int |A_v(g)|\,dg\), since the \(X\)-integral
of the two indicators is the intersection volume multiplying
the line coefficient in \(A_v(g)\).
The full-rank matrices have full additive measure, in both
components at a split place. On them disintegrate the
additive measure over the right \(U(S)(F_v)\)-orbits in
\(\operatorname{GL}_r(E_v)\). The quotient integration
formula for Haar measure applies because additive measure there is Haar
measure times \(|\operatorname{Nm}_{E_v/F_v}\det X|_v^r\)
up to a constant, and this factor is constant on each orbit.
Both groups are unimodular.

For an orbit with representative \(Y\), the integral in
the two group variables is
\[
  \left\|\int_{U(S)(F_v)}\mathbf 1_C(Yg)\,T_w(g)a\,dg\right\|^2 .
\]
In fact expansion of the square and the substitution
\((u,v)=(u,ug)\) give exactly the preceding pairing.
On a fixed orbit the integrand of this vector integral
has compact support. In fact \(Yg\in C\) bounds the entries of \(g\),
and the unitary relation with the fixed unimodular sign
matrix bounds the entries of \(g^{-1}\) as well.
Thus every orbit contributes nonnegatively.
For any orbit meeting \(C^\times=\operatorname{GL}_r
(\mathcal O_{E_v})\), choose \(Y\in C^\times\).
Then \(Yg\in C\) precisely for \(g\in K_v\), and these
are exactly the points there with \(Yg\in C^\times\).
Its contribution is \(\|w\|^2\). Since \(K_v\) has
mass one the measure of these orbits is the additive
mass of \(C^\times\), which is \(\beta_v\) by reduction
modulo the uniformizer. This proves the inequality.
\end{proof}

\begin{theorem}[The definite auxiliary group reproduces the real integral]
\leavevmode\label{thm:flatreproduce816}
Let \(J_S\) be the principal subgroup modulo the squared
old marks on \(U(S)\), with the full lattice stabilizers
elsewhere. On \(U(V^\flat)\) use the analogous \(J_\flat\).
Use characters \((\upsilon^m,1)\), the indicator \(\delta_e\)
of the standard coordinate inclusion \(e\) of \(r\) columns
in \(N\) rows modulo the squared marks, and the unmarked
unit vacua. At infinity take in every factor the
\(\Delta'\)-th power of the minor on the first \(2m\)
rows of this inclusion and all positive source columns.
Call the product \(Q^\flat\).

For each finite class \(c\) of the definite target at
level \(J_\flat\), choose its representative \(y_c\) and
its prescribed full-rank matrix of short rows. In every
minor of \(Q^\flat\) one may choose any \(2m\) rows from
that matrix applied to the variable matrix. Denote by
\(Q_{c,\rho}\) the products over all such selections,
with a separate choice in every factor and embedding.
For any \(g_1\in U(S)(\mathbb A_F)\) there exist complex
numbers \(d_{c,\rho}\), independent of \(\psi\), such that
for every smooth rapidly decreasing \(J_S\)-fixed function
\(\psi\) on \([U(S)]\),
\[
 \int_{G_S}\psi(g_1h)H_{u_1}(h)\,dh
 =
 \sum_{c,\rho}d_{c,\rho}\int_{[U(S)]}
       \psi(x)\theta^\flat_{\delta_e,Q_{c,\rho}}(x,y_c)\,dx .
\]
The representatives have identity real component.
\end{theorem}
\begin{proof}
The real coefficient of the vector \(Q^\flat\) at
\((h,1)\) is \(\|Q^\flat\|^2 H_{u_1}(h)\). Under the two
compact blocks its weights are \((-\Delta',N)\) by the
positive full-column determinant and the \(N\) positive
target rows. On a boost its normalized coefficient is
\(\prod_i(\cosh\xi_i)^{-(N+\Delta')}\) per embedding,
since the contractions kill the positive-letter polynomial
and creation terms are orthogonal to it. These are precisely
the compact characters and boost values of \(H_{u_1}\), as
in Lemma~\ref{lem:carriercoefficient815}. This coefficient
is integrable by Lemma~\ref{lem:carrierweightmass815}.

At an old mark the point test has coefficient equal to
its squared norm on the principal subgroup and zero
elsewhere. The identity submatrix forces \(g\equiv I\)
modulo the squared mark if both \(X\) and \(Xg\) are in
its support. To see this even before \(g\) is known to
be integral, the selected submatrices of \(X\) and \(Xg\)
are invertible over the integer ring and congruent to
the identity. Their quotient is \(g\).
Right convolution therefore acts on the level space
by a strictly positive scalar at that mark.

At the other places the operators
\(U_v=\int A_v(g)R(g)\,dg\) on the level space are
self-adjoint. Here \(R\) is right translation and
\(A_v(g^{-1})=\overline{A_v(g)}\). Bi-invariance preserves
that space. Lemma~\ref{lem:thresholdpositive816} gives
\(U_v\ge\beta_v I\). For a fixed field the sum of
\(\|U_v-I\|\) is finite by Lemma~\ref{lem:thresholdtail816}.
Indeed it is bounded by a constant depending on \(r\)
times \([F:\QQ]\sum_{n\ge2}n^{-3/2}\). These operators
commute at different places. Their product converges in
operator norm to an invertible operator \(U\). To see
invertibility, put aside the finitely many factors with
\(\|U_v-I\|>1/2\). Each has a bounded inverse by the
positive lower bound. The other inverses satisfy
\(\|U_v^{-1}-I\|\le2\|U_v-I\|\), so their product
converges in norm as well. The absolute local coset sums
also give an integrable adelic coefficient.

Let \(\Theta^\flat\) be the auxiliary theta lift using
\(Q^\flat\) and the indicated finite data.
Apply the convergent doubled Siegel--Weil identity and
unfolding in the proof of Theorem~\ref{thm:complementgram812}.
Here \(N=2r+1>2r\), the target
is definite, an old split point test has full column rank,
and the real and finite coefficients together are in
\(L^1\). In particular the full-rank test eliminates
every nongraph term even for a rapidly decreasing input
that has a nonzero constant term. With a positive measure
constant \(c_*\), the identity reads
\[
  (\Theta^\flat)^*\Theta^\flat=c_*\,U\,R_\infty,\qquad
  (R_\infty\psi)(y)=\int_{G_S}\psi(yh)H_{u_1}(h)\,dh .
\]
For precision this follows first as a pairing on rapidly
decreasing test functions. The \(L^1\) convolution bounds
the pairing, so the lift extends boundedly from the dense
smooth compactly supported level space to its \(L^2\)
completion. Thus if \(\Theta^\flat\psi=0\) then
\(R_\infty\psi=0\) in \(L^2\), by invertibility of \(U\).
For the stated inputs this last convolution is continuous
by boundedness of \(\psi\) and \(L^1\)-domination. The
quotient Haar measure has full support, so it vanishes
at \(g_1\).

If the periods on the right of the theorem are all zero
then \(\Theta^\flat\psi\) vanishes everywhere. In fact the
definite target quotient has finitely many finite classes
at the given level. Its right translates at infinity
rotate the rows of \(Q^\flat\) by unitary matrices and
multiply by a scalar character. Any rotated row is a
complex linear combination of the prescribed full-rank
rows at that class. Multilinearity of the determinant
expresses each rotated minor in their \(2m\)-rowed minors.
Taking products gives exactly the span of the \(Q_{c,\rho}\).
Finite level invariance and rational invariance give the
asserted vanishing at every point.
Consequently evaluation of \(R_\infty\psi\) at \(g_1\) is
a linear functional vanishing on the common kernel of a
finite list of period functionals. It descends to their
joint image in a finite-dimensional coordinate space.
Extending the functional to that space supplies the
coefficients in the assertion.
\end{proof}

\begin{corollary}[Detection of each good-form peak by an adelic first factor]
\leavevmode\label{cor:adelicfirst816}
Fix \(g_1\), a peak of the carrier, and let \(g\) be any peak
of a nonzero good form on \(U(W_j)\) at its specified level.
For adelic \(z\in U(W_j)\) write \(\widetilde z\in U(T)\)
for its extension by the identity on the padding.
Let \(\mathcal B_i(z)\) be the following finite collection.
For \(i=(\Lambda,c,\rho)\), integrate \(\phi(x)\) on
\([U(S)]\) against the theta kernel for \(T\oplus V^\flat\)
at \((x,\widetilde z\oplus y_c)\), with characters
\((\upsilon^m,\chi^0\upsilon^{\,m+j+(t+k-1)/2-2js^8})\),
polynomial \(vQ_{c,\rho}\) on the disjoint row blocks,
and elementary residue character indexed by the matrix
\(\Lambda\). All unmarked data are the unit vacua.
There are complex \(a_i\) such that, with \(B_0\) of
Theorem~\ref{thm:integratedpeak815}, for every real
\(u\in U(W_j)(F_\infty)\),
\[
 \left|\sum_i a_i\mathcal B_i(gu)-B_0(\widetilde u)\right|
   \le 2 b_x^{-19n_*}\|u\|_{\max}^{4n_*}|B_0(1)| .
\]
In particular some \(\mathcal B_i(g)\ne0\) at every
sufficiently large node on the given intervals.
\end{corollary}
\begin{proof}
Use the residue function \(\Psi\) of
Theorem~\ref{thm:sparseapprox814} for \(g_1,\widetilde g\).
For fixed adelic \(z\) put
\(\psi_z(x)=\phi(x)\theta_{\Psi,v}(x,\widetilde z)\).
This is smooth and rapidly decreasing, since \(\phi\)
is cuspidal and the theta kernel and its derivatives have
moderate growth for fixed \(\widetilde z\). The principal
subgroup preserves every marked matrix residue and the
unmarked stabilizers fix the lattice vacua. Since the
carrier is fixed by all these subgroups, \(\psi_z\) is
\(J_S\)-invariant.
Apply the preceding theorem to \(\psi_z\).

In each of the resulting integrands the two theta kernels
multiply to the kernel for \(T\oplus V^\flat\) with residue
function \((X,Y)\mapsto\Psi(X)\delta_e(Y)\). This follows
by multiplying the two sums on the direct sum of their
rational position spaces. The oscillator vectors tensor,
and the character on the source is the product of the
two source actions, with the auxiliary determinant twist
trivial. On the two target blocks restriction of the
determinant character \(\upsilon^m\) gives the original
characters. The tensor product of the unmarked vacua
is the unit vacuum.

Expand this residue function in the additive characters
of all the matrices. They are exactly the trace characters
indexed by \(\Lambda\). On every local integer quotient
modulo the squared mark the defining additive character
is primitive, so the pairing \((x,y)\mapsto\chi(xy)\)
is nondegenerate because every nonzero principal ideal contains
the last nonzero ideal, on which the character is nontrivial.
The entrywise trace pairing and the finite product over
marks retain this property. Finite Fourier inversion gives
the required expansion with coefficients independent of
\(z\). All the sums are finite and the quotient integrals
converge absolutely by rapid decrease. It follows that
the real integral at \(z\) is \(\sum_i a_i\mathcal B_i(z)\)
simultaneously for all \(z\). At \(z=gu\) it is exactly
\(B_\Psi(\widetilde u)\) of Theorem~\ref{thm:integratedpeak815}.
Since \(\|\widetilde u\|_{\max}=\|u\|_{\max}\), this gives the
estimate and the nonvanishing at \(u=1\).
\end{proof}

These first factors have now been obtained on the original
adelic quotient. Detecting the good form itself requires
pairing its product with a first factor against the real
projection of Corollary~\ref{cor:complementreproduce812}.
The linear combination above is over \(\CC\).
\clearpage

\subsection{Holomorphy of the adelic first factors}
\label{sec:carrierholomorphy817}
Let \(\mathcal B_i\) be the first factors of
Corollary~\ref{cor:adelicfirst816}. We compute their right
derivatives on \(G_\infty=U(W_j)(F_\infty)\) from their defining carrier
integrals. Put \(a=k+j-1,\ b=m-j,\ s=a+b\),
\(p=s+9m+1,\ N_0=5s-9m-1\), and let \(H=s^8\).
The source of the theta kernel in that integral is
\(S\) of signature \((2m,m)\). Its target
\(T\oplus V^\flat\) contains \(W_j\) and has exactly \(m\)
negative rows, counting the \(j\) padded ones. This equality
with the number of negative source columns is essential below.

At a real place use the following names for matrices of Fock
variables, with target rows on the left and source columns above,
\[
\begin{array}{c|cc}
         & + & -\\ \hline
 +       & X & V\\
 -       & U & Y
\end{array}
\]
The Fock complex structure takes the original complex
coordinate in a like-sign slot and its conjugate in an
opposite-sign slot. Each of the four arrays consists of
ordinary polynomial variables in this complex structure.
\(Y\) is square of size \(m\). The archimedean polynomial
used in \(\mathcal B_i\) depends only on \(X,Y\) and is
divisible by \(\det Y\) at every real place. Indeed its
negative determinant has exponent \(p+m+\mu_\tau\ge1\)
in Corollary~\ref{cor:paddedminor814}, and the short-row
factors of Theorem~\ref{thm:flatreproduce816} use only \(X\).

\begin{lemma}[Transfer of a noncompact root derivative]
\leavevmode\label{lem:roottransfer817}
For this pair, let \(P\) be a polynomial in \(X,Y\) divisible
by \(\det Y\) at a fixed place. Let \(E^{W}_{ij}\) be the
complex root matrix with entry one in positive target row \(i\)
and negative target column \(j\), and \(E^S_{\alpha\beta}\)
the root with positive source row \(\alpha\) and negative source
column \(\beta\). There are polynomial vectors \(Q_{\alpha\beta}\)
such that
\[
     d\omega(E^{W}_{ij})P
       =\sum_{\alpha,\beta}d\omega(E^S_{\alpha\beta})
                                   Q_{\alpha\beta}.
\]
The identity holds with the prescribed determinant twists.
\end{lemma}
\begin{proof}
For any polynomial \(L\) in the like-sign variables,
the two root actions satisfy
\[
 \begin{aligned}
 d\omega(E^W_{ij})L
   &=c_1\sum_{\alpha\text{ positive}}X_{i\alpha}U_{j\alpha}L,\\
 d\omega(E^S_{\alpha\beta})L
   &=c_2\sum_{l\text{ negative}}U_{l\alpha}Y_{l\beta}L .
 \end{aligned}
\]
The constants \(c_1,c_2\) are nonzero and independent of the
indices. To verify the actions, decompose a real boost in an
opposite-sign plane into its two complex root operators. For
one spectator column or row the infinitesimal boost in Fock
coordinates is a nonzero multiple of creation in the pair
and its negative adjoint annihilation. This follows by
differentiating the two-mode boost in the proof of
Lemma~\ref{lem:oscillatorcartan710}. Its off-diagonal
operators are quadratic creation and annihilation,
and its diagonal factor has derivative zero at \(t=0\).
The two operators have opposite torus weights, so taking
the root component keeps exactly one. On the target torus
the weights of the \(X,Y\) letters are the corresponding
row weights, and those of \(U,V\) are their negatives.
For the root \(e_i-e_j\), a positive spectator column
therefore contributes multiplication by \(X_{i\alpha}U_{j\alpha}\).
A negative spectator column contributes differentiation
in \(V_{i\beta}\) and \(Y_{j\beta}\), which kills \(L\).
On the source torus the like-sign letters have the negative
of their column weights since Hom acts through the inverse
source matrix. For \(e_\alpha-e_\beta\) a negative spectator
row contributes multiplication by \(U_{l\alpha}Y_{l\beta}\).
A positive spectator row differentiates in \(X\) and \(V\)
and gives zero. The independent spectator actions add.
Permutations within each sign block give the same constants
for its indices. Determinant characters have zero derivative
on off-diagonal root matrices.

In the row order used for \(Y\) take
\[
        Q_{\alpha\beta}=(c_1/c_2)\,
              X_{i\alpha}(Y^{-1})_{\beta j}P .
\]
The inverse here denotes the adjugate divided by \(\det Y\),
so this expression is polynomial and involves only \(X,Y\).
Applying the second action and summing contracts \(Y\) with
its inverse on the negative target indices. The result is
the first action applied to \(P\), first on \(\det Y\ne0\)
and then identically as a polynomial.
\end{proof}

\begin{proposition}[The compact weights of the first factors]
\leavevmode\label{prop:firstcompact817}
At each embedding number coordinates \(0,\ldots,s-1\) in their
original order, the first \(a\) positive. In each sign block
use decreasing original index as highest-weight order.
Write \(\sigma_0,\ldots,\sigma_{s-1}\) for the decreasing
Harish--Chandra list of Theorem~\ref{thm:gluedcost811}
at this embedding. The first-factor vector is compact
highest on \(W_j\) of weight
\[
  v_i=
  \begin{cases}
   \sigma_0-\sigma_i-i,&0\le i<a,\\
   p+\mu_\tau,&a\le i<s .
  \end{cases}
\]
In particular this is the right compact weight of each
period \(\mathcal B_i\), whenever the period is nonzero.
\end{proposition}
\begin{proof}
The splitting character on the target has weight \(m\).
The inverse determinant of its action on negative Hom
coordinates contributes \(-m\) on positive rows and
\(-2m\) on negative rows. These are respectively the
numbers of opposite-sign source columns with a minus sign.
After the splitting character the vacuum weights are
therefore \(0,-m\). In like-sign letters a compact target
matrix \(k\) sends the variable row array to \(k^{\mathsf t}\)
times the array. This is the standard action on row-labelled
generators, since the linear creation operator transforms by
the target action on its basis vector. On the negative
rows the polynomial uses all rows of \(Y\) in one determinant
of power \(p+m+\mu_\tau\). Its contribution on the negative
block of \(W_j\) is that determinant character irrespective
of the ordering of the padding.

For the positive block, a minor of height index \(h<a\)
has exponent \(c^+_{h,\tau}\). The consecutive initial
\(k-1\) entries in the Harish--Chandra list make this
exponent zero for \(1\le h<k-1\). For \(k-1\le h<a\) the
number \(a-h\) is at most \(j\le m\). The \(2m\) rows
of the minor consequently include every row \(h,\ldots,a-1\)
of \(W_j\), and its other rows lie in the padding. A
triangular matrix upper in the specified reversed order
acts through its transpose on that suffix of rows, preserving
it and multiplying its determinant by the suffix diagonal
product. Hence the positive weight at \(i\) is
\(\sum_{1\le h\le i}c^+_{h,\tau}\).
Through \(\sigma_{a-1}=H-w_{j,\tau}\) every entry exceeds
\((t-1)/2-w_{j,\tau}\). Thus at \(1\le h<a\) the column
count is \(\sigma_{h-1}-\sigma_h-1\). Its sum gives
the stated formula. All strict compact upper roots kill
the vector, also when a sign block is empty. The auxiliary
short-row factors use rows outside \(W_j\).
\end{proof}

\begin{theorem}[Holomorphic weight of the first-factor periods]
\leavevmode\label{thm:firsthol817}
Let \(u\) be the weight of the conjugate of a good form on
\(W_j\) described in Theorem~\ref{thm:jointminima813},
and \(u^\circ\) the weight of the complementary real
coefficient \(K_\infty\) in Lemma~\ref{lem:negativecoefficient812}.
At any adelic point \(z\) the function
\(h\mapsto\mathcal B_i(zh)\) is holomorphic of weight
\(u-u^\circ\) for the order above, with the positive
coordinates first. In particular
\(h\mapsto\mathcal B_i(zh)K_\infty(h)\) is holomorphic
of the same weight as the conjugate good form.
\end{theorem}
\begin{proof}
Differentiating the integral defining the first factor
in an upper noncompact root direction of \(W_j\) replaces
its polynomial by \(d\omega(E^W)P\). By
Lemma~\ref{lem:roottransfer817} this gives a sum of integrals
of the form \(\int_{[U(S)]}\phi(x)R_{E^S}\theta_Q(x,y)\,dx\).
Right invariance of measure turns this into
\(-\int_{[U(S)]} (R_{E^S}\phi)(x)\theta_Q(x,y)\,dx=0\).
Indeed Theorem~\ref{thm:carrier813} makes \(\phi\)
holomorphic with the positive coordinates first, so every
such upper root kills it. For real derivatives the
integration by parts follows by right translating under
the integral and differentiating at the identity.
The carrier and its derivatives decrease rapidly, and
the theta kernel and its derivatives have moderate growth
uniformly for the translates in a compact neighborhood.
This justifies the differentiations, and complex linearity
gives the root identity. Compact upper roots already kill
the polynomial by Proposition~\ref{prop:firstcompact817}.

We identify the resulting weight \(v\). Write \(\delta=j-m/2\).
On a positive index \(u^\circ_i=\delta\). On a negative
index \(i=s-1-l\), \(0\le l<b\), the complementary weight is
\[
     u^\circ_i=\delta+N_0+\sum_{h=l+1}^b c^-_{h,\tau}.
\]
This follows from the determinants on the ordered negative
columns in Lemma~\ref{lem:negativecoefficient812}.
In its block formula the Schur complement \(S_-=D-CA^{-1}B\)
of the positive block transforms under right block upper
matrices by \(S_-\mapsto S_-b_{--}\). In particular its
leading minors transform by the corresponding upper
triangular characters. This also proves that \(K_\infty\)
extends holomorphically with the stated weight.

The conjugate Harish--Chandra list \(-\sigma_i\)
decreases in reversed index order within each sign block.
For the weight \(u=-\sigma+\rho_{\rm n}-\rho_{\rm c}\),
the half-sum of positive noncompact roots has entries
\(-b/2\) in positive coordinates and \(a/2\) in negative
ones. The compact half-sum at a positive \(i\) is
\(i-(a-1)/2\), and at the \(l\)-th negative coordinate it
is \((b-1)/2-l\). Therefore
\[
 u_i=
 \begin{cases}
   -\sigma_i-i+(a-b-1)/2,&i<a,\\
   -\sigma_i+l+(a-b+1)/2,&i=s-1-l\ge a .
 \end{cases}
\]
At a positive index the equality \(u_i=v_i+u^\circ_i\)
now follows from \(\sigma_0+\delta=(a-b-1)/2\).
When \(b>0\), put \(b'=(t-1)/2\) and suppress \(\tau\)
from \(w_j\). For \(1\le h<b\) the negative count is
\(\sigma_{s-h-1}-\sigma_{s-h}-1\), since both entries
are at most \(\sigma_a=-w_j-H\). At \(h=b\) the upper
endpoint is truncated to \(-w_j-b'\), so that
\(c^-_{b,\tau}=H-b'-1=\Delta'\). Thus the sum in
\(u^\circ_{s-1-l}\) is \(-w_j-b'-\sigma_i-b+l\).
Using \(\mu_\tau=w_j-j-(t+k-1)/2\) from
Lemma~\ref{lem:carriercoefficient815}, \(p+N_0=6s\),
\(t=5s\), and \(k+m=s+1\) gives \(u_i=v_i+u^\circ_i\)
on the negative block as well. For \(b=0\) the positive
calculation covers the whole group.

We recall why the differential equations give holomorphy
in the sense used here. For a fixed translate write its
value as \(F(h)\), let \(B\) be the ordered complex upper
triangular group, and extend to the open set \(G_\infty B\)
by \(F(hb)=F(h)\chi_v(b)\), using the diagonal character
of weight \(v\). The intersection of \(G_\infty\) and
\(B\) is the compact diagonal torus, on which \(F\)
has the stated character. The extension is smooth by
local Hermitian Gram--Schmidt on the flags with positive
first \(a\)-plane. It is holomorphic on each right \(B\)
fiber. At a real group point its right differential on
the complexified real Lie algebra agrees, on every upper
triangular direction, with the derivative of this character
times \(F\). The strict roots and the diagonal weights give
these equations. To check the Cauchy--Riemann condition,
the derivative of the extension on such a direction is
the sum of its holomorphic derivative there and its
antiholomorphic derivative on the conjugate direction
for the unitary real form. The holomorphic derivative
already equals the character derivative by the fiber law,
so the antiholomorphic derivative vanishes on those
conjugate directions as well. The upper triangular algebra
and its unitary conjugate, the lower triangular algebra,
span the complex matrix algebra. Thus the antiholomorphic
differential vanishes. The fiber transformation carries
this condition to all of \(G_\infty B\), proving the claim
at every real place.
\end{proof}
\clearpage

\subsection{Separation by products of theta periods}
\label{sec:complexseparation818}
Fix the previous fields, levels, and \(1\le j\le m\).
We prove that a nonzero good form \(F_0\) on \(U(W_j)\)
has a nonzero projected product period using one of the
carrier integrals \(\mathcal B_i\) of
Corollary~\ref{cor:adelicfirst816}. The peak estimates will
be used for a circle of points on the same field.

Write \(G_\infty=\prod_{\tau}U(a,b)\), where \(a=k+j-1\),
\(b=m-j\), and \(s=a+b\). Within the positive and
negative blocks use the orders of Theorem~\ref{thm:firsthol817},
with the positive block first. At one place let \(D_l(h)\)
be the principal minor on the first \(l\) coordinates
in this order. For an integral weight \(w\) decreasing
inside each compact block put
\[
     M_w(h)=(\det h)^{w_s}
                 \prod_{l=1}^{s-1}D_l(h)^{w_l-w_{l+1}},
\]
numbering entries of the ordered list from \(1\) to \(s\).
The only proper minor that can have a negative exponent is
\(D_a=\det A\), where \(A\) is the positive square block of
\(h\). It is invertible by the unitary relations. This
defines a continuous function on the group and a holomorphic
function on \(U(a,b)B\), where \(B\) is complex upper
triangular. Its right character on \(B\) and its left
character on the compact diagonal torus are the product of
diagonal entries to the weights \(w_i\). For a family of
weights use the product of these functions at all embeddings,
and write \(\chi_w\) for the product character.

\begin{lemma}[The circle average of a holomorphic group function]
\leavevmode\label{lem:flagmean818}
At all real places let \(d(t)\) have successive diagonal
entries \(1,t,\ldots,t^{s-1}\) in the indicated order,
where \(|t|=1\). If \(L\) extends holomorphically with
right upper triangular weight \(w\), then
\[
   \int_{|t|=1}\frac{L(d(t)h)}{\chi_w(d(t))}\,dt
         =L(1)M_w(h)
\]
for normalized circle measure. In particular, whenever
\(L\overline{M_w}\) is integrable and \(M_w\in L^2\),
\[
       \int_{G_\infty}L(h)\overline{M_w(h)}\,dh
               =L(1)\int_{G_\infty}|M_w(h)|^2\,dh .
\]
\end{lemma}
\begin{proof}
First suppose all the leading minors of \(h_\tau\) are
nonzero. Factor each matrix as \(n_\tau b_\tau\) with
\(n_\tau\) lower unitriangular and \(b_\tau\) upper
triangular. Ratios of consecutive leading minors are
the diagonal entries of \(b_\tau\), so \(\chi_w(b)=M_w(h)\).
The lower triangular matrix
\(n(z)=d(z)n d(z)^{-1}\) has polynomial entries for \(z\ne0\)
and extends with \(n(0)=I\). It lies in the holomorphic
flag domain for \(|z|\le1\). Indeed that domain consists
of matrices whose first \(a\) columns span a positive
plane, by Hermitian Gram--Schmidt in the given order.
The first \(a\)-plane of \(n_\tau\) equals that of \(h_\tau\),
the graph of a contraction \(Z_\tau\) of operator norm
less than one. For \(0<|z|\le1\) its new graph is
\(d_-(z)Z_\tau d_+(z)^{-1}\), of norm at most
\(|z|\|Z_\tau\|\) when both blocks occur. At zero the
plane is the positive coordinate plane. For \(b=0\)
the domain condition holds everywhere on the invertible
matrices.

Thus \(z\mapsto L(n(z))\) is holomorphic in a neighborhood
of the closed disk. On its boundary the right character
law gives \(L(d(t)h)/\chi_w(d(t))=M_w(h)L(n(t))\).
The mean value formula proves the first identity. The
nonvanishing condition on the minors is dense, since the
individual minors are real analytic on the connected group
and equal one at its identity. Continuity extends the
identity to every \(h\). Finally average \(L(h)\overline{M_w(h)}\)
under left translation by the same circle. The left
character of \(M_w\) cancels the factor \(\chi_w(d(t))\).
Haar invariance and the assumed absolute integrability
justify the exchange of the two integrals.
\end{proof}

\begin{lemma}[The complementary mass relative to the good-form weight]
\leavevmode\label{lem:flagmass818}
Use the weights \(u,u^\circ\) of Theorem~\ref{thm:firsthol817}
and put \(J_u=\int_{G_\infty}|M_u(h)|^2\,dh\). Then
\(M_{u^\circ}=K_\infty\), \(0<J_u<\infty\), and, uniformly
in \(j\) at large nodes on the original intervals,
\[
    \int_{G_\infty}|K_\infty(h)|\,\|h\|_{\max}^{4n_*}\,dh
                   \le b_x^{n_*}J_u .
\]
\end{lemma}
\begin{proof}
By Lemma~\ref{lem:negativecoefficient812} the complementary
coefficient at each place is the product of a determinant
character, a power of \(\det S_-\), and leading minors of
\(S_-=D-CA^{-1}B=(D^*)^{-1}\). Block elimination gives
\(D_{a+l}(h)=\det A\,D_l(S_-)\) and
\(\det h=\det A\det S_-\). Substitution of the weights
in Theorem~\ref{thm:firsthol817} gives \(K_\infty=M_{u^\circ}\).
For \(b>0\), \(\|S_-\|\le1\), and the negative count
at height \(b\) is \(\Delta'=s^8-(t+1)/2\). Therefore
\[
   |K_\infty(h)|\le\prod_\tau|\det A_\tau|^{-(N_0+\Delta')},
   \qquad N_0=5s-9m-1 .
\]
For \(b=0\) it is a unitary character. For the good
weight, every proper minor \(D_l(h_\tau)\) has modulus
at most \(|\det A_\tau|\). For \(l<a\) the product of
the \(l\) largest singular values of \(A_\tau\) bounds
the minor, and all the other singular values are at least
one. For \(l>a\) use block elimination and contraction
of \(S_-\). Thus for \(b>0\),
\(|M_u(h_\tau)|\le|\det A_\tau|^{u_1-u_s}\) in the ordered
list at this place. All the nonboundary powers are
nonnegative, and the boundary factor is exactly \(\det A_\tau\).
The gap \(u_s-u_1\ge2s^8-2s\) was computed in the proof
of Theorem~\ref{thm:jointminima813}. When \(b=0\), the
group is compact and \(|M_u|\le1\).

Use on \(G_\infty\) the product Haar measure, with the
compact factors in the matrix ball decomposition of mass
one. Put \(\mathcal D=\{Z\in M_{b,a}(\CC):\|Z\|<1\}\).
At a place, the measure in the chart \(Z=CA^{-1}\) is
\(c_\tau\delta^{-s}\,dZ\,dk\), where
\(\delta=\det(I-ZZ^*)=|\det A|^{-2}\). The Jacobian
calculation in Lemma~\ref{lem:carrierweightmass815} applies
with these two block sizes. When \(b=0\) give the one-point
ball volume one. The last bounds on \(M_u\) prove
square integrability. Also
\(\|h\|_{\max}\le2\prod_\tau|\det A_\tau|\), by the singular
values of the Cartan boosts. Since \(N_0+\Delta'-4n_*\ge2s\)
at every sufficiently large node, in both cases
\[
    \int |K_\infty(h)|\|h\|_{\max}^{4n_*}\,dh
         \le 2^{4n_*}\prod_\tau c_\tau\operatorname{vol}(\mathcal D).
\]

For the lower bound put \(\varepsilon=[1000s^2(1+b_x)]^{-1}\).
Take \(Z\in\varepsilon\mathcal D\) and \(\|k-I\|\le\varepsilon\).
The group representative of \(Z\) with positive Hermitian
diagonal blocks has these blocks \((I-Z^*Z)^{-1/2}\) and
\((I-ZZ^*)^{-1/2}\). Its lower left block is \(Z(I-Z^*Z)^{-1/2}\)
and its upper right block is the adjoint of the lower left.
It consequently differs from the identity by at most
\(3\varepsilon\) in operator norm. After right multiplication
by \(k\) the bound is \(4\varepsilon\). Every leading
submatrix then has singular values in
\([1-4\varepsilon,1+4\varepsilon]\), giving
\(\big|\log |D_l|\big|\le8s\varepsilon\).
All weights of \(u\) have absolute value at most \(5b_x\),
by Theorem~\ref{thm:jointminima813}. Summing the absolute
exponents shows that the local \(|M_u|^2\) is at least
\(e^{-1}\) on this set.

The indicated compact neighborhood has mass at least
\((\varepsilon/3)^{2s^2}\). Indeed a maximal
\(\varepsilon\)-separated subset of the compact group has
at most \((1+2/\varepsilon)^{2s^2}\) elements, by packing
operator-norm balls of radius \(\varepsilon/2\) in the
ambient complex matrices. Its translates of the given
neighborhood cover the group. The smaller matrix ball
has fraction \(\varepsilon^{2ab}\) of the full Euclidean
volume and \(\delta^{-s}\ge1\). It follows that \(J_u>0\)
and the mass ratio in the assertion is at most
\[
        2^{4n_*}e^f(3/\varepsilon)^{(2s^2+2ab)f},
               \qquad f=[F:\QQ].
\]
Here \(\log(3/\varepsilon)=O(\log b_x)\) and
\(fs^2/n_*=s^2/(2t^3)\to0\). The logarithm of the last
bound is at most \(n_*\log b_x\) for all sufficiently
large nodes, uniformly in the block sizes.
\end{proof}

\begin{theorem}[Separation of good forms by the product periods]
\leavevmode\label{thm:complexseparation818}
At every sufficiently large node of the fixed-\((q_0,d)\)
intervals, let \(F_0\) be a good form on \(U(W_j)\) at the
prescribed level and on its specified real line. If
\[
 \int_{[U(W_j)]} F_0(x)\mathcal B_i(x)
                  \theta^\circ_{\delta_e,Q}(x,y)\,dx=0
\]
for all first-factor indices \(i\), all complementary
finite class representatives \(y\), and the complementary
negative-column tests \(Q\) of
Corollary~\ref{cor:complementreproduce812}, then \(F_0=0\).
In particular every nonzero good form has
\(P_\infty(F_0\mathcal B_i)\ne0\) for some \(i\).
\end{theorem}
\begin{proof}
Each \(F_0\mathcal B_i\) is a smooth rapidly decreasing
input at level \(J\) of Section~\ref{sec:complementperiod812}.
Indeed the carrier is rapidly decreasing and the kernel
in \(\mathcal B_i\) has moderate growth jointly in its
two group variables. Integrating in the carrier variable
leaves a function of moderate growth, also after real
differentiation. The finite data are fixed by the prescribed
principal subgroup at the old marks and the full stabilizers
elsewhere, and the good form is cuspidal. The common-kernel
assertion of Corollary~\ref{cor:complementreproduce812} gives
\[
   \int_{G_\infty} F_0(xh)\mathcal B_i(xh)K_\infty(h)\,dh=0
                 \qquad\text{for all }x,i .
\]
The convolution is continuous because the input is bounded
and \(K_\infty\in L^1\). Thus vanishing in \(L^2\) gives
the pointwise statement.

Suppose \(F_0\ne0\), write \(F=\overline{F_0}\), and choose
a peak \(z\) of \(F\). The function \(F(zh)\) is holomorphic
of weight \(u\) by Theorem~\ref{thm:jointminima813}. Every
\(z_t=z d(t)\) is another peak since the compact torus acts
on the right by a unitary character. Fix a carrier peak \(g_1\).
The bounded function \(B_0\) of Theorem~\ref{thm:integratedpeak815}
depends on this carrier peak and the polynomial, and its
definition uses no good-form peak. Hence the same function
\(D(h)=B_0(\widetilde h)\), with \(D(1)\ne0\), occurs in
Corollary~\ref{cor:adelicfirst816} at every \(z_t\). For each
\(t\) choose a combination
\(L_t(h)=\sum_i a_i(t)\mathcal B_i(z_t h)\) satisfying
\[
    |L_t(h)-D(h)|\le\eta\|h\|_{\max}^{4n_*},\qquad
    \eta=2b_x^{-19n_*}|D(1)|.
\]
Set \(L=L_1\), where \(1\) is the identity circle element.
With \(W_t(h)=\overline{F(z_t h)/\chi_u(d(t))}\), the
convolution identity gives
\(\int W_t(h)L_t(h)K_\infty(h)\,dh=0\).
Since \(\sup|F|<2|F(z)|\) and
\(|L-L_t|\le2\eta\|h\|_{\max}^{4n_*}\),
Lemma~\ref{lem:flagmass818} gives, for each \(t\),
\[
   \left|\int W_t(h)L(h)K_\infty(h)\,dh\right|
      \le4\eta|F(z)|b_x^{n_*}J_u .
\]

The left member inside the absolute value now uses the
fixed \(L\), and its integrand is jointly integrable in
\(h,t\). This follows from boundedness of \(W_t,D\) and
the mass estimate. Thus the next averaging step needs
no measurable choice of the coefficients \(a_i(t)\).
Lemma~\ref{lem:flagmean818} gives the mean of \(W_t(h)\)
as \(\overline{F(z)M_u(h)}\).
By Theorem~\ref{thm:firsthol817} the product
\(L(h)K_\infty(h)\) is holomorphic of weight \(u\).
Its pairing with \(\overline{M_u}\) is integrable since
\(|M_u|\le1\) by the proof of Lemma~\ref{lem:flagmass818}.
The second assertion of Lemma~\ref{lem:flagmean818} therefore
evaluates that pairing as \(L(1)J_u\). Averaging the
displayed bound and cancelling \(|F(z)|J_u>0\) gives
\(|L(1)|\le4\eta b_x^{n_*}\).

On the other hand \(|D(1)-L(1)|\le\eta\). Dividing by
\(|D(1)|>0\) now gives
\[
        1\le 2(4b_x^{n_*}+1)b_x^{-19n_*},
\]
which fails for all sufficiently large nodes. This
contradiction proves separation, including for \(j=m\)
when the real group is compact.
\end{proof}

\begin{corollary}[A nonsingular matrix of complex periods]
\leavevmode\label{cor:complexminor818}
If the good space has dimension \(d_j>0\), then for any
complex basis one can select \(d_j\) of these product
periods whose matrix on that basis has nonzero determinant.
The complementary polynomials can be taken to be products
of minors with integral coefficients.
\end{corollary}
\begin{proof}
The list of first factors and of finite classes is finite.
By multilinearity, every compact translate of the indicated
complementary polynomial is a combination of products of
minors on the prescribed source columns and choices of
target rows. These products have integral coefficients and
there are finitely many of them. The theorem implies that
the joint evaluation matrix has column rank \(d_j\), so it
has the required nonzero square minor.
\end{proof}

The determinant obtained here is a complex determinant for
the periods involving the carrier integrals. An estimate for
the arithmetic determinant also requires their comparison to
the integral periods of the definite single lift, with the
normalizing factors at the old and new marks retained under
the twists.
\clearpage

\subsection{Integral coordinates of the definite first-factor polynomial}
\label{sec:stackedcoordinates819}
The periods in Theorem~\ref{thm:complexseparation818} are
defined with the carrier. We examine the polynomials used
to compare them to the definite single lift. The size of
a column in these polynomials is nearly the full rank.
Taking complementary minors makes the relevant denominator
depend only on the number of boxes missing from full columns.

For a \(p\)-by-\(p\) matrix \(X\) write
\(\Delta_{I,h}(X)=\det X_{I,[h]}\), where \([h]\) is its
first \(h\) columns and \(I\) has \(h\) elements in increasing
order. Fix \(c_h\) columns of each height \(1\le h<p\).
As in Lemma~\ref{lem:dualcartan811}, \(C(c)\) is the integral
span of products of these flag minors with the prescribed
numbers, and \(D(c)=\operatorname{Hom}_{\ZZ}(C(c),\ZZ)\).
Use the canonical map from the Weyl lattice to flag polynomials,
\[
 T_c\ell(X)=\ell_Y K_c(Y,X),\qquad
 K_c(Y,X)=\prod_h
       \det\!\left(Y_{:,[h]}^{\mathsf t}X_{:,[h]}\right)^{c_h}.
\]
Cauchy--Binet expresses \(K_c\) in \(C(c)\otimes C(c)\).
Over the real or complex numbers it defines a positive
Hermitian metric by taking complex conjugation in one
argument. Products over embeddings use tensor products
and distinct variables.

\begin{lemma}[A denominator from the complementary diagram]
\leavevmode\label{lem:complementfischer819}
Let \(c_{\tau,h}\) be the diagrams at the embeddings, put
\(\widehat c_{\tau,p-h}=c_{\tau,h}\), and
\(d=\sum_{\tau,h}(p-h)c_{\tau,h}\). There is a positive
integer \(H_c\le\max(1,d)^{3d}\) depending only on these
diagrams and \(p\) such that for any subring \(A\subset\CC\)
and any \(P\in C(c)\otimes A\), the unique complex vector
\(T_c^{-1}P\) belongs to \(H_c^{-1}(D(c)\otimes A)\).
The same statement holds with coefficients polynomial
in auxiliary variables over \(A\).
\end{lemma}
\begin{proof}
First take an arbitrary diagram \(\nu\) in one block.
Let \(m_\nu(X)=\prod_h\Delta_{[h],h}(X)^{\nu_h}\).
The bilinear Fischer pairing is
\([F,G]=F(\partial)G(0)\), with the differential operator
formed from the same matrix entries. Monomials are
orthogonal and have self-pairing the product of their
factorials. If \(G\in C(\nu)\) then
\[
       [m_\nu,G]=[m_\nu,m_\nu]G(I).
\]
Indeed by diagonal row scaling only products of the
same row multidegree as \(m_\nu\) can contribute. Row one
must occur in every factor. Once the rows strictly
before \(k\) have been placed, row \(k\) must occur in
every factor of height at least \(k\). This inductively
forces the principal product. Every other product also
vanishes at \(I\).

Row substitution by \(Y^{\mathsf t}\) is adjoint for this
pairing to substitution by \(Y\), by the chain rule.
Hence \([K_\nu(Y,\cdot),G]=[m_\nu,m_\nu]G(Y)\).
For \(\ell\) equal to evaluation at \(Y\), the functional
\([T_\nu\ell,\,\cdot\,]/[m_\nu,m_\nu]\) is exactly \(\ell\).
Evaluations span the dual over \(\CC\), since a polynomial
vanishing at all matrices is zero. Thus \(T_\nu\) is
invertible there and its inverse has this normalized
pairing formula. In particular the positive integer
\([m_\nu,m_\nu]\) clears the denominator on \(C(\nu)\).

Let \(J_p\) reverse the \(p\) coordinates and \(w=\sum_h c_h\).
On invertible matrices put
\[
        (S_cF)(X)=\det(X)^w F(X^{-\mathsf t}J_p).
\]
The complementary-minor identity sends each factor,
after its one determinant multiplier, to a signed flag
minor on the complementary rows and first \(p-h\) columns.
It follows that \(S_c\) maps \(C(c)\) onto \(C(\widehat c)\).
The inverse is integral by the same formula, whose double
application is a sign. Furthermore applying \(S_c\) in
both variables to \(K_c\) gives \(K_{\widehat c}\).
Indeed the new Gram matrix before the determinant multipliers
is \(J_p^{\mathsf t}(Y^{\mathsf t}X)^{-\mathsf t}J_p\);
its leading \(h\)-minor is the leading complementary
\((p-h)\)-minor of \(Y^{\mathsf t}X\) divided by
\(\det(Y^{\mathsf t}X)\). Thus \(T_{\widehat c}=S_c T_c S_c^*\),
with \(S_c^*\) the integral dual map. We can use the
denominator of the complementary pairing.

For all embeddings take \(H_c=\prod_\tau
[m_{\widehat c_\tau},m_{\widehat c_\tau}]\).
The product polynomial in these independent variables has
total degree \(d\). The sum of the absolute values of its
coefficients is at most \(\prod_{\tau,h}((p-h)!)^{c_{\tau,h}}
\le d!\) by the determinant expansion. The self-pairing
is at most \(d!\) times the square of this sum. This proves
the asserted bound, including \(d=0\). The constructions
are integral linear maps after multiplication by \(H_c\),
so extension of coefficients gives the last two assertions.
\end{proof}

Return to \(p=s+9m+1\), \(a=k+j-1\). Besides the \(a\)
positive columns of \(W_j\), use the added columns from
\[
       I_{3m-j}\oplus I_{6m+1}\oplus I_{3m}
\]
in that order. In total these two sets have \(p+2m\)
columns. For each positive diagram column of height \(h<a\)
at \(\tau\) put \(2m\) constant rows above the variable
matrix \((X_\tau\mid Z_\tau)\) of \(p\) rows. Numbering
from zero, constant row \(l\) selects coordinate \(h+l\)
of \(W_j\) if \(h+l<a\), and otherwise coordinate \(l\)
of the first added summand. For a column of height \(a\)
use instead the \(2m\) selected rows of the short matrix
\(\tau(Y)\) of the definite class in \(I_{6m+1}\), supported
in the second added summand. There are \(c^+_{h,\tau}\) columns
of the former kind and \(\Delta'\) of the latter. Indeed the
positive count at \(a\) is
\((H-w_j)-((t-1)/2-w_j)-1=\Delta'\), by the list in
Theorem~\ref{thm:gluedcost811}, and subsequent counts vanish.
Here \(H=s^8\). Denote the product of these square determinants
at all embeddings by \(P(X,Z;Y)\), using one fixed column
ordering. Write \(c\) for the full list of these column
multiplicities.

\begin{proposition}[The coefficient vector in the added columns]
\leavevmode\label{prop:stackednorm819}
As a polynomial in \(X\), \(P\) belongs to \(C(c)\) with
coefficients polynomial in \(Z\) and the conjugate entries
\(\tau(Y)\). Its associated vector \(\zeta=T_c^{-1}P\) is
homogeneous of degree \(d=\sum_{\tau,h}(p-h)c_{\tau,h}\)
in \(Z\). Put \(q_*=\sum_{\tau,h}c_{\tau,h}\).
At any complex specialization where the entries of \(Z\)
and \(\tau(Y)\) have modulus at most \(M^4\), one has
\[
          \|\zeta\|^2\le M^{9(12m+1)q_*},
                 \qquad M=p_0p_1 .
\]
The metric is the one defined by \(K_c\).
\end{proposition}
\begin{proof}
In a factor of height \(h<a\) all the columns \(h,\ldots,a-1\)
are selected by distinct coordinate rows, since \(a-h\le j\)
for every nonzero count. Subtract multiples of these rows
to clear the corresponding entries in the bottom \(p\)
rows. None of the other columns changes in this operation.
For height \(a\) there are no columns to clear. The first
\(h\) columns are now the first \(h\) columns of \(X_\tau\)
with \(2m\) zero entries appended above each. The remaining
matrix \(D\) of \(p+2m-h\) columns has entries among
\(0,1\), the entries of \(Z_\tau\) and \(\tau(Y)\).
Expanding on the first \(h\) columns gives a sum of the
flag minors \(\Delta_{I,h}(X_\tau)\). Their coefficients
are signed maximal minors of \(D\), taking all the top
rows and the \(p-h\) bottom rows outside \(I\).
Every contributing monomial has exactly \(p-h\) entries
of \(Z_\tau\), one from each remaining bottom row.
This gives the integral polynomial statement, and taking
products gives the degree. The inverse \(T_c^{-1}\) acts
only on the \(X\)-coefficients.

For one factor the single-column Weyl metric is the
usual exterior-power metric, as follows from Cauchy--Binet
for \(K_c\). The sum of the squared moduli of the indicated
minors of \(D\) is at most \(\det(D^*D)\), the sum over all
row selections. The latter is at most the product of the
squared column lengths. Each length squared is at most
\((p+2m)M^8\le M^9\) for large nodes, and
\(p+2m-h\le12m+1\) since \(h\ge k-1=s-m\).

To pass to the product, tensor the exterior-power spaces
of all its columns and embeddings. Evaluating the first
\(h\) columns of \(X_\tau\) gives their wedge in each
factor. The inner product of the resulting tensors at
two matrices is \(K_c\) with conjugation in the first
argument. On their span this identifies the tensor metric
with the metric defined by \(K_c\). Representing a product
of linear functionals amounts to restricting their tensor
to this span, so its squared norm is bounded by the product
of the individual squared norms. This gives the estimate.
\end{proof}

\begin{theorem}[A common arithmetic denominator on the prescribed family]
\leavevmode\label{thm:stackedheight819}
On the same fields and intervals, at each \(j\) a single
positive integer \(H_c\) works for all definite classes,
row choices, old residue matrices and tame twists. Every
coordinate of \(H_c\zeta(Z;Y)\) in the Weyl lattice is a
polynomial of degree \(d\) in the added columns with
coefficients in the ring generated over \(\ZZ\) by the
\(\tau(Y)\). One has
\[
     d\le(10m+1)q_*,\qquad
     \log H_c\ll_{q_0,d_{\rm K}} f m b_x\log b_x ,
\]
where \(d_{\rm K}\) is the fixed positive quadratic
discriminant parameter of the field family. The short-row
matrices may all be chosen with entries of modulus
\(\le M^4\) in every embedding. At any specialization
with the same bound on the added entries, the squared
Weyl norm of \(H_c\zeta\) is at most
\(\exp(O_{q_0,d_{\rm K}}(f m b_x\log b_x))\).
These assertions about size apply to every conjugate
for which the indicated specialization bounds hold.
Multiplication by the prescribed added residue datum
preserves both the integrality and this estimate.
\end{theorem}
\begin{proof}
Apply Lemma~\ref{lem:complementfischer819} over the
coefficient ring of Proposition~\ref{prop:stackednorm819}.
In every nonzero diagram column \(p-h\le10m+1\).
The total column count is \(\sum_\tau\mu_\tau=O(fb_x)\)
by Lemma~\ref{lem:carriercoefficient815} and the weight
bound in Theorem~\ref{thm:gluedcost811}. Thus
\(\log H_c\le3d\log\max(1,d)=O(fmb_x\log b_x)\).
Here \(f\ll x^{12}\), \(m\ll\sqrt x\), and \(b_x\ge x^{400}\).
All dimensions and multiplicities used for \(H_c\) are
fixed before the choices and twists. The norm estimate
follows from Proposition~\ref{prop:stackednorm819} and
\(\log M\le22\log b_x\).

To obtain the short rows, move the self-dual lattice of
\(I_{6m+1}\) by the finite representative of its class.
The Hermitian length of any nonzero lattice vector is
a totally positive integer of \(F\), since the finite
lattice is integral. Its trace is at least \(f\) by the
arithmetic--geometric mean inequality and the positive
integral norm. Thus its Euclidean length over all
embeddings is at least one. As in
Lemma~\ref{lem:sparsedistance814}, the shortest Euclidean
dual length is at least \(2/(q_1d_{\rm K})\), using the
absolute different and Hermitian self-duality. The
Euclidean transference inequality in that lemma supplies
\(2f(6m+1)\) independent lattice vectors of length at
most \(f(6m+1)q_1d_{\rm K}\). They span over \(E\).
Choose an \(E\)-independent subset of \(6m+1\) and use
these vectors as the rows of \(Y\). The bound is polynomial
in \(x\), since \(q_1\le8 f^2\), and is smaller than
\(M^4\) at large nodes. The Euclidean bound covers one
embedding of each conjugate pair and thus every embedding
in absolute value. The row construction can be used at
all the twists. At a finite place the polynomial conclusion
applies in particular to any integral specialization of
its row and added entries in the indicated frames.

Finally let \(R=\mathcal O_E/\mathfrak a^2\) be the residue
ring of the old marks, and allow the modulus of the added
matrix to include the new marked primes of the twist.
Its first two added blocks are required to be the
prescribed inclusion columns in that modulus. If this
test holds, its residue value is
\[
      \sum_{\gamma\in\operatorname{GL}_{3m}(R)}
          \psi_R\!\left(\operatorname{tr}(\Lambda Z_{\rm last}\gamma)\right),
\]
where \(Z_{\rm last}\) is the last \(3m\) columns reduced
to \(R\) and \(\psi_R\) is the primitive additive character.
It is zero if the inclusion test fails. Each summand is a
root of unity. Thus the value is an algebraic integer and
its modulus in every conjugate is at most
\(|R|^{9m^2}=M^{18fm^2}\). Multiplication adds
\(O(fm^2\log M)\) to the logarithmic bound for the squared
norm. This is absorbed in the stated estimate.
\end{proof}

This result controls the coefficients inserted in the
definite theta kernels. To obtain the integral period
comparison one must also control the arithmetic evaluations
of those kernels and the scalar relating their class sums
to the carrier periods.
\clearpage

\subsection{Fourier coefficients at the twist primes and definite theta values}
\label{sec:markedfourier820}
Fix a field and an index \(j\) from Theorem~\ref{thm:stackedheight819}.
Set \(p=s+9m+1\), \(t_j=9m+1-j\) and \(M=p_0p_1\).
Write \(R=\OO_E/\mathfrak a^2\) for the ring of the two old
marks. For the ideal \(\mathfrak q\) of an allowed tame twist put
\[
 S=\OO_E/(\mathfrak a^2\mathfrak q)=R\times R_{\mathfrak q},
 \qquad R_{\mathfrak q}=\OO_E/\mathfrak q .
\]
The equality is the Chinese remainder isomorphism. For the unit
ideal the second additive group has one element. Put
\(G=\operatorname{GL}_{3m}(R)\). The following calculation keeps the
orbit multiplicities in the trace characters. This matters even
for the old matrix \(\Lambda=0\).

For a finite additive group \(A\), let \(\widehat A\) be its
group of complex characters, and use the Fourier convention
\[
 f(x)=\sum_{\lambda\in\widehat A}\widehat f(\lambda)\lambda(x),
 \qquad
 \widehat f(\lambda)=|A|^{-1}\sum_{x\in A}f(x)\overline{\lambda(x)} .
\]
The three positive definite tests have the following residue
functions. The first is the unmarked constant vacuum. The second
is the point indicator on \(S^{p\times j}\) at the prescribed
inclusion, multiplied in its argument by a unit of norm minus
one in the comparison convention. Denote the resulting point by
\(I_1\). The third is a function of
\((y,z)\in A\times B\), where
\(A=S^{p\times t_j}\) and \(B=S^{p\times 3m}\). It is
\[
  {\bf1}_{y=I_2}\sum_{g\in G}\chi_g(z),\qquad
  \chi_g(z)=\psi_R\!\left(\operatorname{tr}(\Lambda\bar z g)\right).
\]
Here \(I_2\) fixes the inclusion on the two constrained blocks
of \(3m-j\) and \(6m+1\) columns, \(\bar z\) means reduction
to \(R\), and \(\psi_R\) is the prescribed additive character.

\begin{proposition}[Exact Fourier coefficients of the three residue tests]
\leavevmode\label{prop:markedfourier820}
The Fourier coefficients of the point test and the added test are,
respectively,
\[
   \frac{\overline{\lambda(I_1)}}{|S|^{pj}},\qquad
   \frac{\overline{\alpha(I_2)}}{|S|^{pt_j}}\,
        \#\{g\in G:\chi_g=\beta\}
      \quad(\alpha\in\widehat A,\ \beta\in\widehat B).
\]
The sums of the absolute values of the three coefficient rows,
at every complex embedding of their fields of definition, are
\(1,1,|G|\). The tensor product of the three rows has the same
sum \(|G|\). Its entries become algebraic integers after
multiplication by \(|S|^{p(9m+1)}\). At any prime away from
\(N_{E/\QQ}\mathfrak q\) it suffices to multiply by the single
old integer
\[
       D_{\rm old}=M^{2fp(9m+1)}.
\]
It works simultaneously for all \(j\), old matrices and twists
at the given field.
\end{proposition}
\begin{proof}
Fourier inversion for the indicator at \(I_1\) gives the first
formula. In the added test the Fourier sum in \(y\) gives the
point coefficient at \(I_2\), while orthogonality in \(z\)
contributes one for each equality \(\chi_g=\beta\). Thus all
terms with the same frequency have the same phase, including
when the stabilizer of the trace character is nontrivial.
Their multiplicities sum to \(|G|\). In a complex embedding
each root of unity still has modulus one and all multiplicities
remain nonnegative integers. This proves the norm statements,
since the unmarked vacuum has a single coefficient equal to one.
The tensor denominators have exponent \(pj+pt_j=p(9m+1)\).
The chosen half of the primes above each old mark has total
residue degree \(f\), the old marks are unramified, and hence
\(|R|=M^{2f}\). The remaining factor in \(|S|\) is
\(N_{E/\QQ}\mathfrak q\), which is a unit at the stated primes.
\end{proof}

\begin{proposition}[An integral retraction on the finite Fourier rows]
\leavevmode\label{prop:markedretract820}
For the same reduction of the inclusion points at the old marks,
the passage to the twisted residue tests on the tensor product of
the three finite coefficient spaces is a split injection over the
cyclotomic integers with \(N_{E/\QQ}\mathfrak q\) inverted. At every
complex embedding this injection preserves the sum of absolute
values of the coefficients. It applies to arbitrary linear
combinations of old rows with these fixed new inclusion points.
\end{proposition}
\begin{proof}
Separate the old and new matrix variables by the Chinese remainder
isomorphism. The trace characters factor through the old variables.
On the new constrained columns the residue function is a point
indicator at the pair \(I=(I_{1,\rm new},I_{2,\rm new})\) in the
additive group \(A'=R_{\mathfrak q}^{p\times(9m+1)}\).
On the other new columns it is constant. Thus on Fourier rows the
map tensors with the row \(v_\lambda=\overline{\lambda(I)}/|A'|\)
on \(\widehat {A'}\) and with the unit vector at the trivial
frequency on the other new columns. These factors have norm one
for the sum of absolute values. To retract, keep that trivial
frequency and sum the entries over \(\lambda\) with multipliers
\(\lambda(I)\). This sends the tensor factors back to one.
All coefficients of the two maps lie in the indicated ring.
If the old unit of norm minus one is changed, the old point
row is multiplied coordinatewise by the character values of
the displacement. This diagonal operation and its inverse are
integral and have norm one at every embedding.
\end{proof}

We next evaluate the three positive definite kernels in the
marked CM normalization of Theorem~\ref{thm:vac242}. In the
notation \((r_i,s_i,d_i)\) for the source rank, target rank
and polynomial degree they have parameters
\[
        (m,3m,0),\qquad(j,p,0),\qquad(12m+1-j,p,d).
\]
The real test on the third kernel is the Weyl-vector-valued
polynomial \(P_2=H_c\zeta(Z;Y)\) of
Theorem~\ref{thm:stackedheight819}. The other polynomials
\(P_0,P_1\) are one. The vector-valued evaluation means to
apply the same complex linear theta functional to each
coefficient. Use real identity arguments and allow arbitrary
finite pair arguments, including the block diagonal arguments
occurring in a definite class sum. Denote the resulting kernels
by \(K_i(P_i)\). The three residue tests above specify their
finite vectors in the unit lattice vacuum. Conjugating a positive
kernel value is allowed by conjugating its normalization as well.

\begin{theorem}[Conjugate bounds for the positive definite class kernels]
\leavevmode\label{thm:markeddefinite820}
At each field and \(j\) there are nonzero scalars \(a_i\),
depending on the marked oscillator identifications, the ranks
and the polynomial degrees, such that \(a_i K_i(P_i)\) are
algebraic in the indicated scalar or Weyl coordinates. At
every complex embedding \(\rho\) of their joint coefficient
field,
\[
 \prod_{i=0}^2\max\bigl(1,\|\rho(a_i K_i(P_i))\|\bigr)
      \le \exp\!\bigl(C_{q_0,d_{\rm K}}\,f m b_x\log b_x\bigr).
\]
The same constants and scalars apply to all classes, old matrices
and allowed twists on the original intervals. The quadratic CM
field and marked elliptic differential are fixed as the
field varies in the interval.
\end{theorem}
\begin{proof}
Write \(D_i=fr_is_i\), \(c=\pi/A\) in the marked differential.
The products \(r_is_i\) are even since \(m\) is even and
\(p=s+9m+1\) is even. Theorem~\ref{thm:vac242} identifies each
kernel, multiplied by \(a_i=C_i^{-1}c^{d_i/2}\), with
\(P_i(\partial)G_{\phi_i}(0)\). Here \(G_{\phi_i}\) is the
marked section of its finite Weil vector. The scalar \(C_i\)
is fixed by the unit product vacuum in that rank. Changes of
the finite pair arguments and of the characters in the pair
action change the finite vector in this formula.

In the split position model a character of an integral residue
box multiplies its unit vacuum by that character on the box.
This multiplication is a frequency translation. Indeed the
dual of the quotient of the position lattice by the residue
sublattice is its annihilator quotient in the frequency space.
Any two lifts of the same frequency act identically on the box.
At all places off this box the vector is the vacuum. Taking
products gives the finite expansion in rational torsion
translates of the standard vacuum, with the Fourier coefficients
of Proposition~\ref{prop:markedfourier820}. Rational representatives
of the frequency cosets exist by additive approximation. Correcting
the Weyl convention or changing such representatives can insert
roots of unity. The sums of absolute values are therefore bounded
at every embedding by \(L_0=L_1=1\) and \(L_2=|G|\).

First use finite pair representatives supported at completely
split primes outside the level, residue modulus and discriminant,
as in Theorem~\ref{thm:class61}. The local calculation in
Lemma~\ref{lem:split61} sends the unit vacuum to the normalized
vacuum on the moved Hom lattice. It transports each of the
frequency vectors into a torsion vector relative to that lattice
by conjugation of the Heisenberg action. The determinant character
factors of these pair actions are algebraic and have modulus one
in every embedding, by the ray class argument of
Theorem~\ref{thm:one61}. Theorem~\ref{thm:transport59} now
gives algebraic Taylor coefficients and, for every embedding,
\[
 |G_{\phi_i}^{\rho}(v)|\,e^{-\pi\|v\|_2^2/(2A_\rho)}
   \le L_i\bigl[e\max(1,2r_i s_i\operatorname{rd}(E)/\pi)
                      \,\nu_\rho\bigr]^{D_i/2}.
\]
In fact a torsion translate has the same metric supremum as
its vacuum. Apply the triangle inequality to the translated
sections after conjugation. This argument places no upper bound
on the orders of their torsion points.

For arbitrary finite arguments choose open compact subgroups
fixing the finite test under the two pair actions. Such
subgroups exist by smoothness. At these levels express each
argument as a rational point times a representative just used
times a right level element. The right elements fix the test.
Using rational automorphy transfers the rational points to
the real polynomial. Their effect on the matrix variable
at each embedding is a left and a right unitary substitution
and on the vacuum line a circle character. The substitutions
have algebraic entries in the CM field and their conjugates
are still unitary, since field conjugation commutes with its
complex involution. The circle character on a rational
determinant is algebraic and all its conjugates have modulus
one by the same determinant argument. For the third kernel
a matrix argument in the unit polydisc is thereby sent to one
with entries of modulus at most \(\sqrt{D_i}\), by preservation
of the Hilbert--Schmidt norm. This is at most \(M^4\) on the
interval by the bounds in Lemma~\ref{lem:sparsedistance814}.
Thus for every \(\rho\) the conjugate of the transformed
polynomial has supremum in Weyl norm on the unit polydisc
at most \(B_2=\exp(O_{q_0,d_{\rm K}}(f m b_x\log b_x))\),
by Theorem~\ref{thm:stackedheight819}. For the constant inputs
take \(B_0=B_1=1\). These arguments also prove the asserted
algebraicity for arbitrary finite arguments.

For completeness let \(P(v)=\sum_{|\alpha|=d_i}v_\alpha v^\alpha\)
be one of the resulting conjugate polynomials. The Cauchy
formula on the unit polytorus bounds \(\|v_\alpha\|\) by \(B_i\),
including for the vector polynomial by integration in the
finite-dimensional Hilbert space. On the same torus the last
display bounds the section by its right side times
\(\exp(\pi D_i/(2A_\rho))\). The Cauchy formula for the
section therefore gives
\[
 \|P(\partial)G_{\phi_i}^{\rho}(0)\|
 \le B_i (D_i+1)^{d_i}d_i!\,
  L_i\,[e\max(1,2r_i s_i\operatorname{rd}(E)/\pi)\nu_\rho]^{D_i/2}
       e^{\pi D_i/(2A_\rho)} .
\]
Indeed there are at most \((D_i+1)^{d_i}\) multi-indices of
degree \(d_i\), and \(\prod_a\alpha_a!\le d_i!\).
The case of degree zero gives the value of the section at zero.

Only finitely many marked differentials occur among the conjugates
of the fixed elliptic data, so \(A_\rho^{-1}\) and \(\nu_\rho\)
are bounded with \(K\) and the marking fixed.
Lemma~\ref{lem:sparsedistance814} gives
\(\operatorname{rd}(E)\le q_1d_{\rm K}\), and the interval has
\(f\ll x^{12}\), \(s\ll x^5\), \(m\ll \sqrt{x}\) and
\(b_x\ge x^{400}\). Thus \(D_i\log(2r_is_i\operatorname{rd}(E))\)
and \(D_i\) are absorbed in the claimed logarithm.
The degree satisfies \(d\ll f m b_x\), and \(\log(d+2)\) and
\(\log(D_i+1)\) are \(O(\log b_x)\). Finally
\(L_2=|G|\le M^{18fm^2}\) and \(\log M\le22\log b_x\).
Multiplying the three bounds proves the theorem.
\end{proof}

Proposition~\ref{prop:markedretract820} concerns the finite
coefficient rows before applying the theta functional.
Theorem~\ref{thm:markeddefinite820} controls the archimedean
sizes after evaluation. An integral comparison of the projected
periods still requires finite-place bounds for those evaluated
derivatives and the identity relating the definite class sums to
the projected carrier integrals.
\clearpage

\subsection{An all-prime normalization of the definite class kernels}
\label{sec:definiteintegral821}
The comparison with the marked elliptic curve in
Theorem~\ref{thm:markeddefinite820} uses a rational orthonormal
basis. A suitable choice makes the lattice index in that
comparison small enough at the finite places. We make this
choice once per field and pair of ranks, before specifying any
twist or testing any prime.

Write \(K=\QQ(\sqrt{-d_{\rm K}})\) and
\(E=K(\zeta_{q_1})\). The prime divisors of \(q_1\) form a
fixed set on the original interval, since \(q_1=3q_0l_0 2^e\)
with \(l_0\) a power of three. Put \(u=\operatorname{rad}(q_1)\),
\(E_0=K(\zeta_u)\), and choose the primitive roots compatibly.
Set \(h=q_1/u=b w^2\) with \(b\) squarefree and \(w>0\). There
are finitely many possibilities for \(b\), all divisors of \(u\).

\begin{lemma}[A bounded change of lattice to product coordinates]
\leavevmode\label{lem:tracebasis821}
Let \(r r'\) be even, \(D=f r r'\), and give \(E^{rr'}\) the
\(K\)-Hermitian trace form
\(H(X,Y)=\operatorname{Tr}_{E/K}(\sum X_i\bar Y_i)\).
There is a positive integer \(C_0\) depending only on
\(q_0,d_{\rm K}\), common to all ranks, with the following
property. One can choose a \(K\)-orthonormal basis spanning
a product period module \(B\) over \(\OO_K\), such that for
\(n=C_0 w q_1\) and the standard oscillator period module
\(L\) one has
\[
       nB\subset L,\qquad nL\subset B,qquad
       n V,\ nV^*\text{ have algebraic integer entries}.
\]
Here \(V\) is the matrix of the basis vectors in the oriented CM
embeddings. In particular \(I=[L:L\cap B]\) divides \(n^{2D}\)
and \(\log n=O_{q_0,d_{\rm K}}(\log(2x))\) at each node \(x\).
\end{lemma}
\begin{proof}
The cyclotomic and quadratic discriminants are coprime, and
\([E:E_0]=\varphi(q_1)/\varphi(u)=h\). Thus
\(T^h-\zeta_u\) is the minimal polynomial of \(\zeta_{q_1}\)
over \(E_0\). Moreover \(\OO_E=\OO_{E_0}[\zeta_{q_1}]\):
the integer rings of the quadratic and the cyclotomic fields
multiply to the integer ring of their compositum by coprimality
of the discriminants, both at \(q_1\) and at \(u\).
Consequently
\[
                \OO_E=\bigoplus_{a=0}^{h-1}
                           \zeta_{q_1}^a\OO_{E_0}.
\]
The roots of the minimal polynomial are all the \(h\)-th roots
of \(\zeta_u\). The trace to \(E_0\) of
\(\zeta_{q_1}^a\overline{\zeta_{q_1}^{a'}}\) is \(h\) when
\(a=a'\) and zero otherwise, by summing the powers of the
\(h\)-th roots of unity. The summands in the display are
therefore mutually orthogonal for the trace to \(K\).

Group the \(rr'\) entries into pairs. On the fixed space
\(E_0^2\) choose for each possible \(b\) a \(K\)-orthonormal
basis for \(b\operatorname{Tr}_{E_0/K}(x_1\bar y_1+x_2\bar y_2)\).
For existence, diagonalize the form on one copy over \(K\).
Each positive rational diagonal entry \(a\), including the
factor \(b\), occurs twice. As in Lemma~\ref{lem:basis59},
Meyer's theorem on the indefinite rational form
\(X^2+d_{\rm K}Y^2+Z^2+d_{\rm K}T^2-S^2/a\) gives
\(z_1,z_2\in K\) with \(|z_1|^2+|z_2|^2=1/a\). The columns
\((z_1,z_2)\), \((-\bar z_2,\bar z_1)\) orthonormalize
the pair. The finitely many resulting \(\OO_K\)-spans \(B_b\)
are full lattices. Choose one integer \(C_0\) such that
\(C_0 B_b\subset\OO_{E_0}^2\) and
\(C_0\OO_{E_0}^2\subset B_b\) for every \(b\).
Commensurability of full integer lattices proves existence of
such an integer.

In each pair and each cyclotomic summand multiply the
chosen basis by \(\zeta_{q_1}^a/w\). The trace identity and
\(h/w^2=b\) prove orthonormality, and the direct sum of all
these bases spans \(E^{rr'}\). Their integral span satisfies
\[
 C_0 w B\subset\OO_E^{rr'},\qquad
 C_0\OO_E^{rr'}\subset w B\subset B .
\]
The oscillator period module is obtained from \(\OO_E^{rr'}\)
by the inverse half of the relative different at the specified
split components. Locally its scalar is one or the inverse
of a local generator of that different. Thus
\(\OO_E^{rr'}\subset L\) and \(q_1 L\subset\OO_E^{rr'}\)
by the different calculation in Lemma~\ref{lem:sparsedistance814}.
This proves both inclusions in the statement. The first
inclusion of the last display also makes the entries of
\(C_0wV\), and hence their complex conjugates, algebraic
integers. Finally \(L/(L\cap B)\) embeds in
\((n^{-1}B)/B\), a group of order \(n^{2D}\). The size estimate
follows from \(q_1\ll_{q_0} x^{11}\log(2x)\) and \(w^2\le q_1\).
\end{proof}

For the three positive kernels in the preceding subsection set
\[
   e_0=0,\qquad e_1=pj,\qquad e_2=p(9m+1-j),\qquad
   T=3d_{\rm K} M^2 n^2 .
\]
Use the basis of the lemma in the definition of the marked
scalars \(a_i\) of Theorem~\ref{thm:markeddefinite820}. Define
\[
           A_i=a_i\, M^{2f e_i}\,n^{D_i}\,T^{d_i}.
\]
These scalars and \(H_c\) are chosen before the finite
arguments, the subgroup fixing the residue function, the
embedding of the coefficient field and the prime of reduction.

\begin{theorem}[Integral normalized class kernels at the prescribed levels]
\leavevmode\label{thm:scaleddefinite821}
Let \(\ell\) be any rational prime with
\(\ell\nmid N_{E/\QQ}\mathfrak q\), and fix a place
\(\lambda\mid\ell\) of the algebraic numbers. If the two
finite arguments of the \(i\)-th positive definite kernel are
integral at all places above \(\ell\), then for every homogeneous
polynomial \(P\) of its specified degree with \(\lambda\)-integral
coefficients,
\[
                        A_i K_i(P)\in\OO_{\bar\QQ,\lambda}.
\]
The conclusion applies coordinatewise to an integral lattice
of polynomial-valued vectors. In particular it applies to
\(H_c\zeta\) whenever its row parameters are integral in
the input frames at \(\lambda\). For the actual polynomials
and short rows in Theorem~\ref{thm:markeddefinite820}, the
bound there on the product of conjugate sizes also holds
with these \(A_i\), with a constant depending only on
\(q_0,d_{\rm K}\).
\end{theorem}
\begin{proof}
Fix \(i,\lambda\). First arrange the finite pair arguments
at split primes away from \(\ell\), the discriminant, the level
and all the residue primes. This is possible in their double
classes by Theorem~\ref{thm:class61}. Choose each right level
inside the integral group at \(\ell\), small enough to fix the
test for the given finite pair actions. If the original arguments
were integral above \(\ell\), the rational left factors
relating them to these representatives are integral there
with integral inverses. Automorphy transfers those factors to the
real polynomial. The resulting substitutions have integral
entries at \(\lambda\) and preserve integrality of its
coefficients. Indeed the matrix entries and their inverses
are integral at every place of \(E\) above \(\ell\), which also
covers the permuted embeddings in the CM coordinates.
The determinant characters on the rational factors are units
there as well: their integral powers are products of conjugates
of unit determinants and their inverses. The corresponding
character factors of the finite split representatives are
\(\lambda\)-units. In the ray class argument of
Theorem~\ref{thm:one61} a positive power of a lift supported
away from \(\ell\) is principal times finite units. Its
principal generator is a unit at all places above \(\ell\),
giving the assertion by the same character formula.

We now express the section as a sum of torsion translates on
the product of marked elliptic curves. For the first passage,
from \(B\) to \(L\), Lemma~\ref{lem:basis59} sums the translates
by \(L/(B\cap L)\) with signs and factor \(I^{-1/2}\).
Lemma~\ref{lem:tracebasis821} implies
\(n^{D_i}/\sqrt I\) is an algebraic integer, since its square
is the integer \(n^{2D_i}/I\). For the next passage, from
\(L\) to the moved lattice \(L'\), the normalized sum of
Lemma~\ref{lem:split61} has denominator
\([L':L\cap L']^{1/2}\). The local modules agree above \(\ell\),
so this index is prime to \(\ell\). Its normalization is a
\(\lambda\)-unit. The remaining coefficients from the finite
test are those of Proposition~\ref{prop:markedfourier820}.
For the individual leg they become \(\lambda\)-integral after
multiplication by \(M^{2f e_i}\), since the rest of their
denominator is a power of \(N_{E/\QQ}\mathfrak q\).
The signs, torsion composition phases and changes of rational
frequency representatives contribute roots of unity.
These facts give \(\lambda\)-integral coefficients for the
whole sum after the first two integer factors in \(A_i/a_i\).

Consider the order at \(\ell\) of a translation point after
these operations, in product coordinates modulo \(B\). Write
\(L_\ell=L\otimes\ZZ_\ell\), and likewise for the other modules.
The first set of coset representatives is in \(L_\ell\) and
the second is in \(L'_\ell=L_\ell\). For the further finite
frequency translate the class modulo \(L'_\ell\) is killed by
\(M^2\). Indeed at a place over \(\ell\) it is dual to a
quotient of the integral position box by the old residue ideal
or is trivial. This quotient is killed by the rational integer
\(M^2\), and the new residue ideal is a unit at \(\ell\).
These assertions persist on taking a rational representative
of the adelic coset. Composing the three translations adds
their point coordinates. It follows that multiplying the local
coordinate of the sum by \(nM^2\) puts it in \(B_\ell\).
Equivalently its finite order modulo \(B\) has \(\ell\)-part
dividing the \(\ell\)-part of \(nM^2\). This equivalence
follows by tensoring the finite cyclic group generated by
the point with \(\ZZ_\ell\). It also holds as an upper bound
for each of its elliptic product coordinates.

Use \(v_\lambda(\ell)=1\). For an elliptic translate of
order \(N\), Theorem~\ref{thm:torsion312} bounds its derivative
of order \(k\) below in valuation by
\[
          -k\delta-\lfloor k/2\rfloor
                         (\kappa_\ell-2\delta)_+,
          \qquad 0\le\delta\le v_\ell(N).
\]
Here \(\kappa_\ell\le1\), and it is zero away from
\(3d_{\rm K}\) by Theorem~\ref{thm:elliptic312}.
For the coordinates just found this loss is at most
\(k v_\ell(3d_{\rm K} n M^2)\). This includes \(k=0\).
A multi-derivative on an elliptic product distributes across
the single-coordinate factors. Therefore multiplying by
\((3d_{\rm K} n M^2)^{d_i}\) clears every derivative of
total order \(d_i\) in product coordinates. The change from
the embedding coordinate \(Z\) to product coordinates is
\(z=V^*Z\). Iterating the chain rule introduces products of
\(d_i\) entries of \(V^*\), all cleared by the further
factor \(n^{d_i}\). Hence \(T^{d_i}\) clears the undivided
derivatives in the matrix variables. Contraction with the
integral coefficients of \(P\) and addition preserve
integrality. The marked identity of Theorem~\ref{thm:vac242}
now proves the claimed bound for the kernel. For a vector,
apply this argument to its coordinates in an integral basis.

To retain the archimedean estimate, the extra logarithm for
the three kernels together is
\[
   2fp(9m+1)\log M+\Big(\sum_i D_i\Big)\log n+d\log T.
\]
The bounds of Lemma~\ref{lem:tracebasis821} and
Theorems~\ref{thm:stackedheight819} and~\ref{thm:markeddefinite820}
make this \(O_{q_0,d_{\rm K}}(f m b_x\log b_x)\).
All the added factors are positive rational integers.
Their sizes in every embedding are the same.
\end{proof}

The normalization now controls the three evaluated positive
kernels at integral arguments as well as their complex
conjugate sizes. For the projected period one must also track
the output Weyl lattice at a general finite representative,
the class integration weights and the equality with the
carrier-defined factors.
\clearpage

\subsection{The definite class sum in its Weyl lattice}
\label{sec:definiteclass822}
We pass from integral arguments in Theorem~\ref{thm:scaleddefinite821}
to arbitrary finite representatives, keeping track of the lattice
in the target. Let \(C(c)\), \(D(c)=C(c)^*\) and \(T_c\) be the
flag-minor module, its dual and the map of
Lemma~\ref{lem:complementfischer819}. All three include the
tensor product over the oriented embeddings. For a matrix \(A\)
in one block set
\[
       R_A f(X)=f(A^{\mathsf t} X),\qquad
       \rho_D(A)l(f)=l(R_{A^{\mathsf t}} f).
\]
These formulas define representations since transposition reverses
products. Substitution and Cauchy--Binet show that they preserve
their respective integral lattices for any invertible integral
matrix with integral inverse. The identity
\(K_c(Y,A^{\mathsf t}X)=K_c(AY,X)\) proves
\(R_A T_c=T_c\rho_D(A)\). We use \(R\) and \(\rho_D\)
componentwise for tuples of matrices.

For a finite target argument \(h\in U_p(\mathbb A_{F,f})\)
and a place \(\lambda\mid\ell\) of the algebraic numbers, put
\[
   {\cal D}_{h,\lambda}
      =\rho_D\big((h_{\lambda\circ\bar\tau})_\tau\big)
                  \bigl(D(c)\otimes\mathcal O_{\CC_\ell}\bigr).
\]
Here \(\tau\) is one of the chosen embeddings of \(E\) into
the complex numbers, \(\bar\tau(x)=\overline{\tau(x)}\),
and \(h_{\lambda\circ\bar\tau}\) is the \(E\)-matrix component
at the corresponding finite place transported to \(\CC_\ell\).
We choose a coefficient embedding inducing \(\lambda\) to
read this formula. The conjugate embedding occurs because the
output vector in the class sum is complex conjugated.

Let \(b_0\) be a representative of a definite class of rank
\(6m+1\), and take the short matrix \(Y\) of
Theorem~\ref{thm:stackedheight819} in that lattice. In
particular the columns of \(Y^{\mathsf t}\) belong to
\(b_0\widehat{\OO}_E^{6m+1}\). Let \(v,w\) be finite
arguments of ranks \(m,3m\), and put
\(g=\operatorname{diag}(1_{3m-j},b_0,w)\). Use the positive
kernels, the residue tests and the scalars \(A_i\) from the
last two subsections. Set
\[
 k_0(v,w)=A_0 K_0(1;v,w),\qquad
 k_1(h)=A_1 K_1(1;1,h),\qquad
 V(g,h;Y)=A_2 K_2(H_c\zeta(\,\cdot\,;Y);g,h).
\]
The character on the first source is the tame twist times
the prescribed power of the weight-zero character, and the
first target character is trivial. For the second pair the
source character is trivial. On both members of the third
pair the character has exponent
\[
        e_\tau=j+(t+k-1)/2-w_{j,\tau}=-\mu_\tau .
\]
Here \(\mu_\tau\) is the number of determinant factors in its
stacked polynomial, as computed in the proof of
Lemma~\ref{lem:carriercoefficient815}.

\begin{lemma}[Rational changes of the short rows and target frame]
\leavevmode\label{lem:frameequiv822}
Choose rational unitary matrices \(\gamma_b,\gamma_w,\gamma_h\)
of ranks \(6m+1,3m,p\). Put
\(\Gamma=\operatorname{diag}(1_{3m-j},\gamma_b,\gamma_w)\) and
\(Y'=Y\gamma_b^{\mathsf t}\). Then
\[
  V(\Gamma g,\gamma_h h;Y')
      =\rho_D\big((\tau(\gamma_h))_\tau\big)V(g,h;Y).
\]
The scalar \(k_0\) is invariant under rational left changes of
both its finite arguments. The scalar \(k_1\) is invariant
under a rational left change of its finite target argument.
\end{lemma}
\begin{proof}
In a definite pair with source and target splitting exponents
\(e_s,e_t\), a rational left change \((\gamma,\gamma')\)
of the finite arguments can be removed by the global rational
action. The remaining real element is the inverse pair. Write
\(A=\tau(\gamma'^{-1})\), \(B=\tau(\gamma)\). Its action on
creation vectors in the matrix space is \(L\mapsto A L B\).
On polynomial generators, taking the coefficient of the
matrix unit \(E_{ij}\) consequently substitutes
\(Z\mapsto A^{\mathsf t}Z B^{\mathsf t}\): the image of the
generator is the polynomial with coefficients \((AE_{ij}B)_{kl}\).
Multiplication gives the substitution on every monomial.
The compact splitting on the vacuum line has target determinant
exponent \(e_s\) and inverse source determinant exponent \(e_t\)
for the real pair. For the inverse pair the resulting factor is
\((\det A)^{e_s}(\det B)^{e_t}\).
One can see the character on diagonal elements from the
vacuum weight on each circle. Both it and the determinant
formula are characters of the compact groups, so equality
there gives equality everywhere by unitary diagonalization.
Multiplying the factors over embeddings yields the change
for the theta kernel. For \(k_0\) both exponents are zero.
For \(k_1\) only the target varies, its source exponent is
zero and the test polynomial is constant. This proves the
two scalar assertions.

For the third pair take \(B=\tau(\Gamma)\) and adjoin its
identity block on the first \(a=k+j-1\) columns, writing
\(B'=\operatorname{diag}(1_a,B)\).
For each determinant of the stacked polynomial \(P_Y(X,Z)\)
in Proposition~\ref{prop:stackednorm819}, the constant rows
for \(Y'\) are the old ones times \((B')^{\mathsf t}\).
Indeed the coordinate selectors are on the fixed columns,
and the rows from \(Y\) are on the \(\gamma_b\) block.
Substituting \(A^{\mathsf t}X\) and
\(A^{\mathsf t}Z B^{\mathsf t}\) in this new determinant
therefore multiplies the original full square matrix on the
left by \(\operatorname{diag}(1_{2m},A^{\mathsf t})\)
and on the right by \((B')^{\mathsf t}\).
Its determinant gains \(\det A\det B\). In the product of
the determinants this has exponent \(\mu_\tau\) at \(\tau\).
The multiplier from the vacuum is its inverse, since both
exponents are \(-\mu_\tau\).

Now apply the kernel only in the \(Z\) variables and retain
\(X\) as polynomial variables. The cancellation just proved
says that applying \(R_{(A)_\tau}\) to the new flag polynomial
gives the old flag polynomial. Taking \(T_c^{-1}\) converts
\(R\) into \(\rho_D\), since
\(P_Y=T_c\zeta(\,\cdot\,;Y)\) as a polynomial in \(X\).
The constants \(H_c A_2\) occur on both sides. Inverting
\(A=\tau(\gamma_h^{-1})\) gives the desired formula.
\end{proof}

Let \(\mathcal C_t\) be the full-level definite classes of rank
\(t\) and \(d_z\) the order of the rational stabilizer of
the corresponding lattice. Use one finite representative
for each class. Put \(B_*=|\operatorname{GL}_{6fm}(\mathbb F_3)|\).
Define the positive definite class expression, in the
normalizations above, by
\[
 \mathcal Y(h)=B_*^2\sum_{w\in\mathcal C_{3m}}
             \sum_{v\in\mathcal C_m}
       \frac{\overline{k_0(v,w)\,\overline{k_1(h)}\,V(g,h;Y)}}
            {d_w d_v}.
\]
The weights are those of full-level class integration with the
real compact group and the finite level each of volume one.
The estimates below hold for every choice of representatives.

\begin{theorem}[Integral class vectors and their conjugate sizes]
\leavevmode\label{thm:definiteclass822}
For every permitted twist and every finite target \(h\), the
algebraic class vector \(\mathcal Y(h)\) belongs to
\({\cal D}_{h,\lambda}\) at every \(\lambda\mid\ell\) with
\(\ell\nmid N_{E/\QQ}\mathfrak q\). At every complex embedding,
\[
       \|\rho(\mathcal Y(h))\|^2
            \le \exp\!\bigl(C_{q_0,d_{\rm K}} f m b_x\log b_x\bigr).
\]
These conclusions use the same choices of \(A_i,H_c,B_*\)
for all twists and the same fields and intervals as
Theorem~\ref{thm:complexseparation818}.
\end{theorem}
\begin{proof}
We first move the finite arguments into the integral group
above \(\ell\). Rational unitary matrices are dense in the
product of the local unitary groups at these finitely many
places of \(F\). To see this directly, use the Cayley chart
\(X\mapsto (1-X)(1+X)^{-1}\) on the skew-Hermitian matrices.
Its inverse has the same formula in the group variable
where \(1+g\) is invertible. Every local open set meets
that chart. One can multiply a point by central norm-one
scalars arbitrarily close to one, avoiding the finitely many
scalars that give eigenvalue \(-1\). There are infinitely many
such local scalars by the same one-dimensional Cayley formula.
Ordinary weak approximation in the rational vector space of
skew-Hermitian matrices now gives simultaneous approximation
inside all the prescribed open sets. This works also above two.

Choose rational left multipliers moving \(v,w,b_0,h\)
into the integral groups at these places. Use the same
multiplier \(\gamma_w\) for both occurrences of \(w\).
For the added pair use the resulting \(\Gamma\) of
Lemma~\ref{lem:frameequiv822}. The transformed short rows
are integral above \(\ell\) since
\[
       (Y')^{\mathsf t}
          =(\gamma_b b_0)(b_0^{-1}Y^{\mathsf t})
\]
is there a product of integral matrices. By
Theorem~\ref{thm:stackedheight819}, \(H_c\zeta(\,\cdot\,;Y')\)
has integral coefficients at each coefficient place above
\(\ell\). Theorem~\ref{thm:scaleddefinite821} applies to all
three kernels in their moved frames. In particular \(k_0\) and
\(k_1\) themselves are integral there by their invariances.

To treat the conjugated vector at \(\lambda\), use that theorem
also at the place composed with the original complex conjugation.
The vector \(\overline{V(\Gamma g,\gamma_h h;Y')}\) is consequently
integral in the standard Weyl basis at \(\lambda\).
Lemma~\ref{lem:frameequiv822} gives this vector as the action
of \(\rho_D((\bar\tau(\gamma_h))_\tau)\) on the original
\(\overline{V(g,h;Y)}\). Since \(\gamma_h h\) is integral
above \(\ell\) with integral inverse, it preserves the
standard lattice under each of the indicated substitutions.
Applying the inverse action of \(\gamma_h\) therefore puts
the original vector exactly in \({\cal D}_{h,\lambda}\).
Multiplication by the two integral scalars preserves this
membership.

Every stabilizer used in the sum acts faithfully on the
underlying integer lattice of rank \(2fm\) or \(6fm\).
It is finite since the trace form is positive definite.
Extending by the identity if necessary gives a finite subgroup
of \(\operatorname{GL}_{6fm}(\ZZ)\). Reduction modulo three
is injective on each such subgroup. Indeed a nonidentity element
of prime order in the kernel would have the form \(1+3^a U\)
with \(a\ge1\) and \(U\) nonzero modulo three. Its power at
a prime \(l\ne3\) differs from one modulo \(3^{a+1}\).
Its cube differs from one modulo \(3^{a+2}\), by expansion.
A nontrivial finite kernel would contain an element of prime
order, a contradiction. Hence \(d_v,d_w\) both divide \(B_*\),
and \(B_*^2/(d_vd_w)\) is integral. The local assertion
follows on summing.

For the archimedean bound use each of the individual bounds
of Theorem~\ref{thm:scaleddefinite821}, at \(\rho\) or at its
composition with complex conjugation as required. The Weyl
metric is defined by the real Gram kernel, so the norm is the
one used in that theorem. The logarithm of \(B_*\) is at most
\((6fm)^2\log3\). The full-level genus count in the proof
of Lemma~\ref{lem:mass32} bounds the logarithm of the number
of classes in either rank by
\(O((6fm)^4\log(2\cdot 6fm q_1d_{\rm K}))\). Its hypotheses
hold since \(E\subset\QQ(\zeta_{q_1d_{\rm K}})\).
With \(f\ll x^{12}\), \(m\ll\sqrt x\) and \(b_x\ge x^{400}\),
these two logarithmic costs are absorbed in \(f m b_x\log b_x\).
The triangle inequality proves the result, also for a vanishing
sum. All coordinates are algebraic by construction.
\end{proof}

This establishes integrality in the moving output lattice for
the positive definite double class expression. To apply the
bound to the projected periods of
Theorem~\ref{thm:complexseparation818}, the change of sign in
the source space and the successive global theta integrations
must identify that expression with the prescribed first factor,
including its scalar across the twists.
\clearpage

\subsection{The change of sign in the definite class series}
\label{sec:signeddefinite823}
The middle kernel in the prescribed definite class series has
source \(-I_j\) and target \(I_p\). We relate it to the positive
kernel of Theorem~\ref{thm:definiteclass822}. The scalar in this
relation must be chosen before the level, the finite target
argument and the twisting characters.

Fix the integral finite unit \(u_f\) with \(u_f\bar u_f=-1\)
used to frame the negative source, and write \(u=(1,u_f)\) at
the real and finite places. Such a unit exists on these fields.
Indeed \(E/F\) is unramified at every finite place: away from
\(d_{\rm K}\), inertia acts trivially on \(K\), and at
\(d_{\rm K}\) it acts trivially on \(\QQ(\zeta_{q_1})\).
Complex conjugation acts nontrivially on each of these two
fields, so it belongs to none of these inertia groups.
At a split place take a pair of units with product \(-1\).
At an unramified quadratic place the norm on residue units is
surjective, and on successive quotients of the principal units
it induces the surjective residue trace. Lifting successively
and taking the limit gives a unit of norm \(-1\), also above two.
We use the reductions of the one chosen \(u_f\) at every modulus.

Use the pair splitting whose two external characters \(a,b\)
multiply the action on the Hom space by
\(a(\det h)b(\det g)^{-1}\) at arguments \((g,h)\).
Write \(\upsilon_1\) for the restriction of the fixed
unitary splitting character to norm-one ideles. Denote by
\(K^+_{a,b}(\Psi;h)\) and \(K^-_{a,b}(\Psi;h)\) the kernels
of \((I_j,I_p)\) and \((-I_j,I_p)\) at polynomial one,
real identity, source identity and finite target \(h\).
Here \(\Psi\) is the position function on the marked integral
residue box, and the unit vacuum is taken at the other finite
places.

\begin{lemma}[One phase for reversing the source form]
\leavevmode\label{lem:signkernel823}
There is a number \(\epsilon_j\) of absolute value one such
that for every such datum and every \(h,a,b\),
\[
 \overline{K^+_{a,b}(\Psi;h)}
       =\epsilon_j\,K^-_{a^{-1}\upsilon_1^j,
                                 b^{-1}\upsilon_1^p}
          \bigl(x\mapsto\overline{\Psi(u_f^{-1}x)};h\bigr).
\]
The unit in this argument is reduced in the indicated residue
ring. The character of the source determinant in the displayed
negative kernel has no effect at the identity source argument.
\end{lemma}
\begin{proof}
First consider just the Heisenberg data on the two Hom spaces.
In the negative coordinates let \(D\) be the identity on the
real coordinates and multiplication by \(u_f^{-1}\) on the
finite coordinates, mapping to the positive coordinates.
At infinity the signs of the Hermitian form have been reversed.
In the standard finite frames the identity
\(u_f^{-1}\overline{u_f^{-1}}=-1\) changes the finite Hermitian
pairing by a minus sign. Consequently the Weyl multiplication
factors in the two coordinates are complex conjugates. The
Gaussian has the same magnitude in the two real models, and
the standard finite box is preserved since \(u_f\) is a unit.
There is an antiunitary map \(J\) taking the positive unit
vacuum to the negative unit vacuum and taking a Weyl translate
by \(v\) to the translate by \(D^{-1}v\). To construct it on
linear combinations of vacuum translates, conjugate the
coefficients and send each translate to the one specified above.
Their inner products are conjugated: the Weyl multiplication
factor has just been checked, the real vacuum coefficient is a
Gaussian, and the finite vacuum coefficient is the lattice
indicator in the corrected Weyl convention. Any quadratic
correction in the symmetric convention is conjugated by the
same pairing calculation. Thus the map is isometric on the
span and extends, since these translates span densely by
irreducibility. At a split residue box the position action is
\(f(x)\mapsto\overline{f(u_f^{-1}x)}\), as is also seen on the
Fourier expansion of Proposition~\ref{prop:markedfourier820}.
The real polynomial one maps to one.

The map \(D\) identifies the rational translation subgroups in
the two framed spaces. In fact the coordinates of a rational
map \(L\) with source frame \(x_S\) and target frame \(x_T\)
are \(x_T^{-1}L x_S\). In the negative space the source frame
is \(u x_S\), so the same rational map has coordinates
\(u x_T^{-1}Lx_S\), which are taken back to the first coordinates
by \(D\). It follows that the complex-linear functional
\(f\mapsto\overline{\vartheta_-(Jf)}\) is invariant under
the same rational Weyl translations as \(\vartheta_+\).

For clarity, such invariant distributions span one line. In a
rational Schr\"odinger polarization, invariance under rational
frequencies forces support in the rational position points.
Indeed the annihilator of the rational frequencies in the
adelic positions is exactly the rational positions, by adelic
additive duality. Wherever a rational modulation differs from
one, multiplication by its difference from one is locally
invertible. At a single rational support point a distribution
in the finite variables depends only on the value of the
locally constant test in a neighborhood of that point. In the
real variables a distribution supported at a point is a finite
sum of derivatives of the point mass. Translate the point to
zero. The Fourier transform of this real point distribution
is a polynomial. Rational modulations in a rational basis of
the real dual give it a full set of independent periods: for
these finitely many modulations take the finite neighborhood
small enough that all their finite characters are one.
A polynomial with these periods is constant, by restricting
successively to the corresponding lines. This rules out the
derivatives of positive order. Invariance under rational
position translations then makes all masses equal. The equality
with a multiple of the rational summation distribution follows
first on compact tests and then by continuity on Schwartz tests.
For the two normalized functionals the proportionality factor
has modulus one: their mean squares on the quotient of the
Heisenberg translation space by its rational lattice both
equal the squared test norm by Lemma~\ref{lem:zak242}.
The isomorphism \(D\) preserves quotient Haar probability and
\(J\) preserves the norm. We obtain
\(\overline{\vartheta_+(f)}=\epsilon_j\vartheta_-(Jf)\)
with a single \(\epsilon_j\) for all data.

It remains to compare the finite pair actions under \(J\).
The linear action \(G\) of the finite pair on its Hom space
is the same in the two frames and commutes with scalar
multiplication by \(u_f\). The two Weil operators, after
conjugating one by \(J\), intertwine the same Heisenberg
automorphism. Their ratio is a scalar by irreducibility.
We compute it on the vacuum coefficient, retaining its phase.
Write \(L_f\) for the standard finite lattice in normalized
coordinates, and choose a finite matrix \(B\) with
\(B L_f=L_f\cap G^{-1}L_f\). In the unramified lattice
splitting the determinant multiplier is
\[
       \Delta(G)=\chi_f(\det B),\qquad
       \chi_f(z)=\upsilon_f(z)|z|_f^{-1/2}.
\]
Equivalently the unitary lift of the lattice change fixes the
phase of the vacuum overlap to \(\Delta(G)/|\Delta(G)|\).
The squared magnitude of this overlap is
\([L_f:L_f\cap G^{-1}L_f]^{-1}\), by orthogonality of
the coset translates in the normalized lattice sum of
Lemma~\ref{lem:basis59}. Since
\(|\Delta(G)|^2=[L_f:L_f\cap G^{-1}L_f]\), the overlap
itself is \(1/\overline{\Delta(G)}\) in both finite frames.

Let \(z=\det B\). The fractional ideals of \((\det G)z\)
and \(\bar z\) agree. At a finite place with Smith exponents
\(a_k\) of \(G\), the exponent of \(z\) is
\(\sum_k\max(0,-a_k)\), and adding \(v(\det G)\) gives
\(\sum_k\max(0,a_k)\). The exponents at the conjugate place
are the negatives of the original ones as a multiset,
because \(G^{-1}=G^*\) in the standard finite frame.
The character \(\chi_f\) is unramified and
\(\chi_f(\bar z)=\overline{\chi_f(z)}\). For the latter
identity the positive norm factors agree, and the unitary
splitting character is trivial on \(z\bar z\), since its
restriction to the base ideles is the quadratic character of
the extension and \(z\bar z\) is a norm. Also \(|\det G|_f=1\).
Consequently
\[
        \frac{\overline{\Delta(G)}}{\Delta(G)}
                          =\upsilon_f(\det G).
\]
This is the ratio of the conjugated vacuum overlap to the
negative one. At the arguments of this lemma the real
operators are the identities, and \(G\) acts on the \(j\)
columns by \(L\mapsto hL\). Hence \(\det G=(\det h)^j\).
Conjugating the external unitary character at \(h\) replaces
\(a\) by \(a^{-1}\); the overlap ratio then supplies exactly
\(\upsilon_1(\det h)^j\). This proves the stated change of
splitting. Combining it with the functional identity and
the residue-box calculation proves the lemma.
\end{proof}

\begin{theorem}[Arithmetic bounds for the prescribed definite class series]
\leavevmode\label{thm:signedseries823}
Fix the same field, \(j\), old matrix, class with short rows
and tame twist as in Theorem~\ref{thm:definiteclass822}. Write
\(\nu\) for the third pair's character, and let \(\Upsilon(h)\)
be the class series with its original signs. Explicitly it
is the negative kernel of \((-I_j,I_p)\) at the inclusion
indicator, characters \((\upsilon_1^j,\nu)\) and finite
arguments \((1,h)\), times
\[
     \sum_{w\in\mathcal C_{3m}}\sum_{v\in\mathcal C_m}
       \frac{K_0(1;v,w)\,K_2(\zeta(\,\cdot\,;Y);g,h)}{d_v d_w}.
\]
All terms here use the unscaled kernels and polynomial. There
is one nonzero scalar \(S_j\), independent of the old matrix,
short-row choice, twist and \(h\), for which
\[
       \overline{S_j\Upsilon(h)}=\mathcal Y(h).
\]
In particular the left side belongs to \({\cal D}_{h,\lambda}\)
at every prime away from the new twist conductor, and its
squared norm in every complex embedding has the bound in
Theorem~\ref{thm:definiteclass822}.
\end{theorem}
\begin{proof}
In the positive middle kernel take \(a=1\),
\(b=\nu^{-1}\upsilon_1^p\) and
\(\Psi(x)={\bf1}_{u_f x=I_{\rm incl}}\). Then
\(b^{-1}\upsilon_1^p=\nu\), and
\(\overline{\Psi(u_f^{-1}x)}\) is the prescribed negative
inclusion indicator. The lemma makes the conjugate of the
unscaled positive kernel \(\epsilon_j\) times the negative
one at every \(h\). The first positive kernel is the summand
defining the carrier coefficient at \(w\). Expanding that
coefficient gives precisely the inner sum over \(v\)
in the statement. The third kernel is linear in its
polynomial, hence inserting \(H_c\zeta\) multiplies it by
\(H_c\). Substitute these identities into the definition
of \(\mathcal Y\), without changing either class weight.
The claimed equality follows with
\[
             S_j=B_*^2 H_c A_0\overline{A_1} A_2\epsilon_j .
\]
Each factor is nonzero and was fixed for the given field and
\(j\) before the residue data or twist were chosen.
The arithmetic conclusions follow by equality with the very
same algebraic vector of Theorem~\ref{thm:definiteclass822}.
\end{proof}

The series in this theorem is defined on the definite group
of rank \(p\). The first factors of
Theorem~\ref{thm:complexseparation818} are carrier integrals
on \(U(W_j)\). Their comparison requires the intervening
global theta integrals with their levels, polynomials and
character projections.
\clearpage

\subsection{The row stabilizer in the definite theta integral}
\label{sec:refinedclass824}
Fix one field and one interval of Theorem~\ref{thm:complexseparation818},
a lift index \(j\), a choice \(i\) of old matrix \(\Lambda\),
short rows \(Y_0\) and definite class \(b_0\), and a prescribed
twist \(\vartheta\). Write \(s=\dim W_j\), \(p=s+9m+1\) and
\(a=k+j-1\) for the positive index of \(W_j\).
The rank \(p\) space
\[
 W'=W_j\oplus(-I_j)\oplus I_{3m-j}\oplus I_{6m+1}
\]
has \(m\) negative directions. In its ordered coordinate set
let \(e\) be the \(p\)-by-\(s\) matrix including the coordinates
of \(W_j\). Put \(R=\OO_E/\mathfrak a^2\) and
\(S=\OO_E/(\mathfrak a^2\mathfrak q)\) in the chosen split
components. Here \(\mathfrak q\) is the new part of the twist
modulus. For residue matrices a bar in this subsection means
the reduction \(S\longrightarrow R\). We write
\(\operatorname{conj}\) for ordinary complex conjugation.

Let \(\chi\) be the residue character of \(\chi_M^j\vartheta\),
extended by zero on nonunits. Let \(w\) be any matrix of rows
over \(R\) for which \(\Lambda=u w\) with \(u\) a
\(3m\)-rowed matrix over \(R\). We can always take \(w=\Lambda\).
The freedom to choose \(w\) allows the same calculation with
the rows occurring in an integral projection. Define
\[
 F_\Lambda(y)=\sum_{\gamma\in\operatorname{GL}_{3m}(R)}
       \psi_R\bigl(\operatorname{tr}(\Lambda\bar y\gamma)\bigr),
 \qquad y\in\operatorname{Mat}_{p,3m}(S).
\]
Thus \(F_\Lambda(yB)=F_\Lambda(y)\) for
\(B\in\operatorname{GL}_{3m}(S)\), by reindexing \(\gamma\).
Write \(G_S=\operatorname{GL}_p(S)\), \(P=1-ee^{\mathsf t}\), and set
\[
 H=\{A\in G_S: AP=P,\ w\bar A=w\}.
\]
This is a subgroup, since both conditions are stabilizer conditions.
Let \(K\) be the full integral level of the definite unitary group
\(U(I_p)\), and \(K_H\) the inverse image of \(H\) in \(K\).
The reduction map from \(K\) onto \(G_S\) is surjective. At each
chosen split place lift the residue matrix arbitrarily to an
integral matrix. Its determinant is a unit, so the inverse
transpose gives the integral action on the dual component. Use
the identity at the unmarked places.

For a \(p\)-by-\(s\) matrix \(x\) over \(S\), set
\[
 D_w(x)={\bf1}_{\,w\bar x=w\bar e}\,\chi(\det(e^{\mathsf t}x)).
\]
For \(z\in\operatorname{Mat}_{p,p}(S)\) define
\[
 \begin{split}
 \Phi(y,z)&=\begin{cases}
       \chi(\det z)F_\Lambda(z^{-1}y),&z\in G_S,\\
       0,&z\notin G_S,
                 \end{cases}\\
 \Phi_H(y,z)&=D_w(ze)\,{\bf1}_{\,zP=P}\,F_\Lambda(y).
 \end{split}
\]
Write \(\Pi(A)f(y,z)=\chi(\det A)f(A^{-1}y,A^{-1}z)\).

\begin{lemma}[Division into cosets with coefficient one]
\leavevmode\label{lem:rowcosets824}
One has \(\Phi_H(y,z)={\bf1}_{z\in H}\Phi(y,z)\) and
\[
                 \Phi=\sum_{A H\in G_S/H}\Pi(A)\Phi_H .
\]
Each term is independent of the representative of its coset.
\end{lemma}
\begin{proof}
In coordinates with the \(e\)-columns first, \(zP=P\) means
that \(z\) has block form
\(\left(\begin{smallmatrix}D&0\\ B&1\end{smallmatrix}\right)\).
Thus \(\det z=\det D=\det(e^{\mathsf t}ze)\), and
\(w\bar z=w\) is equivalent to \(w\overline{ze}=w\bar e\).
A nonunit determinant contributes zero. If \(z\in H\), then
\(\Lambda\bar z^{-1}=\Lambda\) because \(\Lambda=uw\). This gives
\(F_\Lambda(z^{-1}y)=F_\Lambda(y)\) and proves the first equality.
The determinant multiplier cancels under \(\Pi(A)\) on \(\Phi\),
so \(\Pi(A)\Phi=\Phi\). Hence \(\Pi(A)\Phi_H\) is \(\Phi\)
restricted to \(z\in AH\). These supports partition the invertible
matrices, and all terms vanish elsewhere.
\end{proof}

For a level \(J\subset K\), let \(\mathcal C_J\) be a set of
representatives of \(U(I_p)(F)\backslash U(I_p)(\mathbb A_f)/J\).
For \(h\in\mathcal C_J\) put
\(d_h^J=|h^{-1}U(I_p)(F)h\cap J|\). The stabilizers are finite
because the real group is compact and the rational group is
discrete in the adelic group. These class sets are finite by the
finiteness theorem for definite lattices, as in
Theorem~\ref{thm:definiteclass822}. Write
\(\mathcal A_J(f)=\sum_{h\in\mathcal C_J}f(h)/d_h^J\).

Let \(\Theta_{\Xi}(g,h',h)\) be the theta kernel from
\(I_{3m}\oplus W'\) to \(I_p\), with finite datum \(\Xi(y,z)\)
and real polynomial \(P(X,Z;Y_0)\) of
Proposition~\ref{prop:stackednorm819}. The variables \(y,z\) are
the columns from \(I_{3m},W'\), respectively. The source point
has entries \(h',g,1,1,b_0\) on the five summands, with identity
real part on the definite class variables. The target point has
finite part \(h\) and identity real part. In the convention of
Lemma~\ref{lem:signkernel823}, the two splitting characters
are \((\alpha,\nu)=(\nu\upsilon_1^m,\nu)\). The compact
restriction of \(\alpha\) is \(\chi\). At a real embedding its
exponent is \(m-\mu_\tau\), where \(\mu_\tau\) counts the
determinant factors of \(P\) at that embedding. Indeed \(\nu\)
has exponent \(-\mu_\tau\) there, and
\(\nu\upsilon_1^m/(\chi_M^j\vartheta)\) has trivial tame part.

\begin{proposition}[Refinement of the definite class integral]
\leavevmode\label{prop:refinedmass824}
For every \(g,h'\) the following unnormalized class sums agree:
\[
 \mathcal A_K\bigl(\Theta_{\Phi}(g,h',\cdot)\bigr)
     =\mathcal A_{K_H}\bigl(\Theta_{\Phi_H}(g,h',\cdot)\bigr).
\]
\end{proposition}
\begin{proof}
At the marked split places the integral target change \(k\in K\)
acts on the position function by pullback by its inverse, times
the character of its determinant. This follows by substituting
in the position coordinates of a self-dual lattice box. The
absolute determinant multiplier is one since \(k\) preserves
the lattice. The marked splitting character gives
\(\chi(\det k_S)\), where \(k_S\) is the reduction of \(k\).
The unmarked unit vacuum is fixed. It follows that
\(\Theta_{\Pi(k_S)\Xi}(g,h',h)=\Theta_\Xi(g,h',hk)\).

Both kernels are left invariant under rational changes of the
finite target point. To see this, first use rational automorphy
on the full adelic target. The simultaneous change at infinity
can be dropped, as the vector there is fixed. Indeed, for
\(v\in U(p)\) at one real place, substitution on the \(p\)
variable rows multiplies each stacked determinant by \(\det v\):
the operation on its rows is \(\operatorname{diag}(1_{2m},v^{\mathsf t})\).
In the Fock splitting the vacuum contributes
\((\det v)^{-m}\), one inverse determinant per negative source
column. Multiplication by the external character contributes
\((\det v)^{m-\mu_\tau}\). The product with the \(\mu_\tau\)
polynomial factors is one. This calculation holds at each real
place. The lemma now gives right invariance at \(K\) for the
first kernel, at \(K_H\) for the second, and expresses the
first as the sum of right translates of the second over \(K/K_H\).

Fix a full-level representative \(h\) and put
\(\Gamma=h^{-1}U(I_p)(F)h\cap K\). The refined classes above
\(h\) correspond to the \(\Gamma\)-orbits on \(K/K_H\) by
\(kK_H\mapsto hk\). The summand is constant on each orbit
by the two invariances. The contribution of an orbit, divided
by \(|\Gamma|\), is its common value divided by the order of
\(\Gamma\cap kK_H k^{-1}\). Conjugation by \(k\) identifies
this last group with the rational stabilizer at the refined
class \(hk\). Summing all the orbits gives the equality. In
particular this refinement requires no division by \([K:K_H]\),
even at primes dividing that index.
\end{proof}

Use the unscaled first definite kernel of
Theorem~\ref{thm:signedseries823} to set
\(C_\vartheta(h')=\sum_{v\in\mathcal C_m}K_0(1;v,h')/d_v\).
Define the following function of the actual adelic variable
\(g\in U(W_j)(\mathbb A)\):
\[
 L_{i,\vartheta}(g)=
 \sum_{h'\in\mathcal C_{3m}}\frac{C_\vartheta(h')}{d_{h'}}
                  \mathcal A_K\bigl(\Theta_\Phi(g,h',\cdot)\bigr).
\]
For \(\xi\in D(c)\otimes\CC\) write \(\mathcal L(\xi)\) for
\(T_c\xi\) evaluated on the positive columns of the \(W_j\)
variable in their flag order. All active heights are at most
\(a\), so the other columns of its formal square argument may
be filled with zero. Write \(\Theta_j[D,Q](g,h)\) for the
theta kernel from \(W_j\) to \(I_p\) with characters
\((\upsilon_1^{m-j},\nu)\), residue datum \(D\) and polynomial \(Q\).

\begin{theorem}[The arithmetic coefficient in the theta class integral]
\leavevmode\label{thm:arithmeticclass824}
With the single scalar \(S_j\) and the Weyl vectors
\(\mathcal Y_\vartheta(h)\) of Theorem~\ref{thm:signedseries823},
\[
 S_j L_{i,\vartheta}(g)=
    \sum_{h\in\mathcal C_{K_H}}\frac{1}{d_h^{K_H}}\,
       \Theta_j\!\left[D_w,\,
          \mathcal L\bigl(\operatorname{conj}\mathcal Y_\vartheta(h)\bigr)
                     \right](g,h).
\]
Conjugation uses the real structure of the integral Weyl basis.
The equality also holds after taking any of the cuspidal product
periods of Theorem~\ref{thm:complexseparation818}, using the
same character projection on both sides.
\end{theorem}
\begin{proof}
Refine the inner class sum by
Proposition~\ref{prop:refinedmass824}. Reorder the source summands
as \(W_j,-I_j,I^{\rm add}\), with
\(I^{\rm add}=I_{3m-j}\oplus I_{6m+1}\oplus I_{3m}\).
On these blocks \(\Phi_H\) is a product. Its first factor is
\(D_w\). On the \(j\) negative columns it requires the
inclusion columns. On the first two added blocks it imposes
the remaining inclusion columns, and on the last \(3m\)
columns its value is \(F_\Lambda(y)\). These are precisely
the negative and added tests used in
Theorem~\ref{thm:signedseries823}. The added source point is
\(\operatorname{diag}(1,b_0,h')\).

In any integral basis \((b_\beta)\) of \(D(c)\), write
\(\zeta(Z;Y_0)=\sum_\beta b_\beta\,\zeta_\beta(Z;Y_0)\).
The polynomial identity of Proposition~\ref{prop:stackednorm819}
gives \(P=\sum_\beta\mathcal L(b_\beta)\zeta_\beta\).
The negative block has polynomial one. On an orthogonal
direct sum the three theta kernels multiply for each of these
pure tensors. In detail, take the direct sum of rational
polarizations of the three Hom spaces. The summation set in
the position model is the Cartesian product of the three
rational position spaces. The product Schwartz test gives
the product of their summands. Each rational Schwartz sum is
absolutely convergent at a fixed adelic argument. The
factorization persists after the pair action, since the
oscillator action is the tensor product in this polarization,
with the unit vacua and Gaussian transforms normalized by
products. The three characters on the shared target multiply
to \(\upsilon_1^{m-j}\upsilon_1^j\nu=\alpha\), and the
character \(\nu\) on the block source determinant factors over
the blocks. Thus the normalizations are the ones used in the
three individual kernels.

Now interchange the two finite class sums. At a given \(h\),
the negative scalar kernel and the sum over \(h'\) of
\(C_\vartheta(h')/d_{h'}\) times the added kernel on
\(\zeta_\beta\) are exactly the \(\beta\)-coordinate of
\(\Upsilon_\vartheta(h)\) in Theorem~\ref{thm:signedseries823}.
Linearity of \(\Theta_j\) in the real-place polynomial gives
the asserted sum with \(\mathcal L(\Upsilon_\vartheta(h))\)
before multiplication by \(S_j\). That theorem says
\(S_j\Upsilon_\vartheta(h)=\operatorname{conj}\mathcal Y_\vartheta(h)\),
which proves the formula including the location of the conjugation.

Finally, multiplication by the complementary test, translation,
and integration against a fixed cusp form are linear in the
factor just computed (or conjugate-linear if the conjugate
pairing is used on that factor). The fixed theta tests have
moderate growth and the cusp form is rapidly decreasing on
Siegel sets, as in Theorem~\ref{thm:complexseparation818}.
The products are therefore integrable. Finite class sums
and the finite character projection commute with the integral.
In a conjugate-linear pairing \(S_j\) is conjugated on passing
outside that pairing.
\end{proof}

The theorem inserts the previously bounded algebraic vectors
into an adelic function on the prescribed \(U(W_j)\). Identifying
this function, up to the required scalar, with the separating
carrier integral \(\mathcal B_i\) still requires the global
comparison on the other side of the theta exchange.
\clearpage

\subsection{The matching of the residue data and the carrier integral}
\label{sec:carrierdata825}
Keep the field, interval, \(i,j,\vartheta\), and the matrices
of Section~\ref{sec:refinedclass824}. Write \(r=3m\) and
\(n=p+r\). We identify the actual data on the two sides of the
rank comparison. On the carrier side this includes the
multiplicity with which a single elementary character occurs
in the integral.

\begin{lemma}[Complementary minors over the marked residue ring]
\leavevmode\label{lem:residuepartners825}
Over \(S=\OO_E/(\mathfrak a^2\mathfrak q)\), put
\[
 \Psi\binom{A}{B}=\chi(\det A)F_\Lambda(B),
 \qquad A\in\operatorname{Mat}_{r,r}(S),\
        B\in\operatorname{Mat}_{p,r}(S).
\]
Let the first \(r\) columns of \(k\in\operatorname{GL}_n(S)\)
be \(\binom{A}{B}\), and the last \(p\) rows of \(k^{-1}\)
be \((y\,\,z)\). Then
\[
       \Phi(y,z)=\chi(\det k)^{-1}\Psi\binom{A}{B}.
\]
Every point where \(\Psi\) is nonzero occurs as these columns
of some \(k\). Every point where \(\Phi\) is nonzero occurs
as these rows of the inverse of some \(k\).
\end{lemma}
\begin{proof}
Set \(D=\left(\begin{smallmatrix}1&0\\y&z\end{smallmatrix}\right)\).
The bottom rows of \(Dk\) are \((0\,\,1)\) and its upper
left block is \(A\). Computing the determinant of \(Dk\) in
two ways gives \(\det z\det k=\det A\) over the given ring,
including its old nonreduced factors. Thus \(A\) and \(z\) are
simultaneously invertible. If they are singular in the sense
of having nonunit determinant, both functions vanish. If they
are invertible, \(yA+zB=0\) gives \(z^{-1}y=-BA^{-1}\).
The right invariance of \(F_\Lambda\) under \(-A^{-1}\)
and the determinant identity now give the formula.
For the support assertions, complete the columns by
\(\left(\begin{smallmatrix}A&0\\B&1\end{smallmatrix}\right)\)
when \(A\) is invertible. Complete the rows by taking
\(k^{-1}=\left(\begin{smallmatrix}1&0\\y&z\end{smallmatrix}\right)\)
when \(z\) is invertible.
\end{proof}

The real data associated to these complementary dimensions can
be compared entry by entry. Let \(S_0=I_{2m}\oplus(-I_m)\)
denote the carrier Hermitian space. At an embedding \(\tau\)
let \(X_\tau^+\) be the \((p+2m)\)-by-\(2m\) matrix from
the positive coordinates of \(S_0\) to the positive coordinates
of \(I_r\oplus W'\). Let \(X_\tau^-\) be its separate
\(m\)-by-\(m\) negative-to-negative block. Give the rows in
each block the fixed order used for the minors of
Corollary~\ref{cor:paddedminor814}. For each of the
\(\mu_\tau\) determinant factors in the stacked polynomial
of Proposition~\ref{prop:stackednorm819}, let \(C_{\tau l}\)
be its \(2m\)-by-\((p+2m)\) matrix of constant rows.
The polynomial attached to these rows on the smaller source is
\[
 Q=\prod_\tau\left\{
     \det(X_\tau^-)^{\,\mu_\tau+n-2m}
                 \prod_{l=1}^{\mu_\tau}\det(C_{\tau l}X_\tau^+)
                         \right\}.
\]

\begin{proposition}[The real polynomial on the carrier space]
\leavevmode\label{prop:carrierpolynomial825}
The variables in \(Q\) belong entirely to \(S_0\) and \(W'\).
On these spaces \(Q\) is the product \(vQ_{b_0,\rho}\) of the
padded polynomial and the short-row polynomial defining
\(\mathcal B_i\) in Corollary~\ref{cor:adelicfirst816}.
\end{proposition}
\begin{proof}
All \(m\) negative rows lie in \(W_j\oplus(-I_j)\). Their
exponent is \(\mu_\tau+p+m=\kappa_\tau'\), the negative
exponent in Lemma~\ref{lem:carriercoefficient815}.
For a positive diagram column of height \(h<a\), row \(t\)
of \(C_{\tau l}\), starting with \(t=0\), selects coordinate
\(h+t\) in \(W_j\) when \(h+t<a\), and coordinate \(t\) in
\(I_{3m-j}\) otherwise. Therefore \(C_{\tau l}X_\tau^+\)
lists exactly those rows of the variable matrix and its
determinant is the corresponding padded minor, in its given
order. Its multiplicity is the prescribed \(c^+_{\tau,h}\).
For height \(a\) the \(2m\) rows of \(C_{\tau l}\) are the
selected rows of the short matrix at \(b_0\), supported entirely
on \(I_{6m+1}\). Multiplication against \(X_\tau^+\)
gives the \(2m\)-by-\(2m\) short-row determinant of
\(Q_{b_0,\rho}\). There is one such factor for each prescribed
selection, including any repeated selections. These exhaust
the factors by the description preceding
Proposition~\ref{prop:stackednorm819}. None uses a coordinate
in \(I_r\). The fixed orders of the rows give equality with
the stated product, including its sign.
\end{proof}

Let \(\phi\) be the carrier of Theorem~\ref{thm:carrier813}.
For a finite residue function \(f\) on the \(p\)-by-\(r\)
position box let \(\mathcal I_i(f;g)\) be the integral of
\(\phi(x)\) times the kernel from \(S_0\) to \(W'\) with
datum \(f\), real polynomial \(Q\) just computed and target
point \(\operatorname{diag}(g,1,1,b_0)\). Integrate on
\(U(S_0)(F)\backslash U(S_0)(\mathbb A)\) with the quotient
measure defining \(\mathcal B_i\). Use the same unmarked
vacua and characters \((\upsilon_1^m,\beta_0)\) as there.
In particular the character \(\beta_0\) applied to the
source determinant is trivial on its finite integral units.
For the finite box we allow either the old modulus or the
old modulus enlarged by \(\mathfrak q\).

\begin{theorem}[The exact multiplicity in the carrier period]
\leavevmode\label{thm:carriermultiplicity825}
For every \(g\in U(W_j)(\mathbb A)\),
\[
   \mathcal I_i(F_\Lambda;g)=N_R\mathcal B_i(g),
   \qquad N_R=|\operatorname{GL}_{3m}(R)|.
\]
The integer \(N_R\) is the same for every new twist on the
given field and interval. With \(M=p_0p_1\) one has
\(\log N_R\le18fm^2\log M\).
\end{theorem}
\begin{proof}
First consider \(f_\Lambda(B)=
\psi_R(\operatorname{tr}(\Lambda B))\) at the old modulus.
Proposition~\ref{prop:carrierpolynomial825} and the definition
of the first factor give
\(\mathcal I_i(f_\Lambda;g)=\mathcal B_i(g)\).
Inflating this test from \(R\) to \(S\) gives the identical
finite Schwartz vector. On the old integral boxes the value
is read on the same coset. At a new mark the pullback is
constant in that component on its integral position box and
zero outside. This is the unit vacuum in the split position
model. Thus inflation has scalar one.

For \(\gamma\in\operatorname{GL}_{r}(R)\) lift \(\gamma\)
at the old split marks to an element \(k_\gamma\) of the full
finite lattice stabilizer of \(U(S_0)\). The dual component
is forced by the Hermitian pairing, and away from the old
marks take \(k_\gamma=1\). The split position action for a
source change substitutes \(B\mapsto B\gamma\). The lattice
Jacobian has absolute value one, and the source determinant
multiplier is one by the compact restriction of \(\beta_0\).
The integral with test \(f_\Lambda(\cdot\,\gamma)\) therefore
has as integrand the original integrand translated by
\(x\mapsto xk_\gamma\): the carrier is fixed at this full
level by Theorem~\ref{thm:carrier813}. Right invariance of
the quotient measure gives the same integral. The carrier
is cuspidal, hence rapidly decreasing on Siegel sets, and
the theta function is of moderate growth for fixed target
point. Thus each integral is absolutely convergent.
Summing over all \(\gamma\) gives the formula by linearity,
also when \(\Lambda\) has a nontrivial stabilizer.

At an old split component with residue field of size \(q\),
the reduction from \(\operatorname{GL}_r(\OO/\varpi^2)\)
onto \(\operatorname{GL}_r(\OO/\varpi)\) has kernel
\(1+\varpi\operatorname{Mat}_r(\OO/\varpi^2)\) of order
\(q^{r^2}\). The residue group has order
\(\prod_{h=0}^{r-1}(q^r-q^h)\), by successive choices of
linearly independent columns. Multiplication over the old
components gives \(N_R\); in particular
\(N_R\le |R|^{r^2}=M^{2fr^2}\), as claimed.
\end{proof}

For the untwisted character, Theorem~\ref{thm:complexseparation818}
now applies to the family of functions
\(\mathcal I_i(F_\Lambda;\cdot)\) itself: replacing every
\(\mathcal B_i\) in a nonzero projected product period by this
function multiplies the period by the same positive integer
\(N_R\). In particular at least one of the matched residue and
real tests above has a nonzero carrier product period for each
nonzero good form. The comparison with
\(L_{i,\vartheta}\) requires an identity of global integrals;
the pointwise relation of residue functions in
Lemma~\ref{lem:residuepartners825} supplies its finite tests.
\clearpage

\subsection{Constancy of the first equal-rank theta average}
\label{sec:equalrankaverage826}
The first of the three changes of rank involves two copies of
the definite space \(I_m\). We give a direct argument for the
dependence of its class integral on the source point. The
argument works at the old marks and at the new twist conductor.

Write \(G=U(I_m)\) over \(F\), and \([G]=G(F)\backslash
G(\mathbb A)\). Take a unitary character \(\delta\) of the
adelic norm-one group, trivial on rational points and of
weight zero at every real place. Its finite restriction to
units is unramified outside the marked split places and is
given there by a residue character \(\chi_\delta\). The
application takes \(\delta=(\chi_M^j\vartheta)^{-1}\).
Use the pair \((I_m,I_m)\), splitting characters
\((\delta,\delta)\), and real polynomial one. At every
unmarked place take the standard unit vacuum. At a marked
split place \(v\) use the position test on
\(\operatorname{Mat}_m(F_v)\)
\[
 f_v(X)=
 \begin{cases}
   \chi_{\delta,v}(\det(X\bmod\mathfrak m_v)),&
                           X\in\operatorname{Mat}_m(O_v),\\
   0,&\text{otherwise},
 \end{cases}
\]
where characters of the finite ring are zero on nonunits.
Let \(\theta_\delta(g,h)\) be the resulting rational theta
sum, and let \(I_\delta(g)=\int_{[G]}\theta_\delta(g,h)\,dh\).
With Haar mass one on \(G(F_\infty)\) and on the full finite
stabilizer, this integral is the definite target class sum.
Indeed these groups fix the target vector and the orbit of
a finite representative has volume the reciprocal of its
finite rational stabilizer order. An overall change in the
quotient measure multiplies all these integrals by the same
constant.

\begin{lemma}[The two spherical averages at a split place]
\leavevmode\label{lem:splitaverage826}
Let \(v\) be any split finite place, let \(K_v=\operatorname{GL}_m(O_v)\)
and let \(d\in\operatorname{GL}_m(F_v)\).
Write \(\omega_s,\omega_t\) for the source and target actions,
respectively, in the local position model. For \(a\) in a
double coset use its normalized invariant measure. Then
\[
   \int_{K_v d K_v}\omega_s(a) f_v\,da
       =\int_{K_v d^{-\mathsf t} K_v}\omega_t(a) f_v\,da .
\]
At an unmarked split place \(f_v={\bf1}_{\operatorname{Mat}_m(O_v)}\).
\end{lemma}
\begin{proof}
Choose the same split component for the two copies of \(I_m\).
The position matrix transforms in the test by
\(X\mapsto h^{-1}Xg\). Let \(Tf(X)=f(X^{\mathsf t})\).
We claim \(T\omega_s(a)T^{-1}=\omega_t(a^{-\mathsf t})\).
For the unitary pullbacks on positions, transposition conjugates
\(X\mapsto Xa\) to \(X\mapsto a^{\mathsf t}X\), and the
absolute Jacobians agree. This also gives agreement for the
chosen Weil splitting. In fact the split unitary group of
the whole Hom space is a general linear group on the
\(m^2\) positions. On this group unitary pullback with the
square-root Jacobian implements its action on the Heisenberg
translations. Any other splitting implementing that action
differs by a scalar character, by irreducibility of the
position representation and the multiplication law. A character
has the same value on conjugate position changes. Finally,
the external source multiplier is \(\delta_v(\det a)^{-1}\)
and the target multiplier at \(a^{-\mathsf t}\) is
\(\delta_v(\det a^{-\mathsf t})\). They agree.

The local vector is fixed by both compact actions. On integral
units the position Jacobian and the unramified Hom splitting
are trivial. For the marked vector the substitution on the
source supplies \(\chi_{\delta,v}(\det k)\), and that on the
target its inverse. The two respective external multipliers
cancel them. The unmarked assertion follows by lattice
invariance. More generally, every Schwartz vector \(b\)
fixed by both compact actions satisfies
\[
 b(uXv')=\lambda(\det u)\lambda(\det v')b(X)
               \qquad (u,v'\in K_v),
\]
where \(\lambda\) is the restriction of \(\delta_v\) to units.
The same calculation of the two compact actions gives this
formula for any \(b\). By Smith normal form, every matrix over
\(F_v\) is \(u D v'\) with \(D\) diagonal and \(u,v'\in K_v\),
also for deficient rank: first clear a common denominator,
diagonalize over the discrete valuation ring and divide it
out again. The factors in the displayed formula are unchanged
on passing from \(uDv'\) to \((v')^{\mathsf t}Du^{\mathsf t}\).
Thus \(Tb=b\) on this entire space of compact-fixed vectors.

The average on the left of the lemma is fixed by the source
compact group by invariance of double-coset measure and by
the target compact group since the two actions commute.
Transposition therefore fixes this average. Applying \(T\)
inside it, using \(Tf_v=f_v\), and changing \(a\) to
\(a^{-\mathsf t}\) proves the equality. This automorphism
preserves \(K_v\) and carries one double-coset probability
measure to the other.
\end{proof}

\begin{theorem}[Constancy on the full definite source quotient]
\leavevmode\label{thm:globalequalrank826}
For every twist and tame character just specified,
\(I_\delta(g)=I_\delta(1)\) at every adelic source point.
In particular the actual square class sum with character
\((\chi_M^j\vartheta)^{-1}\) in the first exchange is a
constant function.
\end{theorem}
\begin{proof}
Both rational quotients are compact, since the group is
anisotropic over \(F\). Rational theta sums of Schwartz tests
are continuous and bounded on their product. One can see
continuity on compact sets of representatives by rapid
Schwartz decay in a rational position lattice, with the
finite support contained in a fixed fractional lattice
on each such set. Consequently \(I_\delta\) is continuous
and square integrable. For an unmarked nonsplit finite place
the lattice vacuum and the unramified characters make
\(I_\delta\) fixed by the full integral source stabilizer.
At split places this invariance follows from the compact
calculation in Lemma~\ref{lem:splitaverage826}. At a real place
the compact definite source and target both fix the Gaussian:
the Hom signs are all positive, so the untwisted Fock vacuum
has compact weight zero in the splitting of
Theorem~\ref{thm:signedseries823}, and the external weights
are zero.

Fix a split \(v\) and \(d\). Apply the lemma before rational
summation and target integration, with the other components
held fixed. The local double-coset average of a smooth
compact-fixed vector is a finite average of translates,
so it commutes with both operations. On the target side
each right translate gives the same integral by right
invariance of quotient Haar measure. We obtain
\[
     \int_{K_vdK_v} R(a)I_\delta\,da=I_\delta,\qquad
                       R(a)I_\delta(g)=I_\delta(ga).
\]
Right translation on \(L^2([G])\) is unitary. In particular
\[
 \int_{K_vdK_v}\|R(a)I_\delta-I_\delta\|_2^2\,da
     =2\|I_\delta\|_2^2-
       2\operatorname{Re}\left\langle
        \int_{K_vdK_v} R(a)I_\delta\,da,I_\delta\right\rangle=0.
\]
The integrand is continuous and the double-coset measure has
full support. It follows that \(R(a)I_\delta=I_\delta\) for
each \(a\) in this coset. Continuity on the rational quotient
upgrades equality in \(L^2\) to pointwise equality. Varying
\(d\) gives invariance under the full group at every split
finite place.

Let \(J\) be the subgroup generated by these split local
groups, all finite integral stabilizers, and the full real
group. We finish by showing \(G(\mathbb A)=G(F)J\).
For \(g\in G(\mathbb A)\), the element
\(\operatorname{diag}(\det g,1,\ldots,1)\) belongs to \(J\).
At a nonsplit finite place a norm-one element is a unit.
At the split places there are only finitely many nonintegral
components. Dividing on the right by this diagonal element
puts \(g\) in \(G^1(\mathbb A)\) for \(G^1=\operatorname{SU}(I_m)\).

Choose one of the old split places \(v_0\). For \(m\ge2\),
\(G^1\) is simply connected and absolutely almost simple:
over an algebraic closure it is \(\operatorname{SL}_m\).
At \(v_0\) its local group \(\operatorname{SL}_m(F_{v_0})\)
is noncompact. The strong approximation theorem for simply
connected groups over number fields therefore makes \(G^1(F)\)
dense in the finite adeles of \(G^1\) away from \(v_0\)
(take the omitted places to be \(v_0\) and all real places).
The full hypothesis on the omitted places is precisely
noncompactness of their product for this almost simple group.
We use the number-field theorem of
\href{https://www.mathnet.ru/eng/im2220}{V.~P.~Platonov,
\emph{The problem of strong approximation and the Kneser--Tits
conjecture for algebraic groups}}.
The integral stabilizers away from \(v_0\) form an open
subgroup of this restricted product. Thus for any such \(g\)
of determinant one there is \(\gamma\in G^1(F)\) with
\(\gamma^{-1}g\) integral away from \(v_0\) at the finite
places. Its \(v_0\) and real components already belong to
\(J\). This proves the factorization. In rank one the diagonal
determinant argument alone gives it. Rational left invariance
and right \(J\)-invariance now imply \(I_\delta(g)=I_\delta(1)\).
\end{proof}

This proves the constant-function assertion before normalizing
the first exchange. Its value is the scalar
\(I_\delta(1)\). Using its reciprocal in the chain requires
nonvanishing and a formula for how its absolute value varies
with the tame twist.
\clearpage

\section{The first and second theta exchanges}\label{nav:first-second-exchange}

\subsection{The relative finite doubling factors of the first exchange}
\label{sec:squarefactors827}
Let \(m=2d\) and retain the positive definite pair of
Theorem~\ref{thm:globalequalrank826}. The variation with the
twist at all split places can be computed before using a
global inner product formula. Normalize multiplicative Haar
measure on \(\operatorname{GL}_m(k)\) by giving
\(K=\operatorname{GL}_m(O)\) mass one, and additive Haar measure
on \(\operatorname{Mat}_m(k)\) by giving the integral matrices
mass one. Here \(k=F_v\), \(q\) is its residue cardinality,
and \(\varpi\) a uniformizer at a split finite place. Write
\[
  P_m(q)=\prod_{j=1}^m(1-q^{-j}),\qquad
  H_m(q,s)=\prod_{j=1}^m(1-q^{-2s-j}).
\]
For \(a\in\operatorname{GL}_m(k)\) with elementary divisor
valuations \(b_1,\ldots,b_m\), set \(\ell(a)=\sum_i|b_i|\).
Let \(C_{\delta,v}(a)=\int(\omega_s(a)f_v)(X)\overline{f_v(X)}\,dX\),
where the subscript on \(\omega_s\) means the source action.
The complex parameter in the next formula occurs only in the
scalar \(q^{-s\ell(a)}\). Put
\[
      Z_{\delta,v}(s)=\int_{\operatorname{GL}_m(k)}
                  C_{\delta,v}(a)\,q^{-s\ell(a)}\,d^\times a .
\]

We record why this is the flat-section doubling integral for
the constant source vector. At the theta point \(s=0\),
evaluation of the doubled Schwartz tensor
\(f_v\otimes\overline{f_v}\) on the diagonal graph is the
inner product of its two factors: in diagonal and difference
coordinates set the difference to zero and take the zero
Fourier coefficient in the diagonal coordinate. This is the
calculation of Proposition~\ref{prop:doubledkernel711}.
At \(a\) it gives the displayed \(C_{\delta,v}(a)\), including
the external source character. Moving the section from zero to
\(s\) multiplies its value at each point by
\(\lvert\det A/\det D\rvert^s\), where \(A,D\) are the
diagonal Levi blocks of the split Siegel parabolic in a
parabolic-times-compact factorization.

Here that absolute value is \(q^{-\ell(a)}\). Indeed choose an
integral change of basis in the doubled split component carrying
the diagonal plane to the first \(m\) coordinates. In these
coordinates the bottom rows of the embedded point
\(\operatorname{diag}(a,1)\) have the form \((1-a,1)\) up to
integral invertible changes of basis, equivalently \((a,1)\).
For a matrix in parabolic-times-compact form the fractional
ideal of maximal minors of its bottom rows is \((\det D)\):
the last rows of the compact factor span a direct summand.
For \((a,1)\) the same ideal has valuation
\(\sum_i\min(0,b_i)\), by diagonalizing \(a\) with integral
units. The determinant of the full point has valuation
\(\sum_i b_i\). Thus the valuation of \(\det A/\det D\) is
\(\sum_i b_i-2\sum_i\min(0,b_i)=\ell(a)\).
This calculation also fixes which normalization of the
complex variable is used below.

\begin{theorem}[The marked and unmarked split integrals]
\leavevmode\label{thm:splitzeta827}
Put \(\nu=\delta^{-1}\upsilon_1^d\). At an unmarked split
place write \(u=\nu_v(\varpi)\), reading the norm-one
torus on the chosen component. Then
\[
 Z_{\delta,v}(s)=
 \frac{H_m(q,s)}
 {\displaystyle\prod_{j=1}^m
       (1-uq^{d-j-s})(1-u^{-1}q^{d-j-s})}.
\]
This holds for \(\operatorname{Re}s>d\) and by meromorphic
continuation. At a marked split place with the prescribed
determinant test one has \(Z_{\delta,v}(s)=P_m(q)\) for all \(s\).
\end{theorem}
\begin{proof}
At a mark the support of \(f_v\) is \(K\), regarded as a
subset of the matrix space. The support of a source translate
is \(Ka^{-1}\), disjoint from \(K\) unless \(a\in K\).
When \(a\in K\) the full pair action fixes the test by
Lemma~\ref{lem:splitaverage826}. Its squared norm is the
additive volume \(P_m(q)\) of \(K\), by reduction modulo
\(\varpi\) and the count of invertible residue matrices.
Since \(\ell=0\) on \(K\), the assertion follows at every
complex parameter. This includes an old marked test written
modulo \(\varpi^2\) with a tame unit character.

At an unmarked place the source action on positions is
\(f(X)\mapsto \xi(\det a)|\det a|^d f(Xa)\), with \(\xi\)
an unramified unitary character. To identify it in the
lattice convention, take \(a=\operatorname{diag}(\varpi,1,\ldots,1)\).
On the Hom points the action is \(L\mapsto La^{-1}\).
At the chosen component it has \(m\) elementary divisors
\(\varpi^{-1}\) and at its conjugate \(m\) divisors
\(\varpi\), all others being units. In the vacuum overlap
of Lemma~\ref{lem:signkernel823}, the intersection with the
inverse image lattice thus has determinant ideal
\(\varpi^m\) on the chosen component and a unit on the other.
The overlap \(1/\overline{\Delta}\) is
\(\upsilon_v^{E}(\varpi)^m q^{-m/2}\), where
\(\upsilon_v^{E}\) is the splitting character of the chosen
\(E\)-component. On the two split components the restrictions
of this character are inverse, since its restriction to the
diagonal \(k^\times\) is trivial. Therefore
\(\upsilon_v^{E}(\varpi)^m=\upsilon_{1,v}(\varpi)^d\).
The external multiplier is \(\delta_v(\varpi)^{-1}\).
This gives \(\xi(\varpi)=u\). For general \(a\) the unit
vacuum has phase \(u^{v(\det a)}\), and its overlap in
absolute value is \(q^{-d\ell(a)}\), by diagonalization and
intersection of the two position lattices. Hence the desired
integral is the integral of
\(u^{v(\det a)}q^{-(d+s)\ell(a)}\).

For an explicit evaluation put \(t=d+s\) and
\(J(2s)=\int_{\operatorname{Mat}_m(O)}|\det X|^{2s}\,dX\).
For fixed \(a\) one has
\[
 |\det a|^t
 \int_{\substack{X\ {\rm integral}\\ Xa\ {\rm integral}}}
                 |\det X|^{2s}\,dX
            =J(2s)q^{-t\ell(a)} .
\]
To verify this, reduce \(a\) to its elementary divisors.
Column \(i\) of \(X\) must be divisible by
\(\varpi^{\max(0,-b_i)}\). Scaling this column contributes
\(q^{-(m+2s)\max(0,-b_i)}\), which proves the identity
because \(m+2s=2t\). Singular \(X\) have additive measure
zero, as follows by successively choosing independent columns.
Integrate the identity against \(u^{v(\det a)}d^\times a\)
and put \(Y=Xa\). Using \(dX=P_m(q)|\det X|^m d^\times X\)
gives
\[
 J(2s) Z_{\delta,v}(s)=P_m(q) D(t,u)D(t,u^{-1}),\qquad
 D(t,u)=\int_{K'\!}|\det X|^t u^{v(\det X)}\,d^\times X ,
\]
where \(K'=\operatorname{Mat}_m(O)\cap\operatorname{GL}_m(k)\).
These changes of variables are absolutely convergent for
\(\operatorname{Re}s>d\). Indeed the following integral
calculation with the real parts and \(|u|=1\) bounds both
factors. For additive probability on integral matrices the
generating function of \(v(\det X)\) is
\(\prod_{j=1}^m(1-q^{-j})/(1-yq^{-j})\) for \(|y|<q\).
To see this, the first column is \(\varpi^h\) times a
primitive column with probability \((1-q^{-m})q^{-mh}\).
Carry the primitive column to the first basis vector by an
integral row operation. Given this column, the other columns
are still independent uniform vectors after the operation.
Expanding the determinant reduces the problem to dimension
\(m-1\), with a factor \(y^h\). The geometric sum and
induction give the formula. Consequently
\[
 D(t,u)=\prod_{j=1}^m(1-uq^{m-j-t})^{-1},\qquad
 J(2s)=P_m(q)/H_m(q,s).
\]
Substitution proves the claimed rational function.
\end{proof}

Let \(\delta_0,\delta_1\) be two of the prescribed
weight-zero characters, with sets \(\mathcal S_i\) of
ramified split places (one chosen component per pair),
and write \(Q_i=\prod_{v\in\mathcal S_i}q_v\).
The common old marks are included in both sets.
Their tame characters, like the characters at the new
marks, are nontrivial on residue units. Put
\(\nu_i=\delta_i^{-1}\upsilon_1^d\), and let \(L_i(z)\)
be the finite Hecke \(L\)-function over \(E\) of
\(a\mapsto\nu_i(a/\bar a)\). Its ramified Euler factors
are one. Define
\[
             \mathcal R(s)=\prod_{v\ {\rm split}}
                    \frac{Z_{\delta_1,v}(s)}{Z_{\delta_0,v}(s)}
                 \qquad(\operatorname{Re}s>d).
\]

\begin{theorem}[The conductor power in the relative doubling integral]
\leavevmode\label{thm:squareconductor827}
The product \(\mathcal R\) has a meromorphic continuation
regular and nonzero at zero. Its absolute value there is
\[
       |\mathcal R(0)|=
        \left(\frac{Q_1}{Q_0}\right)^{d^2}
             \prod_{z=1}^d\left|\frac{L_1(z)}{L_0(z)}\right|^2 .
\]
\end{theorem}
\begin{proof}
At an unmarked split place the two factors in each pair in
Theorem~\ref{thm:splitzeta827} are the inverse local Euler
polynomials at the two conjugate places of \(E\): the two
base-change values on a uniformizer are \(u,u^{-1}\).
At a marked place the quotient by the same \(H_m(q_v,s)\)
is \(P_m(q_v)/H_m(q_v,s)\). At a nonsplit place the two
\(\delta_i\) agree locally, since the norm-one group there
is compact and the characters are unramified. Hence their
Euler factors cancel. For \(\operatorname{Re}s>d\) the
absolutely convergent Euler products yield
\[
 \mathcal R(s)=
 \frac{\prod_{v\in\mathcal S_1}P_m(q_v)/H_m(q_v,s)}
      {\prod_{v\in\mathcal S_0}P_m(q_v)/H_m(q_v,s)}
       \prod_{j=1}^{2d}\frac{L_1(s+j-d)}{L_0(s+j-d)} .
\]
The first fraction is regular with value one at zero.

The primitive Hecke functional equation gives the completed
functions \(\Lambda_i(z)=A_i^{z/2}G_\infty(z)L_i(z)\),
where both have the same archimedean factor since the twists
have weight zero. The finite conductor norm is \(Q_i^2\):
the base change is tamely ramified with exponent one on
each of the two split components at a mark, and unramified
elsewhere. Thus \(A_i=|\operatorname{disc}E|Q_i^2\). At
positive real integers \(z\), \(L_i(z)\ne0\). For \(z>1\)
this follows from the Euler product. At \(z=1\), let
\(\lambda\) be the base-change character and omit its
ramified factors in all three functions below. The Euler
logarithms give
\(\zeta_E(\sigma)^3|L(\sigma,\lambda)|^4
 |L(\sigma,\lambda^2)|\ge1\) for real \(\sigma>1\),
by \(3+4\cos t+\cos(2t)=2(1+\cos t)^2\ge0\).
The angular type of \(\upsilon_1^d\) is nonzero, so
\(\lambda\) and \(\lambda^2\) both have nonzero angular
type. Their Hecke functions are holomorphic at one while
the omitted-factor zeta function has a simple pole. A zero
of \(L(1,\lambda)\) would force the product to tend to zero,
which proves the assertion.

The integers \(j-d\) occur in pairs \(z,1-z\) for
\(1\le z\le d\). For real arguments the functional equation
relates \(\Lambda_i(1-z)\) to the complex conjugate of
\(\Lambda_i(z)\) by a factor of absolute value one.
The continued ratios at \(1-z\) are consequently regular
and nonzero, since their common archimedean factors cancel
as meromorphic functions. The contribution of a pair to
the absolute value of the product is
\[
        (A_1/A_0)^{z-1/2}\left|L_1(z)/L_0(z)\right|^2 .
\]
Here the two conductors on the completed positive side
give the power \(z\), and removing the completion at \(z\)
and \(1-z\) subtracts \(1/2\). Summing
\(z-1/2\) from \(1\) to \(d\) gives \(d^2/2\). This proves
the theorem.
\end{proof}

The theorem evaluates the relative product of the actual
split local data at the equal-rank point. Passing from this
continued product to the squared absolute values of the
constants \(I_{\delta_i}(1)\) requires the global doubled
identity at that point.
\clearpage

\subsection{Nonvanishing and normalization in the first theta exchange}
\label{sec:globalfirst828}
Retain \(m=2d\), the definite group \(G=U(I_m)\) and the
weight-zero character \(\delta\). Let \(\eta\) be the quadratic
Hecke character of \(F\) associated with \(E/F\). For each
prescribed \(\delta\), put \(Q_\delta=\prod_{v\in\mathcal S_\delta}q_v\)
and
\[
 \lambda_\delta(a)=(\delta^{-1}\upsilon_1^d)(a/\bar a),\qquad
 D_\delta=Q_\delta^{d^2}\prod_{z=1}^d|L_{E,f}(z,\lambda_\delta)|^2 .
\]
The subscript \(f\) denotes a product over finite places, with
the local factor one at a place where the character is ramified.

\begin{lemma}[The nonsplit and real factors of the square integral]
\leavevmode\label{lem:squarespherical828}
At an inert finite place \(v\) the normalized unramified
doubling integral for the square pair, with both full-level
vectors of norm one and compact volume one, is
\[
 Z_v(s)=
 \frac{\displaystyle\prod_{j=1}^{m} L_{E_w}(s+j-d,1)}
      {\displaystyle\prod_{j=1}^{m} L_{F_v}(2s+j,\eta_v^{\,j})}.
\]
At a real place the corresponding compact integral is a
positive constant independent of \(s\) and \(\delta\).
\end{lemma}
\begin{proof}
There is no finite ramification of \(E/F\) by
Lemma~\ref{lem:signkernel823}. At an inert place every
norm-one element is integral, so \(\delta_v\) is trivial.
The finite space is the standard unimodular one. Each pair
of basis lines of \(I_m\) splits a hyperbolic plane over
\(E_w\): a unit of norm \(-1\) supplies a primitive isotropic
vector, and the unramified residue trace lifts to a surjection
of integer rings. In detail choose an integral vector pairing
to one with the first. Subtract a multiple of the isotropic
vector with trace equal to the norm of this second vector to
make it isotropic as well. This argument also works in residue
characteristic two. In the doubled lattice the diagonal has
an integral isotropic complement, given by the vectors
\((t x,(t-1)x)\) for \(t\in\mathcal O_{E_w}\) with
\(t+\bar t=1\). Its basis with the diagonal has unit determinant.
Thus the doubled spherical section is the section for the
hyperspecial group, with value one at the identity by the
unit-vacuum overlap. The Hom lattice and the additive pairing
are self-dual in the prescribed finite frames.

We use the unramified doubling integral for a normalized
spherical vector. In the unitary convention it gives the
standard base-change Euler factor at \(s+1/2\) divided by
\(\prod_{j=1}^m L_{F_v}(2s+j,\eta_v^j)\) in this equal,
even-rank case. This is the spherical integral of
\href{https://doi.org/10.1515/crll.1992.428.177}{J.-S.~Li,
Theorem~3.1}, in the normalization of
\href{https://arxiv.org/pdf/1207.4709}{Gan, Qiu and Takeda,
\S11.6}. Its hypotheses here are unramified groups, characters,
hyperspecial vectors and the spherical section with value one.
In general the denominator uses the restriction of the
space-splitting character times \(\eta_v^{m-j}\) at the
\(j\)-th argument. That restriction is \(\eta_v^m=1\)
here, and \(m-j\) has the parity of \(j\).

For the constant representation the base-change Satake
eigenvalues are \(q_w^{(m+1)/2-j}\), \(1\le j\le m\).
Indeed the spherical constant occurs in the principal series
with inducing parameter the inverse square root of the
Borel modulus, up to Weyl conjugacy. On the diagonal the
half-sum of the roots over \(E_w\) has entries
\((m-1)/2,(m-3)/2,\ldots,(1-m)/2\). Equivalently, on a
split diagonal coordinate of \(U_m\) the two conjugate
coordinates are inverses and the two opposite entries of
this list contribute their difference to the relative
modulus exponent; \(q_w=q_v^2\). This gives the displayed
Satake multiset after base change. The unramified space
character on the common uniformizer of \(F_v,E_w\) equals
\(\eta_v(\varpi_v)^m=1\). Hence the standard factor at
\(s+1/2\) has the \(m\) arguments \(s+j-d\) in the lemma.
At infinity the entire embedded definite group belongs to
the doubled maximal compact group. The flat-section height is
one there and the prescribed Gaussian is fixed by the compact
pair, by Theorem~\ref{thm:globalequalrank826}. Its squared
norm integrated over this compact group is positive.
\end{proof}

\begin{theorem}[The normalized first exchange on every prescribed twist]
\leavevmode\label{thm:normalfirst828}
On each of the given fields there is a positive number \(B_{F,m}\),
independent of \(\delta\), such that on the entire adelic source
\[
             |I_\delta(g)|^2=B_{F,m}D_\delta .
\]
In particular \(\gamma_\delta=I_\delta(1)^{-1}\) is defined,
\(\gamma_\delta I_\delta(g)=1\) for every \(g\), and
\(|\gamma_\delta|^2=B_{F,m}^{-1}D_\delta^{-1}\).
\end{theorem}
\begin{proof}
Use the same Haar choices on the two copies of \(G\).
The integral \(J_\delta(h)=\int_{[G]}\theta_\delta(g,h)\,dg\)
is constant by the proof of Theorem~\ref{thm:globalequalrank826}
with source and target interchanged. Indeed the local
double-coset equality of Lemma~\ref{lem:splitaverage826}
is reversible since \(a\mapsto a^{-\mathsf t}\) is an
involution. Full integral and real invariance hold on both
sides. The same determinant and strong approximation argument
therefore applies. Integrating both ways on the compact
product shows \(J_\delta(h)=I_\delta(1)\). Thus the squared
norm of \(J_\delta\) is \(\operatorname{vol}([G])|I_\delta(1)|^2\).

Double the source as \(I_m\oplus(-I_m)\), with its diagonal
maximal isotropic subspace. In the doubled construction use
the Schwartz tensor of the test and its conjugate. On the
two factors the rational theta sum of this tensor restricts
to \(\theta_\delta(g_1,h)\overline{\theta_\delta(g_2,h)}\).
This follows from the product polarization and from
complex conjugation of the oscillator on the opposite
copy, as in Proposition~\ref{prop:doubledkernel711}. The
external target characters cancel, and the external source
characters give \(\delta(\det g_1)^{-1}\delta(\det g_2)\).
Equivalently one applies these characters to the two source
arguments after doubling with the underlying untwisted
splitting. Keep the usual correction on the second source
argument in the restriction of the doubled splitting. It is
the character making the two actions there contragredient:
on the diagonal the graph evaluation is the Hermitian inner
product and is invariant under the simultaneous unitary action.
With this convention the value at \((a,1)\) is precisely the
oscillator coefficient used in Section~\ref{sec:squarefactors827}.
The diagonal modulus stays one throughout the flat-section
family since a unitary determinant has norm one.

Let \(\mathcal Z_\delta(s)\) be the integral on \([G]^2\)
of the flat Siegel Eisenstein series with that doubled test
and these source characters. Here the Siegel rank \(n\) is
\(m\), the dimension of the space integrated in the theta
kernel is also \(m\), and its rational Witt index is zero.
The unitary Siegel--Weil formula at \(m=n\) says that the
Eisenstein series is holomorphic at \(s=0\) and its value
is twice the theta integral in the measures of that formula.
There is no regularization of the theta integral over our
compact quotient. These are exactly the \(m=n\) hypotheses of
\href{https://msp.org/pjm/2013/264-1/pjm-v264-n1-p11-s.pdf}
{S.~Yamana, \emph{The Siegel--Weil formula for unitary
groups}, Theorem~2}. Changing the fixed Haar measures gives
a positive common factor. Integration on the compact source
quotients commutes with evaluation at zero: the meromorphic
family of smooth Eisenstein series, holomorphic here, is
locally uniform on compacta. Expanding the product of the
two ordinary theta integrals then shows that
\(\mathcal Z_\delta(0)=c\,|I_\delta(1)|^2\) with \(c>0\)
depending only on the field, spaces, and measures.

For \(\operatorname{Re}s\) sufficiently large unfold the
absolutely convergent Eisenstein series. Its terms are indexed
by the rational maximal isotropic subspaces of the double.
Each is the graph of a rational isometry of \(I_m\): its
intersection with either summand is zero by definiteness,
so both projections are isomorphisms, and isotropy is exactly
the isometry condition. There is one orbit under \(G(F)^2\)
and its stabilizer is the diagonal \(G(F)\). Unfolding thus
integrates the section on the open orbit, with one factor
\(\operatorname{vol}([G])\) from simultaneous translation.
The other variable runs over \(G(\mathbb A)\); on that
variable the section factors into the local doubling
integrals. Absolute convergence permits the product in
this half-plane. This is also the usual Eulerian unfolding
in \href{https://arxiv.org/pdf/1207.4709}{Gan, Qiu and
Takeda, \S\S11.5--11.6}. The section parameter has the
split height computed in Section~\ref{sec:squarefactors827}.

Combining Theorem~\ref{thm:splitzeta827} at all split places
with Lemma~\ref{lem:squarespherical828} at the other places
gives the meromorphic identity
\[
 \mathcal Z_\delta(s)=c'\,
 \frac{\displaystyle\prod_{j=1}^m L_{E,f}(s+j-d,\lambda_\delta)}
      {\displaystyle\prod_{j=1}^m L_{F,f}(2s+j,\eta^j)}
       \prod_{v\in\mathcal S_\delta}\frac{P_m(q_v)}{H_m(q_v,s)} .
\]
Here \(c'>0\) includes the real compact integrals and measure
normalizations. They are common to the twists. On an inert
uniformizer the base-change character has value one as used
in the lemma. The identity holds initially in the unfolding
half-plane and extends by the Hecke functional equations
and uniqueness of meromorphic continuation.

All factors on the right have finite nonzero values at zero.
For the denominator, the \(j=1\) value is \(L_{F,f}(1,\eta)\),
which is positive by taking the ratio of residues in
\(\zeta_E(z)=\zeta_F(z)L_{F,f}(z,\eta)\). The other denominator
arguments exceed one. The correction at each mark is one
at zero. In the numerator, Theorem~\ref{thm:squareconductor827}
gives nonvanishing at the positive integers \(1,\ldots,d\).
At each complex place of \(E\) the angular character of
\(\lambda_\delta\) has exponent \(2d\) or \(-2d\), since
\(\upsilon_1\) has weight one on the norm-one circle and
base change is \(a\mapsto a/\bar a\). Its gamma factor is
\(\Gamma_{\mathbb C}(z+d)\). At \(z=1-d,\ldots,0\) these
gamma factors are finite and nonzero. The primitive Hecke
functional equation now gives finite nonzero numerator
values at these arguments from the paired positive ones.
It follows that \(I_\delta(1)\ne0\).

Take the ratio of the meromorphic identities for any two
twists. All nonsplit and real factors cancel. The ratio at
zero is the ratio of the squared absolute values of the
two actual constants, hence is positive. Its absolute
value is \(D_{\delta_1}/D_{\delta_0}\) by
Theorem~\ref{thm:squareconductor827}. Fixing any prescribed
reference twist gives \(B_{F,m}=|I_{\delta_0}(1)|^2/D_{\delta_0}>0\).
This proves the assertions for the entire family on that field.
\end{proof}

\begin{corollary}[The signed constant for insertion in the second exchange]
\leavevmode\label{cor:signedfirst828}
Let \(c_0=\chi_M^j\vartheta\), take \(\delta=c_0^{-1}\),
and write \(\chi\) for the finite residue character of \(c_0\).
Denote by \(I^-(g)\) the definite target class sum with source
\(-I_m\), positive target \(I_m\), splitting characters
\((c_0\upsilon_1^m,c_0\upsilon_1^m)\), marked test
\(\chi(\det X)\) and real polynomial one. There is a scalar
\(\gamma^-\) such that
\[
       \gamma^- I^-(g)=1\quad\text{for all }g,\qquad
       |\gamma^-|^2=B_{F,m}^{-1}/D_{c_0^{-1}}.
\]
\end{corollary}
\begin{proof}
Use the same finite unit \(u_f\) of norm \(-1\) as in
Lemma~\ref{lem:signkernel823}. Its antiunitary identification
and the uniqueness of the rational summation functional give
a single phase \(\epsilon\) for the pair of ranks \(m,m\).
The finite calculation in that lemma applies with both group
arguments moving. In fact the Hom change \(L\mapsto hLg^{-1}\)
has determinant \((\det h)^m(\det g)^{-m}\), so the ratio of
the conjugated and reversed lattice phases is
\(\upsilon_1(\det h)^m\upsilon_1(\det g)^{-m}\).
Together with complex conjugation of both external characters
these are the two specified negative-pair characters.
At the real places the compact groups fix both tests: on the
reversed pair the \(m\) negative Hom columns give inverse
target and direct source vacuum exponents, canceled by the
respective external weights \(m\). The old and new tests
transform by
\(\overline{\chi^{-1}(\det(u_f^{-1}X))}
=\chi(u_f)^{-m}\chi(\det X)\) on units, with both sides
zero off the units. Here \(\chi(u_f)\) is evaluated in the
marked residue components. Consequently
\(\overline{I_\delta(g)}=\epsilon\chi(u_f)^{-m}I^-(g)\)
for every \(g\), since the class weights are real.
Multiplying by \(\overline{\gamma_\delta}\) gives the assertion
with \(\gamma^-=\overline{\gamma_\delta}\epsilon\chi(u_f)^{-m}\).
The two additional factors have modulus one.
\end{proof}

This gives the signed constant and its full twist dependence
at the first step. The second exchange passes from that
constant to the carrier on \(U(2m,m)\). The third compares
the carrier period with the arithmetic class integral on
\(U(W_j)\).
\clearpage

\subsection{The two cuspidal forms in the second exchange}
\label{sec:secondcarrier829}
Put \(r=3m\), \(S_0=I_{2m}\oplus(-I_m)\),
\(G=U(S_0)\), \(H=U(I_r)\), and \(\kappa=\upsilon_1^m\).
Fix a prescribed \(c=\chi_M^j\vartheta\), with finite
residue character \(\chi\). Let \(C(h)\) be the theta lift
of the constant on \(U(I_m)\) to \(H\) with the two external
characters one, the full lattice vacua, and polynomial one.
Normalize compact measures so that integration on the
definite groups gives the previous stabilizer-weighted
class sums. Both compact real group actions fix this definite
Gaussian. The coefficient \(C_\vartheta(h)\) in
Theorem~\ref{thm:arithmeticclass824} is
\(c(\det h)C(h)\): its finite and real test vectors are the
same and the external character \(c\) acts on the target
determinant alone.

Let \(\Theta_c(x,h)\) be the kernel for \((S_0,I_r)\) with
characters \((c\kappa,c)\), polynomial one, the marked square
test \(\chi(\det X)\) on integral matrices, and unmarked unit
vacua. The unnormalized form in the second exchange is
\[
                   A_c(x)=\int_{[H]} C(h)\Theta_c(x,h)\,dh .
\]
Indeed moving the factor \(c(\det h)\) from the coefficient
into the kernel changes the first splitting character from
\(\kappa\) to \(c\kappa\). Let \(\phi\) be the prescribed
nonzero carrier of Theorem~\ref{thm:carrier813}, a theta
integral from \(-I_m\) to \(S_0\). We compare \(A_c\) and
\(\phi\) on this same group.

\begin{proposition}[Level, weight and cuspidality of the class integral]
\leavevmode\label{prop:secondcusp829}
The function \(A_c\) is a smooth cuspidal automorphic form
on \(G\), fixed by its full finite integral stabilizer. It
is holomorphic with compact weight
\((0,\ldots,0;\,r,\ldots,r)\), in the order of
Theorem~\ref{thm:carrier813}.
\end{proposition}
\begin{proof}
At a marked place the source change by an integral unit \(k\)
multiplies the raw determinant test by \(\chi(\det k)\);
the splitting multiplier is its inverse. The corresponding
target change acquires the inverse factor from the test and
the direct factor from the splitting since \(\kappa\) is
unramified. Unmarked finite vacua are fixed. At a real place
each negative source column contributes one inverse target
determinant on the vacuum, so the target compact action
has weight \(-m+m=0\). For the source action, each target row
contributes one copy of the determinant on the negative
source block. The positive block has weight zero. This gives
the two claimed source weights, since \(c\) has weight zero.
In the polynomial Fock model the noncompact lowering
operators are contractions, each summing a product of
derivatives in a positive and a negative column over the
target row index. They annihilate the polynomial one.
These are the holomorphy equations for the positive-first
order, as in Theorem~\ref{thm:carrier813}. Rational
invariance and smoothness follow from the theta construction
and compactness of \([H]\).

Let \(P\) be a rational maximal parabolic of \(G\), attached
to a nonzero totally isotropic \(E\)-space \(U\subset S_0\).
We show that even the average over the center of its
unipotent radical vanishes. Complete \(U\) to a rational
Witt decomposition \(U\oplus S_1\oplus U'\), with \(U'\)
dual to \(U\). In the mixed position model the oscillator has
position coordinates \(y\in\operatorname{Hom}_E(U,I_r)\)
and a remaining oscillator on \(S_1\) and \(I_r\).
The rational theta sum accordingly sums these coordinates
over their rational values. On the term of coordinate \(y\),
a central shear with parameter \(B\) multiplies by the
additive character paired with the Gram matrix \(y^*y\).
Indeed the shear sends a dual isotropic coordinate to itself
plus \(B\) times an isotropic coordinate. In position
coordinates a shear of a momentum by a linear function of
position multiplies by the additive character of its
quadratic form. Here that form is a fixed nonzero trace
multiple of \(\operatorname{tr}(B y^*y)\); the imaginary
unit of the unitary splitting identifies the skew and
Hermitian conventions. The shear fixes the remaining
oscillator. The trace pairing between the two rational
matrix spaces is nondegenerate, as one sees on diagonal
entries and on each pair of off-diagonal entries. Integration
over the adelic center modulo rational points therefore
requires \(y^*y=0\), by additive character orthogonality.
The rational theta sums may be integrated termwise on this
compact unipotent quotient by uniform Schwartz summability
on compact sets. Definiteness of \(I_r\) at a real place
forces every column of such a rational \(y\) to vanish.

The component at \(y=0\) is zero already at any one marked
split place. In the split position coordinates before passing
to the mixed model, the test is supported on invertible
square matrices \(X\). This persists after arbitrary local
source and target translations, which compose \(X\) with
invertible maps. Relative to the Witt decomposition, passing
to the mixed model keeps the columns \(X|_U\) on the chosen
split component and Fourier transforms the columns on \(U'\);
the dual Fourier variables supply the conjugate components
of \(y\) using the Hermitian pairing. In particular setting
\(y=0\) sets \(X|_U=0\) before integration in the latter
columns. An invertible \(X\) cannot vanish on a nonzero
subspace. The other coordinates \(S_1\) are unrestricted,
so the entire mixed vector at zero is zero. This proves
vanishing of the central average for every \(x,h\), hence
of the full parabolic constant term, also after integration
against \(C(h)\). The old split marks are present for every
twist. Every proper parabolic is contained in a maximal one,
so vanishing of their constant terms proves cuspidality.
\end{proof}

We next compare all the split spherical Hecke operators,
including those at the conductor. At a split finite place
with residue size \(q\), identify the two rank \(r\) groups
with \(\operatorname{GL}_r(k)\) by their prescribed finite
frames. For \(0\le b\le r\) let \(T_b\) denote right
translation averaged with probability on the double coset of
\(\operatorname{diag}(\varpi 1_b,1_{r-b})\) at full integral
level. Write \({N\brack h}_q\) for the number of
\(h\)-dimensional subspaces of a residue vector space of
dimension \(N\), with value zero outside \(0\le h\le N\).
Let \(\xi\) be the value of the unramified Weil splitting
character at a uniformizer in the chosen \(E\)-component.
Define
\[
 \begin{split}
 t_b&=\xi^{-mb}q^{mb/2},\qquad s_l=\xi^{-rl}q^{rl/2},\\
 \lambda_b&=\frac{t_b}{{r\brack b}_q}
  \sum_{\substack{0\le l\le\min(m,b)\\ b-l\le r-m}}
     {r-m\brack b-l}_q\,q^{\,l(r-m-b+l)}
                       {m\brack l}_q\,s_l^{-1}.
 \end{split}
\]

\begin{lemma}[A rectangular lattice count for the lower theta lift]
\leavevmode\label{lem:secondhecke829}
At every split finite place \(T_b C=\lambda_b C\) on \(H\)
and \(T_b\phi=\kappa_v(\varpi)^b\lambda_b\phi\) on \(G\).
\end{lemma}
\begin{proof}
Consider first the \((r,m)\)-matrix test
\(f={\bf1}_{\operatorname{Mat}_{r,m}(O)}\) with the two
external characters one. A target element with \(b\)
elementary divisors \(\varpi\) acts by
\(t_b f(h^{-1}X)\). A source element with \(l\) divisors
\(\varpi^{-1}\) acts by \(s_l f(Xg)\), with the other
elementary divisors one in both statements. The positive
absolute factors are the square-root position Jacobians.
For the phases one may compute the overlap as in
Theorem~\ref{thm:splitzeta827}: a target change of valuation
one gives \(m\) positive Hom exponents on the chosen component,
hence \(m\) negative ones on the conjugate component. A
source change of valuation minus one gives \(r\) positive
Hom exponents on the chosen component. The conjugate
uniformizer has splitting value \(\xi^{-1}\), giving the
two phases written above.

Both averaged tests vanish off the integral matrix lattice.
For an integral \(X\), let \(a\) be its residue rank. The
target average divided by \(t_b\) is the fraction of
codimension \(b\) residue subspaces containing its column
space. It is \({r-a\brack b}_q/{r\brack b}_q\).
The corresponding source average divided by \(s_l\) is
the fraction of \(l\)-dimensional subspaces contained in
\(\ker\bar X\), namely \({m-a\brack l}_q/{m\brack l}_q\).
Indeed \(gO^m\) is \(O^m\) enlarged by \(\varpi^{-1}\)
times the chosen subspace. The two subspaces in these
counts are uniform because the integral compact groups
act transitively on the relevant Grassmannians.

Put \(t=r-m\) and \(n=m-a\). In a space \(k_0^t\oplus k_0^n\)
over the residue field \(k_0\), count a \(b\)-plane by the
dimension \(l\) of its projection to the second summand.
Choose its \(b-l\) dimensional intersection with the first
summand and its \(l\) dimensional image in the second.
For these choices the possible graphs are indexed by the
linear maps from the image to the first summand modulo the
intersection, giving \(q^{l(t-b+l)}\) possibilities. Thus
\[
 {r-a\brack b}_q=
   \sum_l {m-a\brack l}_q {r-m\brack b-l}_q q^{l(r-m-b+l)} .
\]
The support and this identity express the target average of
\(f\) as the corresponding linear combination of the source
averages, with coefficients whose sum is \(\lambda_b\).
Integrate over the definite source against the constant
function. Each source translate has exactly the original
integral by Haar invariance. This gives the identity for
\(C\). For \(\phi\) the finite rectangular test and the two
finite frames are identical: the negative definite source
is framed using the norm-minus-one unit. The different real
polynomial does not enter this local identity. The additional
target multiplier is \(\kappa(\det h)\), constant and equal
to \(\kappa_v(\varpi)^b\) on this double coset. The source
integral is again against the constant. This proves the
second identity.
\end{proof}

\begin{theorem}[Agreement of the two split systems of Hecke equations]
\leavevmode\label{thm:secondcomparison829}
On the prescribed field and for every \(c\), both \(A_c\)
and the carrier \(\phi\) satisfy
\[
                    T_b f=\kappa_v(\varpi)^b\lambda_b f
\]
at every split finite place and every \(0\le b\le r\).
This includes both sets of marked places and determines the
action of the full split spherical Hecke algebra on these forms.
\end{theorem}
\begin{proof}
Use the square test defining \(\Theta_c\). At a split place
it is fixed by both compact actions and is invariant under
transposition. For the underlying Hom action transposition
exchanges the source operator at \(a\) with the target
operator at \(a^{-\mathsf t}\), by the pullback argument of
Lemma~\ref{lem:splitaverage826}. The quotient of the external
source and target factors now equals \(\kappa(\det a)\),
since they are \(c(\det a)^{-1}\) and
\((c\kappa)(\det a)^{-1}\), respectively. The compact
restrictions of \(c\) and \(c\kappa\) agree. As in that
lemma, any vector fixed for both compact actions is fixed
by transposition, by Smith normal form. Hence the source
average for \(T_b\) equals \(\kappa_v(\varpi)^b\) times
the target average on the inverse-transpose coset. This
holds equally for the determinant test at a mark and for
the lattice indicator elsewhere.

Pass this equality through the theta sum and integrate over
\([H]\) with \(C\). Right translation of the target kernel by
\(a^{-\mathsf t}\) transfers to right translation of \(C\)
by \(a^{\mathsf t}\), using the right invariance of the quotient
measure. The transpose of a double coset has the same
elementary divisors and the induced measure is its invariant
probability. Thus the resulting right average on \(C\) is
\(T_b C\). Lemma~\ref{lem:secondhecke829} now gives the stated
equation for \(A_c\), and gives the same equation for \(\phi\)
directly. The spherical Hecke algebra of
\(\operatorname{GL}_r(k)\) is generated by these elementary
double cosets and the inverse of the central coset \(b=r\)
(the elementary symmetric generators under Satake). The latter
coset acts by an invertible translation, so its inverse action
also agrees.
\end{proof}

The calculation locates \(A_c\) in the same cuspidal
holomorphic weight space and split Hecke eigenspace as the
nonzero carrier. The second exchange still requires showing
that the two specified global vectors are proportional with
the prescribed nonzero scalar and computing its norm.
\clearpage

\subsection{The spherical parameters on the two groups of rank \(3m\)}
\label{sec:carrierspectra830}
Let \(G=U(I_{2m}\oplus(-I_m))\) and \(H=U(I_r)\), where
\(r=3m\). Use the coefficient \(C\), the carrier \(\phi\),
the kernel \(\Theta_c\) and the form \(A_c\) of
Section~\ref{sec:secondcarrier829}. For a full-level function
\(u\) on the compact quotient \([H]\), invariant at infinity,
write
\[
 (\mathcal T_c u)(x)=\int_{[H]}u(h)\Theta_c(x,h)\,dh,\qquad
 B_c(h)=\int_{[G]}\phi(x)\overline{\Theta_c(x,h)}\,dx .
\]
Thus \(\mathcal T_c C=A_c\). We use inner products linear
in the first argument. Then \(B_c=\mathcal T_c^*\phi\) in
the sense of these pairings. All the integrals exist:
a cusp form is rapidly decreasing on Siegel sets and the
theta kernel has moderate growth, uniformly for \(h\) in a
compact set of representatives. The proof of
Proposition~\ref{prop:secondcusp829} gives cuspidality and
the same holomorphic compact weight for \(\mathcal T_c u\).
It only uses the compactness of \([H]\) and the marked support
of the kernel. The function \(B_c\) is full-level and fixed
by the compact real group of \(H\).

For a positive integer \(b\) write
\[
             A_b(q)=\{q^{(b-1)/2-i}:0\le i<b\}
\]
as a multiset. A spherical parameter at a split place is
the Satake multiset for \(\operatorname{GL}_r(k)\) in the
chosen component, with the elementary positive coweights
acting by the normalization of Lemma~\ref{lem:secondhecke829}.
At an inert place use the standard base-change multiset
of Theorem~\ref{thm:inert175}.

\begin{lemma}[The definite lower lift is nonzero]
\leavevmode\label{lem:nonzerodefinite830}
At every sufficiently large node of the prescribed family one
has \(C\ne0\).
\end{lemma}
\begin{proof}
Double the definite source \(I_m\) of the lift defining \(C\).
The target \(I_{3m}\) is anisotropic, the Siegel rank is \(m\),
and the Siegel point is \(m>m/2\). Thus the convergent
Siegel--Weil formula and the unfolding of
Theorem~\ref{thm:carrier813} apply. Every rational maximal
isotropic subspace of \(I_m\oplus(-I_m)\) is a graph of an
isometry. Consequently \(\|C\|^2\) is a positive Haar factor
times the integral of the source oscillator coefficient for
the prescribed vacuum. The real source fixes the Gaussian
and gives a positive compact volume. Each finite coefficient
is one on the full stabilizer, of mass one. Off that
stabilizer its integral in absolute value is \(O(q_v^{-m/2})\)
by Lemma~\ref{lem:selfdualtail813}. This lemma uses only the
finite self-dual lattices and the two ranks \(m,3m\), which
are the same for this positive pair. The sum of the errors
is \(O(f2^{-m/2})=o(1)\) on the fields under consideration,
as in Theorem~\ref{thm:carrier813}. Absolute convergence
bounds the difference of the finite product from one by
\(\exp(O(f2^{-m/2}))-1\). Thus it is nonzero at sufficiently
large nodes, proving the assertion.
\end{proof}

\begin{theorem}[Full finite spherical systems of the carriers and adjoint]
\leavevmode\label{thm:allcarriers830}
At every split \(v\), put \(q=q_v\),
\(\xi=\upsilon_w(\varpi_v)\) in the chosen component and
\(\kappa_v(\varpi_v)=\xi^{2m}\). Define
\[
 \begin{split}
 P_{H,v}&=\xi^{\,r-m}A_m(q)\ \sqcup\ \xi^{-m}A_{r-m}(q),\\
 P_{G,v}&=\kappa_v(\varpi_v)P_{H,v}.
 \end{split}
\]
At an inert place \(v\) with \(w\mid v\), set
\[
          P_{H,v}=P_{G,v}=A_m(q_w)\sqcup A_{r-m}(q_w),
                       \qquad q_w=q_v^2 .
\]
Every constituent of \(C,B_c\) at full level has parameter
\(P_{H,v}\), and every constituent of \(\phi,A_c\) at full
level has parameter \(P_{G,v}\). Here a statement about
constituents allows a form to vanish.
\end{theorem}
\begin{proof}
At a split place the probability double coset \(T_b\) of
Section~\ref{sec:secondcarrier829} has eigenvalue
\(q^{b(r-b)/2}e_b(P)/{r\brack b}_q\) at multiset \(P\).
This is the minuscule coweight formula in the Satake
isomorphism: its unnormalized transform is the sum of the
weight monomials with \(b\) distinct indices times
\(q^{b(r-b)/2}\), and its double coset consists of
\({r\brack b}_q\) right cosets. The elementary symmetric
function of the segment of length \(t\) satisfies
\[
            e_l(A_t(q))=q^{-l(t-l)/2}{t\brack l}_q .
\]
For instance, expand \(\prod_{i=0}^{t-1}(1+zq^i)\) by degree,
or induct using the last factor. In the union \(P_{H,v}\)
choose \(l\) entries in the first segment and \(b-l\) in
the second. Their phase is \(\xi^{-mb+rl}\). After multiplying
their symmetric function by \(q^{b(r-b)/2}\), the exponent
of \(q\) in this term is
\[
 \frac{b(r-b)-(b-l)(r-m-b+l)-l(m-l)}2
             =\frac{mb-rl}{2}+l(r-m-b+l).
\]
The resulting coefficients are exactly those in \(\lambda_b\)
of Lemma~\ref{lem:secondhecke829}. Multiplication of every
entry by \(\kappa_v(\varpi_v)\) accounts for the factor
\(\kappa_v(\varpi_v)^b\) in the equations for \(\phi,A_c\).

For \(B_c\), the transpose calculation of
Theorem~\ref{thm:secondcomparison829} before evaluating on
\(C\) gives, on the whole full-level compact source space,
\[
 T_b^G\mathcal T_c u
    =\kappa_v(\varpi_v)^b\,\mathcal T_c T_b^H u .
\]
Take adjoints and apply this to the eigenfunction \(\phi\).
Spherical convolution is a commutative algebra closed under
adjoints in the unitary right regular representation. Thus
each operator is normal, and the eigenvalues of its adjoint
on an eigenvector are the complex conjugates. It follows
that \(T_b^H B_c=\lambda_b B_c\). This argument includes
every marked split place.

We detail the passage from global theta integrals to the
local condition used at an inert place. The spaces of
cuspidal automorphic forms with the given finite level and
real holomorphic weight are finite dimensional, by finite
cuspidal multiplicity at fixed level and infinitesimal
character. On the compact group \(H\) the corresponding
space is a finite class space. Decompose such forms in the
cuspidal Hilbert sum. If a projection of a theta lift pairs
nontrivially with an irreducible summand, there is a
nonzero joint local oscillator pairing with its two
constituents. Indeed fix pure tensors giving a nonzero
global pairing and hold the data outside \(v\) fixed.
The local form on the oscillator and the two test
representations is equivariant, by right translation in the
two quotient integrals. Its associated map from the smooth
oscillator lands in the smooth dual of the tensor of the
two local representations. These representations are
admissible, so the smooth dual is the tensor of their
contragredients; this follows on compact-fixed vectors by
finite dimensionality. The nonzero image is an irreducible
quotient isomorphic to that tensor. One may first project
the input onto one of finitely many summands with nonzero
contribution. Complex conjugation in an adjoint integral
replaces the appropriate representations by their duals.
At an inert place an unramified unitary-group parameter is
invariant under conjugate inversion at the common uniformizer:
its entries occur in reciprocal pairs, with a possible
central entry one. Thus duals give the same multiset here.

All dyadic places of \(F\) split in \(E\). In fact the
imaginary quadratic subfield is split at two, and conjugation
on \(E/F\) exchanges those two components. It cannot fix any
place above two. Every remaining nonsplit finite place is
unramified quadratic by Lemma~\ref{lem:signkernel823}. The
even unitary spaces in question are quasi-split there, with
hyperspecial stabilizers, by the hyperbolic lattice splitting
in Lemma~\ref{lem:squarespherical828}. Thus
Theorem~\ref{thm:inert175} applies at every inert place.
Both source and target space characters on the uniformizer
are one since the ranks are even. The finite restrictions
of \(c,\kappa\) to the compact norm-one group there are
trivial. Starting with the constant in rank \(m\), whose
multiset is \(A_m(q_w)\), the positive lower lift to \(H\)
and the negative lower lift to \(G\) both adjoin \(A_{r-m}(q_w)\).
This gives the systems for \(C\) and \(\phi\). For the square
pair of ranks \(r,r\) the same transfer theorem leaves the
multiset unchanged. Apply the local pairing just proved to
\(\mathcal T_c C\) and to its adjoint on \(\phi\) to obtain the
two remaining assertions.
\end{proof}

At all finite places a spherical isomorphism class is determined
by its Satake multiset here. At an inert place the relative torus
parameter is determined up to signed permutation by its reciprocal
pairs, which is the relative Weyl equivalence. Thus the theorem
also fixes the finite spherical matrix coefficients with the
unit-vector normalization, including at places above two.
\clearpage

\subsection{Uniqueness of the two automorphic lines}
\label{sec:carrierlines831}
Fix a sufficiently large node so that \(C\) and \(\phi\) are
both nonzero, as in Lemma~\ref{lem:nonzerodefinite830} and
Theorem~\ref{thm:carrier813}. Let \(D\) stand for either
\(H=U(I_{3m})\) or \(G=U(I_{2m}\oplus(-I_m))\).
Write \(\mathcal E_D\) for the space of smooth cuspidal
automorphic forms on \([D]\) with full finite level and
Hecke eigencharacter specified by \(P_{D,v}\) in
Theorem~\ref{thm:allcarriers830} at every finite \(v\).
On \(H\) require invariance at infinity. On \(G\) require the
scalar compact weight \((0^{2m};(3m)^m)\) and holomorphy at
every real place. The compact group \(H\) has no proper
rational parabolic, so its cuspidality condition is vacuous.

\begin{lemma}[The matrix coefficient on a prescribed scalar line]
\leavevmode\label{lem:carriermatrix831}
For every nonzero \(f\in\mathcal E_D\) the function
\[
 M_f(a)=\frac{\langle R(a)f,f\rangle}{\langle f,f\rangle},
 \qquad a\in D(\mathbb A),
\]
is the same. It is a product of normalized local coefficients.
At a real place the coefficient is determined by the compact
character and the holomorphy equations. This real assertion
also holds for a smooth unit vector satisfying the same
equations in any unitary representation of the real group.
\end{lemma}
\begin{proof}
At a finite place insert the orthogonal projections onto the
\(K_v\)-fixed vectors on both sides of \(R(a_v)\) in the
pairing. This compressed operator is convolution by probability
on \(K_va_vK_v\). Its eigenvalue is determined by the specified
spherical character. It commutes with translation at all other
places. Removing the finite components one by one multiplies
the archimedean coefficient by these eigenvalues. Only finitely
many factors differ from one for any fixed adele.

At a real place of \(H\) the coefficient is one. For \(G\)
use the complexified Lie algebra
\(\mathfrak p^+\oplus\mathfrak k_{\mathbb C}\oplus\mathfrak p^-\),
where the antiholomorphic derivatives \(\mathfrak p^-\)
annihilate the vector. Complex conjugation of the real form
interchanges \(\mathfrak p^+\) and \(\mathfrak p^-\).
In a unitary representation the adjoint of a complexified
derivative \(Z\) is \(-\bar Z\) on smooth vectors. Order
monomials by the Poincar\'e--Birkhoff--Witt theorem with
\(\mathfrak p^-\) on the right and \(\mathfrak p^+\) on the
left. In a derivative of the coefficient at the identity
any ordered monomial containing a right lowering factor
annihilates the vector. Any remaining monomial with a left
raising factor pairs to zero by the adjoint relation.
Every term left over is evaluated by the differential of
the one-dimensional compact character. Thus all derivatives
of the normalized coefficient at the identity are determined.

Here these derivatives determine the function globally.
Indeed choose real tangent directions in the noncompact
part in orthonormal pairs \(X,JX\) for its invariant complex
structure. For \(Z=X+iJX\),
\(X^2+(JX)^2=(Z\bar Z+\bar Z Z)/2\).
One of \(Z,\bar Z\) is lowering. On the given vector the sum
therefore reduces, up to sign and the factor \(1/2\), to the
commutator in \(\mathfrak k_{\mathbb C}\). Each compact
direction also acts by a scalar. The sum of the squares
of this real basis, including a basis of the compact part,
consequently acts on the vector by a fixed scalar. The
coefficient solves the corresponding left-invariant elliptic
differential equation with real analytic coefficients.
Analytic elliptic regularity makes it real analytic.
The real unitary group is connected by its Cartan
decomposition and the connectedness of its maximal compact
subgroup. Equality of all derivatives implies equality everywhere.
On the product of real groups use the union of the bases.
The derivatives at the identity factor by place on pure
monomials, and the product of the local normalized coefficients
therefore gives the same analytic function. For example one
may use at each place the restriction to that factor of the
coefficient of any one of the specified unit vectors.
\end{proof}

For \(D=H\) let \(L=U(I_m)\) and let \(K(l,z)\) be the
definite rectangular theta kernel defining \(C\).
For \(D=G\) take \(L=U(-I_m)\) and use the rectangular
kernel with polynomial \(\prod_\tau\det(Z_{\tau,-})^{2m}\)
defining \(\phi\). In both cases put
\[
 T u(z)=\int_{[L]}u(l)K(l,z)\,dl,\qquad
 T^* f(l)=\int_{[D]}f(z)\overline{K(l,z)}\,dz .
\]
Write \(f_0=T1\). The integrals against cusp forms converge
absolutely by rapid decrease and the moderate growth of the
theta kernel. Bounds are uniform over the compact quotient
\([L]\). These facts give the adjoint pairing even when the
form \(Tu\) is only known to have moderate growth.

\begin{theorem}[Uniqueness on the two groups of rank \(3m\)]
\leavevmode\label{thm:carrierline831}
One has \(\mathcal E_H=\mathbb C C\) and
\(\mathcal E_G=\mathbb C\phi\).
\end{theorem}
\begin{proof}
We first prove that \(T^*\) is injective on \(\mathcal E_D\).
Let \(V\) be the space defining \(D\), of dimension \(r=3m\),
and double it as \(V\oplus(-V)\) with diagonal maximal
isotropic space. Use the tensor of the specified rectangular
test with its conjugate, in the doubled polarization of
Proposition~\ref{prop:doubledkernel711}. Its theta sum on
the two factors at \(z_1,z_2\) is
\(K(l,z_1)\overline{K(l,z_2)}\), with the usual restriction
character on the second factor and the external determinant
characters on the two arguments. Include these characters in
the restriction throughout the flat Siegel family, as in
Theorem~\ref{thm:normalfirst828}. They are fixed independently
of the choice of \(f\). Simultaneous left translation of the
two arguments by \(D(\mathbb A)\) leaves the corrected graph
section invariant. On its diagonal Levi the determinant has
norm one, and the unitary restriction character cancels
there, equivalently by invariance of the Hermitian pairing.

Integrate the resulting Siegel Eisenstein series against
\(\overline{f(z_1)}f(z_2)\) on \([D]^2\) and denote the
meromorphic pairing by \(Z_f(s)\). The integrated theta
space is definite of dimension \(m\), and the Siegel rank
is \(r\). In the unitary normalization the point is
\(s_0=(m-r)/2=-m\).
\href{https://msp.org/pjm/2013/264-1/pjm-v264-n1-p11-s.pdf}
{Yamana, Theorem~2} applies since \(m\le r\).
The ordinary theta integral suffices because \(L\) is
anisotropic. The Schwartz section is holomorphic in the
Eisenstein series at \(s_0\), and the formula gives twice
this theta integral in the measures of that theorem. Hence
\[
                  Z_f(s_0)=b_L\|T^*f\|^2,\qquad b_L>0 .
\]
The same constant works for all \(f\). For a noncompact
\([D]\) we use here the usual meromorphic continuation of
Eisenstein series as smooth automorphic forms of locally
uniform moderate growth after removal of poles. Pairing
with rapidly decreasing cusp forms is continuous on such
a family. More explicitly, on the two Siegel sets its
bound by fixed powers of the heights on compact parameter
sets is absorbed by choosing the two independent orders
of cuspidal decay larger than those powers and the volume
exponents. This permits evaluation inside the integral.

For \(\operatorname{Re}s\) large, unfold the absolutely
convergent series as in the doubling integral of
\href{https://arxiv.org/pdf/1207.4709}{Gan, Qiu and Takeda,
\S11.5}. To see which terms remain, intersect a rational
maximal isotropic subspace \(Y\) with the two summands and
write \(X_1,X_2\) for the intersections. Its projection to
the \(i\)-th summand lies in \(X_i^\perp\). Comparing
dimensions for both projections gives
\(\dim X_1=\dim X_2\) and equality of the projected spaces
with \(X_i^\perp\). The induced subspace modulo
\(X_1\oplus X_2\) is the graph of an isometry between the
two nonsingular quotients \(X_i^\perp/X_i\). If \(X_1\ne0\),
the unipotent radical of the first parabolic stabilizing
\(X_1\) acts trivially on this quotient and preserves \(Y\).
It acts there with determinant one. The section is invariant
under this radical at adelic points in the unfolding, so
integration over its rational quotient gives the vanishing
constant term of \(\bar f\). If \(X_1=0\), both intersections
vanish and \(Y\) is a graph of an isometry of \(V\). These
graphs form one rational orbit with diagonal stabilizer.
Writing \(z_1=za,\ z_2=z\) in the unfolded integral gives
the corrected section at \((a,1)\) paired with
\(\int_{[D]}\overline{f(za)}f(z)\,dz=\|f\|^2\overline{M_f(a)}\).
It follows from Lemma~\ref{lem:carriermatrix831} that
\(Z_f(s)/\|f\|^2=Z_{f_0}(s)/\|f_0\|^2\) first in this
half-plane and then meromorphically. At \(s_0\) it follows
that \(\|T^*f\|^2/\|f\|^2\) has one common value for all
nonzero \(f\in\mathcal E_D\). For \(f_0\) this value is
positive, since \(\langle 1,T^*f_0\rangle=\|f_0\|^2>0\).
This proves injectivity.

It remains to determine the possible image on \([L]\). The
lower compact real group fixes the rectangular vector,
including the determinant polynomial for the negative source
by Theorem~\ref{thm:carrier813}. The finite integral groups
fix it as well. At a split place use the probability Hecke
operators for the positive coweight \((1,0,\ldots,0)\).
The pointwise rectangular count of Lemma~\ref{lem:secondhecke829}
for \(b=1\), before the integration in \(l\), gives
\[
             T_1^D T u=T(h_0u+h_1T_1^L u),\qquad h_1\ne0 .
\]
In fact the two summands in the count have \(l=0,1\);
the \(l=1\) counting coefficient before dividing by the
nonzero normalization of its source average is
\({m\brack1}_q q^{r-m}\) times \(t_1/{r\brack1}_q\).
The possible target character multiplies both terms by the
same unit. The negative source coset in the kernel becomes
the positive coset on \(u\) by the change of variables in
the quotient integral. Applying this equation to \(u=1\)
shows that \(T_1^D\) acts on \(f_0\), and hence on all of
\(\mathcal E_D\), by \(h_0+h_1\). The Hecke operators on
the full-level subspace are normal because the spherical
algebra is commutative and stable under adjoints. Taking
adjoint pairings against arbitrary full-level, real-invariant
\(u\) therefore yields
\((T_1^L)^*T^*f=T^*f\).

In a unitary representation, if a probability average of
translates fixes a vector \(v\), integration of
\(\|R(a)v-v\|^2\) gives zero. By continuity every element
in the support fixes \(v\). Thus the inverse fundamental
double coset fixes \(T^*f\), as do its inverses and
\(\operatorname{GL}_m(O_v)\). These generate
\(\operatorname{GL}_m(F_v)\) by the elementary-divisor
decomposition and permutation of diagonal coordinates.
This holds at every split finite place. The determinant and
strong approximation argument of Theorem~\ref{thm:globalequalrank826}
now gives constancy on \([L]\). Indeed a norm-one determinant
is integral at every inert place. After its diagonal removal,
strong approximation for the simply connected special unitary
group away from one split place and the real places reduces
the remaining adele to integral components. The same argument
works for the form with the opposite sign since it defines
the same unitary group. The injective map \(T^*\) therefore
has image in the one-dimensional space of constants. Each
\(\mathcal E_D\) contains the specified nonzero \(f_0\),
which proves the theorem.
\end{proof}

\begin{corollary}[The scalar and its adjoint in the second exchange]
\leavevmode\label{cor:carriermultiple831}
Write \(A_c=\mathcal T_c C\) and \(B_c=\mathcal T_c^*\phi\)
as in Section~\ref{sec:carrierspectra830}. For every
prescribed twist there are unique scalars \(a_c,b_c\) with
\[
 A_c=a_c\phi,\qquad B_c=b_c C,\qquad
 a_c\|\phi\|^2=\overline{b_c}\|C\|^2,\qquad
 |a_c|^2=\frac{\|B_c\|^2\|C\|^2}{\|\phi\|^4}.
\]
\end{corollary}
\begin{proof}
Their membership in the two spaces follows from
Theorem~\ref{thm:allcarriers830} and the level, weight and
cuspidality in Proposition~\ref{prop:secondcusp829}.
Apply Theorem~\ref{thm:carrierline831}, then use
\(\langle \mathcal T_c C,\phi\rangle
 =\langle C,\mathcal T_c^*\phi\rangle\)
with the inner product linear in the first argument.
\end{proof}
\clearpage

\subsection{Local norm factors for the adjoint of the second exchange}
\label{sec:secondlocal832}
Use \(m=2d,\ r=3m\), the carrier \(\phi\) on \(G\) and
the square kernel \(\Theta_c\) of Section~\ref{sec:secondcarrier829}.
To take the norm on the definite group \(H\), integrate
\(\Theta_c(x,h)\overline{\phi(x)}\) first in \(x\).
The result is \(\overline{B_c(h)}\). In the doubled norm
the graph section at \((a,1)\) therefore multiplies the
normalized coefficient \(\overline{M_\phi(a)}\) on \(G\).
Lemma~\ref{lem:carriermatrix831} gives its factorization by
place. In particular at a split place the spherical parameter
in this coefficient is \(P_{G,v}^{-1}\), obtained by taking
the inverse of each entry. Indeed complex conjugation in a
unitary representation gives the contragredient, whose
Satake character is inverted on the dual torus.

Here is a split local calculation for a general spherical
input. Let \(k\) be a nonarchimedean local field with ring of
integers \(O\), uniformizer \(\varpi\) and residue size \(q\).
Give \(K=\operatorname{GL}_r(O)\) multiplicative measure one
and \(M_r(O)\) additive measure one. For a split torus phase
\(u\) of absolute value one write \(u(a)=u^{v(\det a)}\).
Use \(P_r(q),H_r(q,s)\) from Section~\ref{sec:squarefactors827}
with \(m\) replaced by \(r\).

\begin{lemma}[A split square integral with a spherical coefficient]
\leavevmode\label{lem:splitspherical832}
Let \(M(a)\) be the unit-vector spherical coefficient of
an irreducible unitary spherical representation with parameter
\(\{\alpha_1,\ldots,\alpha_r\}\) in the positive-coweight
convention. If \(\ell(a)\) is the sum of the absolute values
of the elementary divisor valuations, then
\[
 \int_{\operatorname{GL}_r(k)}
 M(a)u(a)q^{-(s+r/2)\ell(a)}\,d^\times a
 =
 \frac{H_r(q,s)}
 {\displaystyle\prod_{i=1}^r
  (1-u\alpha_i q^{-s-1/2})
  (1-u^{-1}\alpha_i^{-1}q^{-s-1/2})}.
\]
The identity holds in a right half-plane and gives the
meromorphic continuation of the integral.
\end{lemma}
\begin{proof}
Put \(t=s+r/2\), let \(K'=M_r(O)\cap\operatorname{GL}_r(k)\)
and \(J(2s)=\int_{M_r(O)}|\det X|^{2s}\,dX\). Diagonalizing
\(a\) as in Theorem~\ref{thm:splitzeta827} gives
\[
 J(2s)q^{-t\ell(a)}
   =|\det a|^t\int_{X,Xa\,\in\, K'}|\det X|^{2s}\,dX,
 \qquad J(2s)=P_r(q)/H_r(q,s).
\]
Singular matrices have additive measure zero. Multiply by
\(u(a)M(a)\), integrate, and set \(Y=Xa\). The change
\(dX=P_r(q)|\det X|^r d^\times X\) leaves the product
weight \(|\det X|^t|\det Y|^t u(X)^{-1}u(Y)\) on \(K'\times K'\),
times \(P_r(q)M(X^{-1}Y)\). The two integrated operators
are bi-\(K\)-invariant. In the matrix coefficient of their
composition one may therefore insert the projection onto
the one-dimensional spherical line between them. This gives
one factor from \(Y\) with \((u,\alpha_i)\) and another from
\(X\) with \((u^{-1},\alpha_i^{-1})\). The inverse parameter
in the second factor follows from \(M(X^{-1})\) being the
coefficient of the dual spherical representation at \(X\).

For clarity, the one-sided factor is
\[
 \int_{K'}|\det y|^t u(y)M(y)\,d^\times y
       =\prod_i(1-u\alpha_i q^{(r-1)/2-t})^{-1}.
\]
One can evaluate each compact double coset on the normalized
spherical principal series with this parameter. Its spherical
vector in the induction model takes the value
\(\delta_B(a)^{1/2}\prod_i\alpha_i^{v(a_i)}\) at \(nak\),
where \(a=\operatorname{diag}(a_i)\), \(B\) is the upper Borel
and \(n\) is upper unitriangular. The spherical Hecke character
on its \(K\)-line is the one of \(M\), also when the principal
series is reducible. Evaluating the convolution on this vector
at the identity computes its eigenvalue. In Iwasawa integration
the measure is \(\delta_B(a)^{-1}dn\,d^\times a\,dk\), with
integral unipotent, unit diagonal and \(K\) of mass one.
Integrality of \(na\) requires \(v(a_i)\ge0\) and
\(n_{ij}a_j\in O\) for all \(i<j\). The volume in \(n\) is
\(\prod_j|a_j|^{-(j-1)}\). Multiplication by
\(\delta_B(a)^{-1/2}\) turns this into
\(\prod_j|a_j|^{-(r-1)/2}\). Independent geometric sums over
the nonnegative valuations give the formula. For
\(\operatorname{Re}t\) sufficiently large all these manipulations
are absolutely convergent. For the change of variables this
also follows from \(|M|\le1\) and the one-sided integral for
the trivial representation. Dividing the two-factor evaluation
by \(J(2s)/P_r(q)\) proves the lemma.
\end{proof}

Define two unitary Hecke characters of \(E\) by
\[
 \lambda_1(a)=(c\upsilon_1^d)(a/\bar a),\qquad
 \lambda_2(a)=(c^{-1}\upsilon_1^m)(a/\bar a).
\]
Let \(\eta\) be the quadratic character of \(F\) associated
with \(E/F\), and \(\mathcal S_c\) the marked split places of
\(F\). At a place \(v\) put
\[
 N_v(s)=
 \prod_{w\mid v}\left\{
 \prod_{j=1}^m L_{E_w}(s+j-d,\lambda_{1,w})
 \prod_{j=1}^{2m} L_{E_w}(s+j-m,\lambda_{2,w})\right\},\qquad
 D_v(s)=\prod_{j=1}^r L_{F_v}(2s+j,\eta_v^j).
\]
The Euler factor of a ramified character is one.

\begin{theorem}[Local factors for the adjoint square norm]
\leavevmode\label{thm:secondlocal832}
For the flat doubled Schwartz section formed from the actual
test of \(\Theta_c\) and its conjugate, the graph integral
against the normalized coefficient of \(\bar\phi\) is
\(N_v(s)/D_v(s)\) at every unmarked finite place. At each
marked \(v\) it is \(P_r(q_v)\). At every real place it
is independent of \(c\), holomorphic in \(\operatorname{Re}s>0\),
and has a finite strictly positive limit at \(s=0\).
\end{theorem}
\begin{proof}
At a split unmarked place take the chosen component of the
norm-one torus and write \(\xi=\upsilon_w(\varpi)\). The
source action has determinant multiplier \(c^{-1}\).
The lattice overlap for a source elementary divisor \(\varpi\)
has phase \(\xi^r\), one power for each target row. This is
the calculation with the Hom lattice in Theorem~\ref{thm:splitzeta827};
it uses only the number of rows and the chosen finite frames.
Consequently at the theta point the unit-vacuum coefficient
is \(u^{v(\det a)}q^{-r\ell(a)/2}\), with
\(u=\xi^r c_v(\varpi)^{-1}\). The target multiplier is
evaluated at the identity. Varying the section multiplies
the graph value by \(q^{-s\ell(a)}\), by the elementary-divisor
calculation of Section~\ref{sec:squarefactors827}. By
Theorem~\ref{thm:allcarriers830}, multiplication of the
inverse spherical parameter by \(u\) gives the two segments
\[
             c_v(\varpi)^{-1}\xi^{-m} A_m(q),\qquad
             c_v(\varpi)^{-1}\xi^{2m} A_{2m}(q).
\]
In Lemma~\ref{lem:splitspherical832} each segment is accompanied
by its inverse. Since \(\upsilon_{1,v}(\varpi)=\xi^2\),
the first pair carries the two values of \(\lambda_1\) above
\(v\), interchanged in order, and the second pair the two
values of \(\lambda_2\). In a segment of length \(2e\) the
arguments of the Euler factors are \(s+j-e,\ 1\le j\le2e\).
At a split place \(D_v(s)^{-1}=H_r(q,s)\).
This proves the unmarked split assertion.

At a marked place the position test is supported on
\(\operatorname{GL}_r(O_v)\), and its absolute value there is
one. A source translate has disjoint support unless \(a\in K_v\).
On \(K_v\) the character in the substitution is cancelled by
the source determinant multiplier. Thus the overlap is its
additive squared norm, \(P_r(q_v)\). Both the flat height and
the spherical coefficient equal one on \(K_v\). This proves
the asserted integral including at an old mark of modulus
\(\varpi^2\).

At every inert place the groups, level, section and constituent
are unramified. The double has the integral diagonal
polarization used in Lemma~\ref{lem:squarespherical828}.
The local norm-one group is compact, and the unramified
space character of even rank has value one on the common
uniformizer. Apply the spherical unitary doubling formula of
\href{https://doi.org/10.1515/crll.1992.428.177}{J.-S.~Li,
Theorem~3.1}, with the normalization and denominator checked
in that lemma. Its numerator is the standard Euler factor at
\(s+1/2\) for the contragredient. Both segments of
\(P_{G,v}=A_m(q_w)\sqcup A_{2m}(q_w)\) are self-inverse.
The characters \(\lambda_1,\lambda_2\) are unramified and
equal one on \(\varpi_w=\varpi_v\). This gives exactly the
indicated \(N_v(s)\), including the residue size \(q_w=q_v^2\).

At a real place put \(p=2m,\ q'=m\) and use Cartan parameters
\(t_1\ge\cdots\ge t_{q'}\ge0\) of \(U(p,q')\). Let \(h(a)\)
be the normalized source coefficient of the Gaussian of
the square oscillator, including its splitting character.
Its compact character and holomorphy equations are those of
\(\phi\) by Proposition~\ref{prop:secondcusp829}. By the real
assertion of Lemma~\ref{lem:carriermatrix831} the normalized
coefficient of \(\bar\phi\) is \(\overline{h(a)}\).
Consequently the integral at the theta point has integrand
\(|h(a)|^2\) times the positive squared norm of the Gaussian.

For a boost the two-mode oscillator factorization in
Lemma~\ref{lem:oscillatorcartan710} gives
\(\lvert h(a(t))\rvert=\prod_i(\cosh t_i)^{-r}\).
In each row, the lowering factor kills the vacuum, the
nonconstant terms from the raising factor are orthogonal to
it, and the diagonal factor contributes \((\cosh t_i)^{-1}\)
for the paired columns. The two compact factors only change
its phase. The Cartan density is bounded by a rank-dependent
constant times \(\exp(2(r-1)\sum_i t_i)\), by the same lemma.
The resulting square is integrable against this density since
\(r=p+q'\), with an exponential bound
\(O_r(\exp(-2\sum_i t_i))\). It is positive in a neighborhood
of the identity.

To see continuity in the section parameter, take orthonormal
coordinates on the real doubled space for its positive
majorant. In the diagonal and orthogonal difference frame
the lower \(r\) rows of the embedded element \((a,1)\) have
Gram determinant \(\det((1+aa^\dagger)/2)\). This follows by
writing them, up to a right unitary change of basis, as
\(2^{-1/2}(a,-1)\). For an upper-block parabolic times
unitary factor, this determinant is \(|\det D|^2\) of the
lower Levi block. The upper block has the reciprocal determinant
modulus. The flat modulus of the unitary Siegel section is
the \(\mathbb C/\mathbb R\) norm of the upper determinant, hence
the height on the graph is \(b(a)=\det((1+aa^\dagger)/2)^{-1}\).
It lies in \((0,1]\), since \((1+y)/2\ge\sqrt y\) for each
positive eigenvalue and \(|\det a|=1\). In the alternative
dual basis the inverse element gives the same determinant.
Thus moving to \(s\) inserts \(b(a)^s\). The integrable square
dominates for \(\operatorname{Re}s\ge0\), giving continuity at
zero and holomorphy in the open half-plane by Morera's theorem
and Fubini. All real data are independent of the weight-zero
twist. This proves the real assertion.
\end{proof}
\clearpage

\subsection{The second global theta exchange and its normalization}
\label{sec:secondglobal833}
For a prescribed twist \(c=\chi_M^j\vartheta\) put
\(Q_c=\prod_{v\in\mathcal S_c}q_v\), with the old marks
included. For a norm-one character \(\alpha\) ramified at
these marks, write
\[
 \mathcal D_e(\alpha)=Q_c^{\,e^2}
    \prod_{z=1}^e|L_{E,f}(z,\lambda_\alpha)|^2,\qquad
 \lambda_\alpha(a)=\alpha(a/\bar a).
\]
We will use this with \(\alpha=c\upsilon_1^d\) and
\(c^{-1}\upsilon_1^m\), so the marked primes are determined
by \(\alpha\) itself. The conductor power uses one residue
cardinality for each conjugate pair over a split place of \(F\).

\begin{theorem}[The normalized carrier identity]
\leavevmode\label{thm:secondglobal833}
At every sufficiently large node of the prescribed family
there is a number \(K_{12}>0\), depending on the field,
rank and measure choices, such that for every prescribed
twist \(c\) there is a unique \(\gamma_{12}(c)\) with
\[
 \gamma_{12}(c)\int_{[H]}C(h)\Theta_c(x,h)\,dh=\phi(x)
           \quad(x\in G(\mathbb A)),\qquad
 |\gamma_{12}(c)|^2
  =\frac{K_{12}}{\mathcal D_d(c\upsilon_1^d)
                      \mathcal D_m(c^{-1}\upsilon_1^m)}.
\]
The same scalar works for every evaluation point and every
subsequent period test at this node.
\end{theorem}
\begin{proof}
Double \(G\) with the diagonal Siegel parabolic and integrate
on the definite space \(I_r\). Form the doubled section with
the square test of \(\Theta_c\) and its conjugate, with the
determinant restriction conventions of
Theorem~\ref{thm:normalfirst828}. In particular its corrected
graph value at \((a,1)\) at \(s=0\) is the source oscillator
coefficient for this very test. The same choice on the
diagonal makes the section invariant under simultaneous left
translation there for every \(s\). Let \(\mathcal Z_c(s)\)
be the integral of the restricted Eisenstein series against
\(\bar\phi(x_1)\phi(x_2)\) on \([G]^2\), divided by
\(\|\phi\|^2\). In the unitary Siegel--Weil formula the
Siegel rank and the dimension of the integrated space both
equal \(r\). This space is anisotropic. Thus
\href{https://msp.org/pjm/2013/264-1/pjm-v264-n1-p11-s.pdf}
{Yamana, Theorem~2} evaluates the holomorphic value of the
series at \(s=0\) as twice the ordinary theta integral in
its measures. The external target determinant characters on
the two copies cancel. At that point the product of the two
rational theta sums is
\(\Theta_c(x_1,h)\overline{\Theta_c(x_2,h)}\).
The rapid decrease of \(\phi\) absorbs the uniform moderate
growth on the two Siegel sets, as in
Theorem~\ref{thm:carrierline831}. Compactness in \(h\) and
Fubini now give, in our fixed measures,
\[
             \mathcal Z_c(0)=b\,\|B_c\|^2/\|\phi\|^2,\qquad b>0 .
\]
Only the measures enter \(b\).

Unfold in an absolute convergence half-plane. The non-graph
orbits vanish by the proper parabolic constant terms of the
cusp form, as proved in Theorem~\ref{thm:carrierline831}.
In the graph orbit the quotient variable \(x\) supplies
\(\int_{[G]}\overline{\phi(xa)}\phi(x)\,dx\), leaving
the coefficient \(\overline{M_\phi(a)}\) after normalization.
It factors by Lemma~\ref{lem:carriermatrix831}. The graph
section is a restricted tensor product of the specified
local sections, so absolute convergence gives the product
of the local integrals. Equivalently one integrates over
larger and larger sets of places, using compact volume one
outside each set. The finitely many marked data enter with
their actual squared additive norms. Theorem~\ref{thm:secondlocal832}
therefore gives, up to a single positive Haar factor \(b'\),
\[
 \mathcal Z_c(s)=b' Z_\infty(s)
 \frac{\displaystyle
       \prod_{j=1}^{m}L_{E,f}(s+j-d,\lambda_1)
       \prod_{j=1}^{2m}L_{E,f}(s+j-m,\lambda_2)}
      {\displaystyle\prod_{j=1}^{3m}L_{F,f}(2s+j,\eta^j)}
        \prod_{v\in\mathcal S_c}\frac{P_r(q_v)}{H_r(q_v,s)} .
\]
Here \(Z_\infty\) denotes the product of the real graph
integrals. At a mark both numerator characters are ramified
on each of the two components. The Euler numerator is one
there and the reciprocal denominator is \(H_r(q_v,s)\);
the last product replaces it by the actual local answer.

All finite Hecke functions on the right have their classical
meromorphic continuations. The real factor is holomorphic in
\(\operatorname{Re}s>0\) by Theorem~\ref{thm:secondlocal832}.
Thus the identity holds meromorphically there. We can approach
zero along positive real \(s\). The real factor tends to its
strictly positive value, and the correction at each mark tends
to one. In the denominator the \(j=1\) value \(L_{F,f}(1,\eta)\)
is finite and positive by the quotient of the residues of
\(\zeta_E\) and \(\zeta_F\). The factors at the other integers
are nonzero by absolute convergence.

For either numerator character put \(e=d\) or \(e=m\)
respectively. Its angular exponents at the complex places
of \(E\) have absolute value \(2e\), because the finite-order
twist has weight zero and base change is \(a\mapsto a/\bar a\).
Its finite Hecke function \(L(z)\) is nonzero at \(z=1,\ldots,e\).
For \(z>1\) the Euler product suffices. At \(1\), the character
and its square have nonzero angular type and hence holomorphic
Hecke functions. Removing the ramified factors throughout,
the Euler logarithms give
\(\zeta_E(\sigma)^3|L(\sigma)|^4|L(\sigma,\lambda^2)|\ge1\)
in the omitted-factor convention for \(\sigma>1\). This is
the sum-of-squares trigonometric inequality in
Theorem~\ref{thm:squareconductor827}. A zero at one would
outweigh the third power of the simple zeta pole. The
archimedean gamma factors \(\Gamma_{\mathbb C}(z+e)\) are
finite and nonzero at \(1-e,\ldots,0\), since their smallest
argument is one. The primitive Hecke functional equation
now gives a finite nonzero value at each of these integers
from the paired positive argument. The whole right side
therefore has a nonzero limit at zero. Since \(\mathcal Z_c\)
is holomorphic there and its value is a positive factor
times \(\|B_c\|^2\), we have \(B_c\ne0\).

It remains to compare the constants across twists without
introducing any dependence on \(x\) or on the period test.
For either block the finite conductor ideal of the primitive
character has norm \(Q_c^2\), by nontrivial tame ramification
in the two split components at every mark and unramified
data elsewhere. Its completed conductor is
\(A_c^{\rm cond}=|\operatorname{disc}E|Q_c^2\).
Compare to any reference twist \(c_*\) on the same field.
For the two central-symmetric arguments \(z,1-z\),
\(1\le z\le e\), taking absolute values in the two functional
equations gives
\[
 \left|\frac{L_c(z)L_c(1-z)}{L_{c_*}(z)L_{c_*}(1-z)}\right|
 =
 \left(\frac{A_c^{\rm cond}}{A_{c_*}^{\rm cond}}\right)^{z-1/2}
 \left|\frac{L_c(z)}{L_{c_*}(z)}\right|^2 .
\]
The gamma factors are common and cancel. Summing the exponents
\(z-1/2\) gives \(e^2/2\); the conductor power is consequently
\((Q_c/Q_{c_*})^{e^2}\). Apply this calculation to both
blocks of the displayed product. The other factors at zero
are independent of \(c\). It follows that
\[
 \frac{\|B_c\|^2}{\|B_{c_*}\|^2}
  =\frac{\mathcal D_d(c\upsilon_1^d)\mathcal D_m(c^{-1}\upsilon_1^m)}
          {\mathcal D_d(c_*\upsilon_1^d)\mathcal D_m(c_*^{-1}\upsilon_1^m)} .
\]
By Corollary~\ref{cor:carriermultiple831} the same ratio
holds for \(|a_c|^2/|a_{c_*}|^2\), and \(a_c\ne0\).
Take \(\gamma_{12}(c)=a_c^{-1}\), and define \(K_{12}\)
by the claimed identity at \(c_*\). This makes it positive
and independent of the twist. The equality \(A_c=a_c\phi\)
holds in the cuspidal space and hence at every point by
smoothness and full support of Haar measure. This proves
the theorem.
\end{proof}

\begin{corollary}[Removal of the first signed scalar]
\leavevmode\label{cor:isolateexchange833}
Let \(\gamma^-(c)\) be the first signed scalar of
Corollary~\ref{cor:signedfirst828}. On the product in which
that first class sum has already been multiplied by
\(\gamma^-(c)\), the remaining multiplier
\(\gamma_2(c)=\gamma_{12}(c)/\gamma^-(c)\) satisfies
\[
       |\gamma_2(c)|^2=
          K_{12}B_{F,m}/\mathcal D_m(c^{-1}\upsilon_1^m).
\]
\end{corollary}
\begin{proof}
In that corollary \(\delta=c^{-1}\), so
\(D_\delta=\mathcal D_d(c\upsilon_1^d)\).
The ratio of the two squared norms cancels this factor.
\end{proof}

Thus the first two exchanges supply the carrier itself on
the original group at the original level. The comparison
between its period and the integral definite theta class sum
in Theorem~\ref{thm:arithmeticclass824} concerns the third
exchange on \(I_r\oplus W'\).
\clearpage

\section{The third theta exchange}\label{nav:third-exchange}

\subsection{The smaller-rank integral in the third exchange}
\label{sec:smallthird834}
Fix the field and interval of Theorem~\ref{thm:complexseparation818},
a lift index \(j\), and a choice \(i=(\Lambda,b_0,\text{short rows})\)
as in Section~\ref{sec:refinedclass824}. Write
\[
 W=I_r\oplus W',\quad
 z_i(g)=\operatorname{diag}(g,1,1,b_0)\in U(W')(\mathbb A),
 \quad y_i(h',g)=\operatorname{diag}(h',z_i(g)).
\]
Here \(S_0\) has signature \((2m,m)\), \(r=3m\),
\(p=\dim W'\) and \(n=p+r\). At a definite class point
we take the identity real component. Retain
\(\kappa=\upsilon_1^m,\ c=\chi_M^j\vartheta\) and the
character \(\nu\) of the pair \((W,I_p)\) in
Section~\ref{sec:refinedclass824}. That pair has characters
\((\alpha,\nu)\) with \(\alpha=\nu\kappa\). The carrier
kernel on \((S_0,W')\) has characters \((\kappa,\beta_0)\)
where \(\beta_0=\nu\kappa/c\). Indeed removing the twist
and old tame part \(c\) from the first-factor character
\(\nu\) gives its fixed untwisted part, and the carrier
character is this part times \(\kappa\). Thus
\(c\beta_0=\alpha\) and \(\alpha/\nu=\kappa\).

On \((S_0,W)\) use the characters \((\kappa,\alpha)\).
At the marked split places take the position datum
\(\Psi\binom{A}{B}=\chi(\det A)F_\Lambda(B)\) of
Lemma~\ref{lem:residuepartners825}. Elsewhere take the
unit finite vacua. At the real places use the polynomial
\(Q\) of Proposition~\ref{prop:carrierpolynomial825}, with
its order of the rows. Denote this kernel by \(\theta_{\Psi,Q}\).
The smaller rank refers to the space \(S_0\) integrated in
\[
 \mathcal S_{i,c}(g)=
 \sum_{h'\in\mathcal C_r}\frac{C_c(h')}{d_{h'}}
       \int_{[G]}\theta_{\Psi,Q}(x,y_i(h',g))\,dx,\qquad
 C_c(h')=c(\det h')C(h').
\]
The weight \(d_{h'}^{-1}\) is that of the definite class
of the full stabilizer, as in Theorem~\ref{thm:arithmeticclass824}.

\begin{theorem}[The contracted smaller-rank integral]
\leavevmode\label{thm:smallthird834}
For every \(g\) and every prescribed twist one has
\[
     \gamma_{12}(c)\mathcal S_{i,c}(g)
           =N_R\,\mathcal B_i(g),\qquad
           N_R=|\operatorname{GL}_{r}(R)|.
\]
Each inner integral at the stated points is absolutely
convergent. The formula holds for the entire family of
\(i,g\) with the same \(\gamma_{12}(c)\).
\end{theorem}
\begin{proof}
Write \(K_{1,c}(x,h')\) for the kernel on \((S_0,I_r)\)
with characters \((\kappa,c)\), determinant test and constant
real polynomial, and \(K_{2,i}(x,z_i(g))\) for the carrier
kernel on \((S_0,W')\) with characters \((\kappa,\beta_0)\),
datum \(F_\Lambda\), and polynomial \(Q\).
On the displayed block diagonal one has exactly
\[
     \theta_{\Psi,Q}(x,y_i(h',g))
                 =K_{1,c}(x,h')K_{2,i}(x,z_i(g)).
\]
Indeed the position spaces for the two target summands are
a direct sum, and rational summation in the product
polarization runs over all pairs of rational position vectors.
The unmarked and marked Schwartz data are products.
Proposition~\ref{prop:carrierpolynomial825} puts all the
polynomial variables in the second factor, leaving the
Gaussian in the first. The two multipliers on the shared
source multiply to \(\alpha(\det x)^{-1}\). The two target
multipliers are the restriction of \(\kappa\) on the full
target determinant. The underlying oscillator on an orthogonal
sum is the tensor product in these normalizations, as in
the proof of Theorem~\ref{thm:arithmeticclass824}. Absolute
convergence of the rational Schwartz sums at each adelic point
allows their multiplication.

Changing the first character of \(K_{1,c}\) to \(c\kappa\)
multiplies it by \(c(\det h')\) and gives exactly \(\Theta_c\)
of Section~\ref{sec:secondcarrier829}. For each fixed \(h'\),
Proposition~\ref{prop:secondcusp829} gives cuspidality in \(x\)
already for the kernel before class integration. Its
holomorphic scalar weight, finite level and smoothness put
it in the cuspidal space of the given weight, so it is
rapidly decreasing on Siegel sets. The second factor is of
moderate growth in \(x\) for fixed \(z_i(g)\).
This proves absolute convergence. There are finitely many
class terms. Passing the sum through the integral gives
\[
 \mathcal S_{i,c}(g)=
 \int_{[G]}\left(\sum_{h'\in\mathcal C_r}
     \frac{C(h')\Theta_c(x,h')}{d_{h'}}\right)
                     K_{2,i}(x,z_i(g))\,dx .
\]
Full level and compact real invariance in the combined
summand identify the inner sum with \(A_c(x)\). This uses
the same quotient measure as in Theorem~\ref{thm:secondglobal833}.
Multiplication by \(\gamma_{12}(c)\) therefore inserts
\(\phi(x)\). The resulting integral is precisely
\(\mathcal I_i(F_\Lambda;g)\), including the inflation at
the new marks. By Theorem~\ref{thm:carriermultiplicity825}
it equals \(N_R\mathcal B_i(g)\). This proves the identity.
\end{proof}

\begin{corollary}[Separation by the smaller-rank integrals for every twist]
\leavevmode\label{cor:smallthirdseparation834}
Fix any nonzero good form covered by
Theorem~\ref{thm:complexseparation818}. Choose a complementary
test, a character projection and a first-factor index \(i\)
for which the carrier product period is nonzero. Replacing
\(\mathcal B_i\) by \(\mathcal S_{i,c}\) gives a nonzero
projected period for every prescribed twist \(c\). Its
squared absolute value is multiplied by
\[
        \frac{N_R^2}{K_{12}}\,
           \mathcal D_d(c\upsilon_1^d)
           \mathcal D_m(c^{-1}\upsilon_1^m).
\]
\end{corollary}
\begin{proof}
Hold fixed all the data other than the indicated first
factor. The character projection and period functional
in that factor are linear or conjugate-linear according
to the pairing convention. Apply Theorem~\ref{thm:smallthird834}
and take squared absolute values. Theorem~\ref{thm:secondglobal833}
evaluates the scalar, which is strictly positive.
\end{proof}

The larger-rank expression at these same arguments is
\(L_{i,\vartheta}(g)\) in Theorem~\ref{thm:arithmeticclass824}.
Its inner integral is on \(U(I_p)\), has characters
\((\alpha,\nu)\), and uses the residue partner \(\Phi\)
and the stacked polynomial \(P\). In particular the still
needed comparison is between this expression and the
specific \(\mathcal S_{i,c}\) above. An identity
\(\delta_3(c)L_{i,\vartheta}(g)=\mathcal S_{i,c}(g)\)
would put the multiplier \(\gamma_{12}(c)\delta_3(c)/N_R\)
in front of the arithmetic expression. The preceding proof
determines the carrier and the integer \(N_R\); it does not
supply that remaining identity or the absolute value of
\(\delta_3(c)\).
\clearpage

\subsection{Frame integrals at the split conductor of the third exchange}
\label{sec:marksections835}
Retain \(n=p+r\) and \(\alpha=\nu\kappa\) from
Section~\ref{sec:smallthird834}. At a split place write \(k=F_v\),
\(O=O_v,\ K_b=\operatorname{GL}_b(O)\), and let \(\xi\) be the
value of the Weil splitting character on the chosen uniformizer.
Read norm-one characters on the chosen component. For
\(t\in k^\times,\ s\in\mathbb C\) put
\[
 \omega_s(t)=\alpha_v(t)^{-1}\xi^{\,n v(t)}|t|^{\,n/2+s}.
\]
The parameter \(s\) varies a local determinant integral.
At \(s=0\) its weight is that of the actual oscillator.
We use Haar measures with \(K_b\) of mass one for all \(b\).

Fix \(k^n=k^r\oplus k^p\), put \(E_r=\binom{1_r}{0}\),
\(F_p=(0\ \ 1_p)\), and for \(t\in\operatorname{GL}_n(k)\) set
\(Y_t=tE_r,\ X_t=F_p t^{-1}\). Thus \(Y_t\) is a basis of an
\(r\)-plane and \(X_t\) a surjection with precisely this
kernel. Define
\[
 \begin{split}
 J_p(\Phi;t,s)&=\int_{\operatorname{GL}_p(k)}
          \omega_s(\det h)\Phi(hX_t)\,dh,\\
 J_r(\Psi;t,s)&=\int_{\operatorname{GL}_r(k)}
          \omega_s(\det u)\Psi(Y_tu)\,du .
 \end{split}
\]
Begin with locally constant functions compactly supported in
the full-rank matrix loci, including zero extensions of
the marked tests. These integrals then exist for every \(s\).
For example the map \(h\mapsto hX_t\) parametrizes the
surjections with the given kernel. It is a homeomorphism
onto that fiber, which is closed in the full-rank locus.
The inverse sends \(X\) to \(XT\) for any right inverse \(T\)
of \(X_t\). The same argument applies to \(Y_tu\).
Compact support in the full-rank locus therefore gives
compact integration domains.

\begin{theorem}[Equality of the marked sections in complementary frames]
\leavevmode\label{thm:marksections835}
For the pair \((\Phi,\Psi)\) of Lemma~\ref{lem:residuepartners825},
inflated at the marks to the integral boxes, one has
\[
                 J_p(\Phi;t,s)=\omega_s(\det t)J_r(\Psi;t,s)
\]
at every marked split place and every \(t\in\operatorname{GL}_n(k)\).
At \(s=0\) this equality persists under every simultaneous
translation by the actual split group \(U(W)(F_v)\) on the
two position models.
\end{theorem}
\begin{proof}
The additive trace character of the old finite ring factors
over its local components. Moreover \(\operatorname{GL}_r(R)\)
is the product of the local groups, so its sum defining
\(F_\Lambda\) is a product of independent local sums. On a
new component the inflated \(F_\Lambda\) contributes the
constant function on the integral box. Thus the marked tests
in the assertion are the indicated local factors at old
precision two and at new precision one.

Suppose first that \(t\in K_n\). Then \(X_t\) is an integral
surjection with integral right inverse, and \(Y_t\) is an
integral injection with integral left inverse. The support
of \(\Phi\) requires its last square block to have unit
determinant, so \(hX_t\) must be an integral surjection.
This is equivalent to \(h\in K_p\) as far as primitivity
is concerned, since \(X_t(O^n)=O^p\). Similarly \(Y_tu\)
must be an integral primitive frame, forcing \(u\in K_r\).
Indeed \(Y_t\) identifies \(O^r\) with the intersection
of its column space with \(O^n\), and a primitive frame in
that space spans this intersection. Conversely multiplication
by these unit groups preserves primitivity and the tests
obey
\[
     \Phi(hX)=\chi_v(\det h)\Phi(X),\qquad
     \Psi(Yu)=\chi_v(\det u)\Psi(Y).
\]
For the first identity the quotient \(z^{-1}y\) is fixed
by left multiplication of the two blocks. For the second,
the determinant of the top block contributes \(\chi_v(\det u)\)
and the lower test is right invariant. The same equations
hold at the zeros of the tests. On units \(\omega_s\) is
\(\chi_v^{-1}\). Thus the two integrals are respectively
\(\Phi(X_t)\) and \(\Psi(Y_t)\), without a counting factor.
Their relation is the determinant identity of
Lemma~\ref{lem:residuepartners825} on the local residue ring.

For general \(t\), the lattice in the column space given by
its intersection with \(O^n\) is saturated, hence a direct
summand over the discrete valuation ring. An integral basis
of it extends to one of \(O^n\). Consequently
\[
 t=k_0\begin{pmatrix}a&b\\0&d_0\end{pmatrix},
 \qquad k_0\in K_n,\quad a\in\operatorname{GL}_r(k),\
 d_0\in\operatorname{GL}_p(k).
\]
One has \(Y_t=Y_{k_0}a\) and \(X_t=d_0^{-1}X_{k_0}\).
Changes of variables in the two Haar integrals multiply
their values at \(k_0\) by \(\omega_s(\det d_0)\) and
\(\omega_s(\det a)^{-1}\), respectively. Their quotient
relation is therefore the one asserted at \(t\), including
when both sides vanish.

Finally, at \(s=0\) write \(\omega=\omega_0\).
For \(g\in U(W)(F_v)=\operatorname{GL}_n(k)\), the two
translates have the forms
\[
 \Phi_g(X)=b(g)\Phi(Xg),\qquad
 \Psi_g(Y)=a(g)\Psi(g^{-1}Y).
\]
If \(z=\det g\), the position Jacobians and splitting
characters give
\[
 b(g)=\nu_v(z)^{-1}\xi^{p v(z)}|z|^{p/2},\qquad
 a(g)=\kappa_v(z)\xi^{-r v(z)}|z|^{-r/2}.
\]
For the first substitution the \(p\) rows are acted on
from the source, each giving the square-root Jacobian and
one chosen splitting phase. In the second substitution the
\(r\) columns are acted on from the target, with the inverse
Jacobian and conjugate splitting phase. These are the source
and target computations of Lemma~\ref{lem:secondhecke829}.
The external multipliers are the displayed characters,
in their respective source and target conventions.
Since \(\alpha=\nu\kappa\), \(b(g)/a(g)=\omega(z)\).
Now \(X_tg=X_{g^{-1}t}\) and \(g^{-1}Y_t=Y_{g^{-1}t}\).
Applying the identity there and inserting the two multipliers
gives \(J_p(\Phi_g;t,0)=\omega(\det t)J_r(\Psi_g;t,0)\).
\end{proof}

\begin{proposition}[The kernel of integration along bases]
\leavevmode\label{prop:framecoinvariants835}
On the test functions compactly supported in the full-rank
locus, the kernel of \(\Phi\mapsto J_p(\Phi;\cdot,s)\)
is spanned by
\(\omega_s(\det h)\Phi(h\,\cdot)-\Phi\) for
\(h\in\operatorname{GL}_p(k)\).
The kernel on the column side is spanned by
\(\omega_s(\det u)\Psi(\cdot\,u)-\Psi\) for
\(u\in\operatorname{GL}_r(k)\).
In particular any linear evaluation on the row data invariant
under the indicated weighted changes of basis has the same
value for two data with the same frame integral. The
corresponding assertion holds for column data.
\end{proposition}
\begin{proof}
Invariance of the integral under the displayed operations
follows by translation of Haar measure. For the converse,
project the full-rank row space to the Grassmannian by
sending a surjection to its kernel. On each coordinate minor
chart this is a product with \(\operatorname{GL}_p(k)\):
use the frame with the chosen nonsingular \(p\)-column
minor equal to the identity. Use a disjoint finite partition
by compact open subsets inside such charts on the projection
of the support. Such a partition exists by compactness
and the compact open coordinate basis. A locally constant
compactly supported function in a product chart is a finite
sum of indicators of compact open rectangles, by compactness
of its support. Refining the base partition writes it on
each piece as one fixed fiber function \(F(h)\).

If all frame integrals vanish, the ordinary integral of
\(\omega_s(\det h)F(h)\) on this fiber is zero.
It is compactly supported and right invariant under a
sufficiently small compact open subgroup \(J\). Write it
as a finite sum \(\sum_l c_l{\bf1}_{h_lJ}\). The cosets
have equal measure, so \(\sum_l c_l=0\). The function is
a linear combination of differences
\({\bf1}_{h_lJ}-{\bf1}_J\). Undoing the weight identifies
these differences with differences under the stated weighted
left change of basis. Extending each term by zero on the
other base pieces proves the row claim. On the column
side trivialize by a nonsingular row minor, use left
invariance of the weighted fiber function under a compact
open subgroup and take differences of its right translates.
Both groups are unimodular. The statements about invariant
evaluations follow by linearity.
\end{proof}

These assertions describe how a marked local functional
can depend on its matrix test. Equality between the actual
adelic functionals also depends on the spaces, real data
and rational summation.
\clearpage

\subsection{The split frame multiplier at the third critical block}
\label{sec:thirdsplitprod836}
The dimensions in Section~\ref{sec:smallthird834} satisfy
\[
       e=(p-r)/2,\qquad n/2=r+e,\qquad e=(k+7m)/2 .
\]
Indeed \(\dim W_j=k+m-1\), \(p=\dim W_j+9m+1\) and
\(r=3m\); both \(k\) and \(m\) are even. In this subsection
\(k\) in the dimension formula denotes the original source
rank. At a split finite place write \(q=q_v\) for the residue
cardinality and \(O=O_v\) for the local integer ring.

\begin{lemma}[Relative multipliers of the two split lattice sections]
\leavevmode\label{lem:thirdsplitvacuum836}
At an unmarked split place use the two integral matrix
indicators \(\Phi^0,\Psi^0\) of sizes \(p\)-by-\(n\)
and \(n\)-by-\(r\). Put
\(u_v=\alpha_v(\varpi)^{-1}\xi^{\,n}\) at that place.
Then, in the notation of Section~\ref{sec:marksections835},
\[
 J_p(\Phi^0;t,s)=D_{c,v}(s)\,\omega_s(\det t)J_r(\Psi^0;t,s),
 \qquad
 D_{c,v}(s)=\prod_{\ell=1}^{2e}
            (1-u_v q^{-(s+\ell-e)})^{-1}.
\]
This holds for all frames \(t\) in an absolute convergence
half-plane, then as an identity of meromorphic functions.
At a marked split place the analogous identity for the
marked tests has \(D_{c,v}(s)=1\).
\end{lemma}
\begin{proof}
For \(t\in K_n\) the matrix \(X_t\) has integral right
inverse and \(Y_t\) integral left inverse. Thus \(hX_t\)
is integral exactly when \(h\) is, and \(Y_tu\) is integral
exactly when \(u\) is. In dimension \(b\) the resulting
one-sided integral on \(\operatorname{GL}_b(F_v)\cap M_b(O)\)
has weight \(\omega_s(\det h)\). At an unmarked place its
value is
\[
       \prod_{i=1}^b(1-u_v q^{i-1-n/2-s})^{-1}.
\]
This follows either by the elementary determinant generating
function in Theorem~\ref{thm:splitzeta827}, with additive
and multiplicative Haar measures related by
\(dX=P_b(q)|\det X|^b\,d^\times X\), or by the one-sided
integral in Lemma~\ref{lem:splitspherical832} evaluated at
the constant representation. Initially take
\(\operatorname{Re}(n/2+s)>p-1\), so both integrals converge.
Taking the ratio for \(b=p\) and \(b=r\) leaves just
\(i=r+1,\ldots,p\). As \(n/2=r+e\), the exponents in the
remaining denominators are \(-s-e,-s-e+1,\ldots,-s+e-1\).
This gives \(D_{c,v}\) on integral frames, where
\(\omega_s(\det t)=1\). For all frames apply the parabolic
decomposition and change of variables of
Theorem~\ref{thm:marksections835}, initially where both
integrals converge. That argument also gives their meromorphic
extensions frame by frame. The last assertion, with the
actual residue functions, is Theorem~\ref{thm:marksections835}.
\end{proof}

Put \(\sigma_c=\alpha^{-1}\upsilon_1^{\,r+e}\) on the adelic
norm-one torus and \(\lambda_c(a)=\sigma_c(a/\bar a)\) on the
ideles of \(E\). This notation for \(\sigma_c\) is confined
to this third block. Its finite phase at the chosen component
of an unmarked split place is just \(u_v\), since
\(\upsilon_{1,v}(\varpi)=\xi^2\). The conjugate component
has the inverse value. The tame conductor at each old and
new mark has exponent one on each of the two components.
Let \(c_*\) be a reference twist with the same field and
lift index. For \(\operatorname{Re}s\) sufficiently large define
\[
 \mathcal R_3(s;c,c_*)=
 \prod_{v\ {\rm split}}
 \frac{D_{c,v}(s)\overline{D_{c,v}(\bar s)}}
      {D_{c_*,v}(s)\overline{D_{c_*,v}(\bar s)}} .
\]
Taking the conjugate at \(\bar s\) gives the second component
as a holomorphic function of \(s\). This is a relative
product of frame multipliers.

\begin{theorem}[The value of the relative split frame multiplier]
\leavevmode\label{thm:thirdsplitprod836}
The product \(\mathcal R_3\) continues meromorphically and is
regular and nonzero at zero. Its value there is positive and
\[
       \mathcal R_3(0;c,c_*)=
       \frac{\mathcal D_e(\sigma_c)}
            {\mathcal D_e(\sigma_{c_*})},\qquad
 \mathcal D_e(\sigma_c)=Q_c^{e^2}
             \prod_{z=1}^e|L_{E,f}(z,\lambda_c)|^2 .
\]
\end{theorem}
\begin{proof}
The unitary phase in the lemma gives
\[
 \mathcal R_3(s;c,c_*)=
 \prod_{\ell=1}^{2e}
       \frac{L_{E,f}(s+\ell-e,\lambda_c)}
            {L_{E,f}(s+\ell-e,\lambda_{c_*})}
\]
in the initial half-plane. At an unmarked split pair its two
Euler factors are the lemma and its conjugate. At a mark for
one of the two twists its own two numerator-character Euler
factors are one, as is its frame multiplier. This works
also when the two sets of new marks differ. At an inert
place the norm-one unit character is unramified and trivial,
and a common uniformizer maps to one under \(a\mapsto a/\bar a\).
Thus the inert Euler factors coincide. There are no other
finite places by Lemma~\ref{lem:signkernel823}. This proves
the product identity and gives the continuation by the
classical Hecke functional equations.

At a real embedding \(\alpha\) has weight \(m-\mu_\tau\)
as in Section~\ref{sec:refinedclass824}. The weight of
\(\sigma_c\) is \(a_\tau=\mu_\tau+2m+e\), independent of
the twist. Hence its base-change character at the associated
complex place has angular exponent of absolute value
\(2a_\tau\), with \(a_\tau\ge e\). At the integers
\(1-e,\ldots,0\) the common gamma factors
\(\Gamma_{\mathbb C}(z+a_\tau)\) are finite and nonzero.
There are no positive-integer zeros: at \(z>1\) use the
Euler product, and at \(z=1\) use the trigonometric
inequality of Theorem~\ref{thm:squareconductor827};
the character and its square have nonzero angular type.
The functional equation then shows that the finite Hecke
functions are nonzero throughout \(1-e,\ldots,e\).

For the sign in the claimed value, the finite Hecke function
of each of these characters equals that of its inverse,
by conjugation of ideals of \(E\). Indeed conjugation sends
\(\lambda_c\) to its inverse and preserves the absolute
norm, and the ramified set consists of conjugate pairs.
The common \(L\)-function is real on the real axis. Its
root number in the completed functional equation could
a priori contribute at the reflected arguments. Here it
is \(+1\). We give the local-to-global reason. The character
\(\lambda_c=\sigma_c\circ(1-\mathrm{conj})\) is trivial on
the ideles of \(F\). Use a fixed nonzero \(\Delta\in E\)
of relative trace zero and the standard additive character
of \(\mathbb A_E/E\) obtained by composing any nontrivial
additive character of \(\mathbb A_F/F\) with
\(z\mapsto\operatorname{Tr}_{E/F}(\Delta z)\).
It is trivial on \(\mathbb A_F\). At an inert finite place
\(\lambda_c\) is trivial, whose local root number in
unitary normalization is one for every nontrivial additive
character. At a split place the two characters are inverse.
The two additive characters are opposite under the canonical
identifications with \(F_v\), since the two components of
\(\Delta\) are opposites. The product of the two local
epsilon factors at the central point is one by the local
Tate functional equation applied twice, using opposite
Fourier kernels and self-dual measures. This includes tame
ramification. At a real place of \(F\) the associated
character of \(\mathbb C^\times\) and the additive character
are independent of \(c\), so their local root number is the
same for \(c\) and \(c_*\). In fact, for an exponent \(2a\)
the standard complex epsilon factor is \(i^{|2a|}\) up to
the Fourier sign convention, and scaling the real-trace
additive character by a purely imaginary scalar contributes
its angular character value \((\pm i)^{2a}\); their product
is one. Consequently the global root numbers in particular
have ratio one. For the asserted comparison only this last
ratio is needed.

The completed conductor is \(A_c=|\operatorname{disc}E|Q_c^2\).
Apply the primitive functional equations to the pair \(z,1-z\)
for \(1\le z\le e\). The root-number ratio is one, the real
archimedean factors cancel, and at positive integers the
dual value is the conjugate of the original one. The
contribution of the pair to the ratio is
\[
          (A_c/A_{c_*})^{z-1/2}
              \left|\frac{L_{E,f}(z,\lambda_c)}
                          {L_{E,f}(z,\lambda_{c_*})}\right|^2 .
\]
Multiplying these positive quantities gives the formula.
The conductor exponent is \(e^2\) for the product of the
residue sizes on the \(F\)-side, by
\(\sum_{z=1}^e(z-1/2)=e^2/2\).
\end{proof}

The theorem evaluates the split frame contributions at the
rank difference \(p-r\). Identification of the relative
norms of the two global theta periods with this product
requires a comparison of their rational theta sums and
their real polynomial data.
\clearpage

\subsection{One standard residue test for the global comparison}
\label{sec:thirdstandard837}
We now use translations of the common group \(U(W)\) in
the two actual theta integrals. The translations can mix
\(I_r\) and \(W'\). We first justify the integral on the
indefinite group at such points.

\begin{lemma}[Rapid decrease from a split full-rank test on the source]
\leavevmode\label{lem:rankdecay837}
Let \(S\) be a Hermitian space, \(T\) another Hermitian space,
and use the theta kernel on \(U(S)\times U(T)\) with Schwartz
data. Suppose at one split finite place its position test on
\(\operatorname{Hom}(S_v,T_v)\) in the chosen component is
compactly supported in the injective matrices. For a fixed
target point the theta kernel is rapidly decreasing on Siegel
sets of the source. The estimates are locally uniform in the
target point and hold for Schwartz derivatives at infinity.
\end{lemma}
\begin{proof}
If the source is anisotropic its quotient is compact. Otherwise
choose a rational maximal isotropic flag in \(S\) and a dual
flag, giving a rational Witt decomposition with isotropic
basis \(e_1,\ldots,e_l\). In the mixed polarization the
rational summation variables include the images
\(y_1,\ldots,y_l\in T(E)\) of these basis vectors, written
in the two split components at a split place. Denote remaining
rational position coordinates by \(z\). In the chosen split
position model passage to this polarization keeps the columns
of the matrix on the indicated isotropic subspace in the
chosen component and Fourier transforms the dual columns,
as in Proposition~\ref{prop:secondcusp829}. In particular
the mixed test vanishes at \(y_1=0\): before the partial
Fourier transform the retained column on \(e_1\) would be
zero. This remains true after arbitrary local source and
target translations, since both act here by invertible
matrix substitutions and scalar multipliers.

Apply adelic reduction theory to the restriction of scalars
of \(U(S)\), whose center is anisotropic. Its Siegel sets
have the form \(\Omega aK\), using a compact set \(\Omega\)
in the minimal parabolic, a compact \(K\), and the positive
split chamber with the simple roots bounded below. This is
the adelic reduction theorem in
\href{https://www.claymath.org/library/cw/arthur/pdf/61.pdf}
{Arthur, \emph{An introduction to the trace formula},
Theorem~8.1}. In Witt coordinates take \(a\) trivial at
the finite places, with diagonal \(a_1,\ldots,a_l\) on
the isotropic flag at each real embedding, and the inverse
diagonal on the dual. These positive parameters are shared
across embeddings for the rational split component. The
ratios \(a_i/a_{i+1}\) and the last parameter \(a_l\) are
bounded below. All \(a_i\) are therefore bounded below and
are \(O(1+a_1)\), with constants allowed to depend on the
fixed group and Siegel set. The conjugates \(a^{-1}\Omega a\)
lie in a compact set: inverse conjugation contracts positive
root coordinates up to constants, and the Levi centralizes
the split component. Fixed finite translates of these sets
can be absorbed in the compact data.

Move \(a\) to the left in the pair action. The remaining
Schwartz data range in a bounded archimedean Schwartz family,
and their finite parts have common compact support in the
mixed model. This follows from continuity of the oscillator
on the Schwartz space and local constancy at the finite
places. All retain the marked vanishing just proved.
The rational coordinates that can contribute lie in a
fixed product of fractional lattices in \((y_1,\ldots,y_l,z)\).
The action of \(a\) in the mixed model substitutes \(a_i y_i\)
for \(y_i\), fixes \(z\), and multiplies by determinant
characters and the square root of the position Jacobian.
It is the linear action on the isotropic position half and
the identity on the remaining Witt summand. Its absolute
multiplier is bounded by a fixed power of \(1+a_1\).
Every occurring \(y_1\) is nonzero, so its Euclidean norm
over all embeddings is bounded below on its lattice.
It follows that
\(1+\|(a_1y_1,\ldots,a_ly_l,z)\|\) dominates both
\(1+a_1\) and \(1+\|(y_1,\ldots,y_l,z)\|\) up to fixed
constants. Use an arbitrarily high power of Schwartz decay
for the first bound and a power exceeding the lattice
dimension for the second. Summing the latter lattice
majorant gives \(O_N((1+a_1)^{-N})\) for every \(N\),
after absorbing the Jacobian. The height on these Siegel
sets grows at most polynomially in \(1+a_1\). The same
argument applies after differentiation and for target points
in a compact set. This proves the lemma.
\end{proof}

At a marked place, for a locally constant function \(f\)
of compact support on the \(p\)-by-\(r\) matrices, define
the following functions in the open square-block charts:
\[
 \Psi_f\binom{A}{B}={\bf1}_{A\in K_r}\chi_v(\det A)f(BA^{-1}),
 \qquad
 \Phi_f(y,z)={\bf1}_{z\in K_p}\chi_v(\det z)f(-z^{-1}y).
\]
A value outside the indicated condition is defined to be
zero. Put \(f^0={\bf1}_{M_{p,r}(O)}\) and use the associated
pair as the standard test at this mark. At all marks together
the actual tests are obtained by taking \(f\) equal to
\(F_\Lambda\) on the integral box and zero off it. Indeed
\(F_\Lambda\) is invariant under right multiplication by
units, including \(-1\). At the new marks its factor is
exactly \(f^0\).

Let \(V_{\Phi,P}(y)\) be the integral on \([U(I_p)]\) for the
large pair at \(y\in U(W)(\mathbb A)\) and
\(V'_{\Psi,Q}(y)\) the integral on \([G]\) for the small
pair at the same \(y\), with the characters of
Section~\ref{sec:smallthird834}. Write \(V^0_P,V^{\prime0}_Q\)
for the integrals with the standard tests at every mark,
retaining the same unmarked and real data. All have their
ordinary meaning: the large group of integration is anisotropic,
and the small test is supported on injections at each mark
so Lemma~\ref{lem:rankdecay837} applies.

\begin{theorem}[Common finite translates for both global integrals]
\leavevmode\label{thm:thirdstandard837}
For the prescribed \(F_\Lambda\) there are points \(g_C\)
of \(U(W)(\mathbb A)\), supported at the old split places,
and explicit numbers \(b_C\), the same for the two sides,
such that at every \(y\)
\[
       V_{\Phi,P}(y)=\sum_C b_C V^0_P(yg_C),\qquad
       V'_{\Psi,Q}(y)=\sum_C b_C V^{\prime0}_Q(yg_C).
\]
Here \(C\) runs over \(M_{p,r}(R)\) at the old modulus
and, with \(N_R=|\operatorname{GL}_r(R)|\),
\[
       \sum_C |b_C|^2=
            N_R\,|\{\gamma\in\operatorname{GL}_r(R):
                                   \gamma\Lambda=\Lambda\}|.
\]
The translates are independent of the new marks.
\end{theorem}
\begin{proof}
Work first at one mark. A ball in the graph coordinate
has the form \(B_0+\varpi^L M_{p,r}(O)\), with \(L\ge0\).
Put
\[
 g=\begin{pmatrix}1_r&0\\ B_0&\varpi^L1_p\end{pmatrix}
       \in U(W)(F_v),\qquad h=\varpi^L1_p\in U(I_p)(F_v),
\]
using the chosen split coordinates. Use \(a(g),b(g)\) and
\(\omega=\omega_0\) of Theorem~\ref{thm:marksections835}.
On the small kernel the action on the standard test is
\(a(g)\Psi_{f^0}(g^{-1}Y)\). Its upper block is still \(A\),
and the lower block becomes \(\varpi^{-L}(B-B_0A)\).
Thus this vector is \(a(g)\Psi_f\) for the indicator \(f\)
of the ball. On the large kernel apply \(g\) at \(W\) and
\(h\) at \(I_p\). The matrix argument of the standard test
is
\[
 h^{-1}(y_1,z)g=\bigl(\varpi^{-L}(y_1+zB_0),z\bigr).
\]
The condition is \(z\in K_p\) with
\(-z^{-1}y_1\in B_0+\varpi^L M_{p,r}(O)\), and the unit
character is still \(\chi_v(\det z)\).
The two multipliers are \(b(g)\) and
\(\omega(\det h)^{-1}\). The latter is the target
determinant character \(\alpha_v\), the conjugate splitting
phase for \(n\) columns and the inverse square-root
Jacobian. Since \(\det g=\det h\) and \(b(g)/a(g)=\omega(\det g)\),
their product is exactly \(a(g)\). This proves the two
equalities of local Schwartz vectors with the same scalar.

Use \(L=2\) and any lift of the residue center \(C_v\) at
each old place; at new places the standard box uses the
identity. Extend the split matrices to adelic group
elements, with the dual components forced by the unitary
condition and identity elsewhere. Write \(g_C\) for the
product of the \(g\)'s. The old boxes partition the integral
matrices, with \(F_\Lambda\) constant on each box. By
linearity the coefficients are
\[
          b_C=F_\Lambda(C)\prod_{v\ {\rm old}} a(g_{C,v})^{-1}.
\]
In the larger integral the target translation by the product
of the \(h\)'s disappears by right invariance of Haar measure
on \([U(I_p)]\). The smaller integral only has the translation
by \(g_C\) in the external variable. Absolute convergence and
finiteness of the sums justify the result on both quotients.

For the norm, \(|a(g)|=|\det g|^{-r/2}\), so that the
squared modulus of the product of \(a(g)^{-1}\) is
\(|R|^{-pr}\). The primitive additive character of \(R\)
gives a perfect matrix trace pairing: if one matrix entry
is nonzero, multiplying that entry by a suitable element of
\(R\) makes its additive character nontrivial, by primitivity.
Using trace cyclicity, the Fourier frequencies in
\(F_\Lambda(C)\) are \(\gamma\Lambda\), with their actual
multiplicities. Orthogonality on the \(|R|^{pr}\) matrices
\(C\) yields
\[
 \sum_C |F_\Lambda(C)|^2
  =|R|^{pr}\#\{(\gamma,\gamma'):\gamma\Lambda=\gamma'\Lambda\}
  =|R|^{pr}N_R\,|\operatorname{Stab}(\Lambda)|.
\]
Multiplication by \(|R|^{-pr}\) proves the formula.
\end{proof}

In particular an identity between the two standard global
integrals at all \(y\), with one scalar at the fixed twist,
would pass to every prescribed residue choice with precisely
that scalar. Restriction to \(y_i(h',g)\) and summation against
\(C_c(h')/d_{h'}\) then give the two functions of
Section~\ref{sec:smallthird834}.
\clearpage

\subsection{The common real compact type in the third comparison}
\label{sec:thirdreal838}
At a real embedding put \(b=2m,\ N=p+b,\ T=\mu_\tau\) and
\(L=T+n-2m\). The positive block of \(W\) has dimension
\(N\) and its negative block has dimension \(m\).
For arbitrary \(b\)-by-\(N\) complex matrices
\(\eta_1,\ldots,\eta_T\), write the two real polynomial
vectors of Section~\ref{sec:carrierdata825} in the form
\[
 P_\eta(Z)=\prod_{l=1}^T\det\binom{\eta_l}{Z},\qquad
 Q_\eta(Y,D)=(\det D)^L\prod_{l=1}^T\det(\eta_lY).
\]
Here \(Z\) is \(p\)-by-\(N\), \(Y\) is \(N\)-by-\(b\), and
\(D\) is the \(m\)-by-\(m\) negative block on the smaller
side. The orders of the source coordinates in the determinants
are fixed as before. The constant polynomial has its usual
Gaussian in either Fock space. For the norm computation take
the unit Bargmann normalization in which a monomial of
multi-index \(a\) has squared norm \(\prod_j a_j!\).

\begin{theorem}[The real identification and its norm on the polynomial type]
\leavevmode\label{thm:thirdreal838}
The rule \(P_\eta\mapsto Q_\eta\) extends uniquely to an
isomorphism between the two spans of these vectors. With
the characters on the actual two theta kernels it commutes
with \(U(N)\times U(m)\) in \(U(W)(F_\tau)\). Each span is
irreducible under \(U(N)\). In the displayed normalization
the map multiplies every squared norm by
\[
      \frac{G_m(L)\,G_b(T)}{G_p(T)},\qquad
      G_h(t)=\prod_{i=1}^h\frac{(t+i-1)!}{(i-1)!}.
\]
In particular the factor is independent of the row selections
and of the tame twist.
\end{theorem}
\begin{proof}
For a full-rank \(Z\), exterior multiplication of its \(p\)
rows identifies their wedge, using the fixed volume form,
with a decomposable \(b\)-vector in \(\bigwedge^b\mathbb C^N\).
More explicitly, the map on a \(b\)-tuple of covectors
\(\eta\) which sends it to \(\det\binom{\eta}{Z}\) is the
alternating functional obtained by contracting the ordered
volume vector by these \(p\) rows. Its underlying \(b\)-plane
is \(\ker Z\): this follows in a basis in which the last
\(p\) covectors span the rows of \(Z\), and change of basis
preserves the assertion. Choose as columns of \(Y\) a basis
of this kernel scaled to have exactly that exterior product.
Then \(\det\binom{\eta}{Z}=\det(\eta Y)\) for every \(\eta\)
at once. Conversely, every full-rank \(Y\) gives such a
full-rank \(Z\), by choosing and scaling a basis of the
annihilator. Therefore a linear combination of the products
on one side is zero precisely when its counterpart on the
other is zero. One can test on full-rank matrices, which
are dense; the independent factor \((\det D)^L\) can be
tested at \(D=1\). This proves that the map of spans is
well defined and bijective, including \(T=0\).

We detail the compact actions because determinant characters
enter the asserted equivariance. For \(k_+\in U(N)\) on the
larger side \(Z\) is substituted by
\(Z(k_+^{\mathsf t})^{-1}\). All rows and these columns have
positive Hom sign, the positive source vacuum has weight zero,
and the external source multiplier \(\nu^{-1}\) contributes
\((\det k_+)^T\). In each stacked determinant multiply all
columns by \(k_+^{\mathsf t}\). This turns the transformed
lower rows back to \(Z\) and the upper rows into
\(\eta_l k_+^{\mathsf t}\), with a factor
\((\det k_+)^{-1}\). The \(T\) factors cancel the external
multiplier. On the smaller side the variable substitution is
\(Y\mapsto k_+^{\mathsf t}Y\). Its target vacuum contributes
\((\det k_+)^{-m}\) from the negative source columns, cancelled
by the target action of \(\kappa=\upsilon_1^m\).
Thus both actions take \(\eta_l\) to \(\eta_l k_+^{\mathsf t}\).

For \(k_-\in U(m)\) in the negative block of \(W\), the
larger source action has vacuum exponent \(p\) and external
exponent \(T\); the variables of \(P_\eta\) are untouched.
On the smaller target side the \(b=2m\) positive source
columns each give one inverse determinant in the vacuum,
and the external exponent is \(m\). The negative-to-negative
block has positive Hom sign and its determinant contributes
\((\det k_-)^L\). Since \(L-m=T+p\), both final actions
have character \((\det k_-)^{T+p}\). These vacuum and
substitution rules are the same Fock normalization used in
Proposition~\ref{prop:secondcusp829}, with signs assigned
row by row. This proves equivariance for the full compact
group.

To prove irreducibility it suffices on the smaller side to
ignore the fixed factor in \(D\). A polynomial in \(Y\)
of degree \(bT\) which changes by \((\det u)^T\) under
\(Y\mapsto Yu\), \(u\in\operatorname{GL}_b(\mathbb C)\),
lies in a single irreducible \(U(N)\)-space. Indeed the
degree \(bT\) polynomial algebra on the two index sets has
the Cauchy decomposition into
\(S_\lambda(\mathbb C^N)\otimes S_\lambda(\mathbb C^b)\),
one summand per partition of \(bT\) with at most \(b\) rows,
in the index action used here. To see the multiplicity, take
the two tensor powers, apply Schur--Weyl duality on each
set of indices, and then take invariants under simultaneous
permutations. Two symmetric-group factors pair to the trivial
representation once when their partitions agree and never
otherwise. The determinant character selects only the
rectangle \(\lambda=(T,\ldots,T)\) with \(b\) entries.
The span of the products is a nonzero invariant subspace:
the repeated frame \(\eta_l=(1_b\ \ 0)\) gives the \(T\)-th
power of a nonzero minor of \(Y\). It is therefore the whole
irreducible space. The proof includes \(T=0\).

The two invariant Hermitian norms on an irreducible compact
type are proportional by Schur's lemma. Evaluate their
ratio at the repeated frame. There \(P_\eta\) is the \(T\)-th
power of the determinant of the last \(p\) columns of \(Z\),
and \(Q_\eta\) is the indicated power of \(\det D\) times
the \(T\)-th power of the determinant of the first \(b\)
rows of \(Y\). For an \(h\)-by-\(h\) matrix of independent
standard complex Gaussians the \(2t\)-th absolute determinant
moment is \(G_h(t)\). In fact Gram--Schmidt makes the squared
determinant a product of squared projected column lengths.
Conditional on the preceding columns the next length is the
sum of \(i\) independent exponential variables of mean one,
with \(i=h,h-1,\ldots,1\). Its conditional \(t\)-th moment
is \((t+i-1)!/(i-1)!\). Successive conditioning proves the
product formula. This Gaussian norm is exactly the stated
monomial norm, by one-variable integration. Independence of
the two square blocks in \(Q_\eta\) now gives the required
ratio.
\end{proof}

At every real place choose the repeated standard frame in
the fixed ordered positive coordinates of \(W\), and take
the tensor product over the places. Denote the resulting
vectors by \(P_*,Q_*\). Use the standard residue tests of
Section~\ref{sec:thirdstandard837} at all old and new marks.
The pair of adelic functions so obtained is \(V^0_{P_*}\) and
\(V^{\prime0}_{Q_*}\) on \([U(W)]\).

\begin{corollary}[A single pair of global functions for all choices]
\leavevmode\label{cor:thirdseed838}
Fix a field, a lift index and a twist. If a scalar
\(\delta_3(c)\) satisfies
\[
             \delta_3(c)V^0_{P_*}(y)=V^{\prime0}_{Q_*}(y)
               \quad\text{for every }y\in U(W)(\mathbb A),
\]
then for every prescribed \(i\) and \(g\)
\[
 \delta_3(c)L_{i,\vartheta}(g)=\mathcal S_{i,c}(g),\qquad
 \frac{\gamma_{12}(c)\delta_3(c)}{N_R}\,L_{i,\vartheta}(g)
                   =\mathcal B_i(g).
\]
Whenever the given node has a nonzero good form as in
Theorem~\ref{thm:complexseparation818}, the smaller standard
function \(V^{\prime0}_{Q_*}\) is nonzero. In that case a
scalar fulfilling the first equality of functions on \([U(W)]\)
is necessarily nonzero and unique.
\end{corollary}
\begin{proof}
The compact orbit of the repeated frame spans the irreducible
space in Theorem~\ref{thm:thirdreal838} at each place. Since
the map takes \(P_*\) to \(Q_*\) equivariantly, the same finite
linear combination of compact translates takes the two vectors
to any prescribed pair \(P_\eta,Q_\eta\). Tensor the finite
combinations over the real places. Translation of a real test
on the shared \(U(W)\) is right translation in \(y\) for
both global functions. All integrals exist by
Lemma~\ref{lem:rankdecay837} and compactness on the larger
side. Thus the standard identity passes to the actual
polynomials. Apply Theorem~\ref{thm:thirdstandard837} at
the finite places. Restrict to \(y_i(h',g)\) and sum with
\(C_c(h')/d_{h'}\). The two resulting sides are precisely
\(L_{i,\vartheta}\) and \(\mathcal S_{i,c}\) by their
definitions. The second identity follows from
Theorem~\ref{thm:smallthird834}.

If the smaller standard function vanished identically, its
real and finite translates would all vanish. The span and
translation formulas would force \(\mathcal S_{i,c}=0\)
for every \(i\), contradicting Corollary~\ref{cor:smallthirdseparation834}
in the presence of a nonzero good form. This proves the
nonvanishing assertion and the uniqueness statement.
\end{proof}

The remaining scalar problem here concerns the actual pair
\(V^0_{P_*},V^{\prime0}_{Q_*}\) on the rational quotient.
The compact norm ratio and the relative split factor in
Theorem~\ref{thm:thirdsplitprod836} have been computed on
their specified local data.
\clearpage

\subsection{The norm of the larger standard integral}
\label{sec:thirdlargenorm839}
On the common adelic quotient \([U(W)]\) of the third
exchange write \(F_c=V^0_{P_*}\) for the standard function
of Corollary~\ref{cor:thirdseed838}. Keep \(r=3m,\ n=p+r\),
\(e=(p-r)/2\) and
\(\lambda_c(a)=(\alpha^{-1}\upsilon_1^{n/2})(a/\bar a)\).
The determinant block in the marked matrix test has size
\(p\) on this side. Measures in this subsection are fixed
across all twists on the field at this lift index.

\begin{theorem}[Global norm of the larger standard theta integral]
\leavevmode\label{thm:thirdlargenorm839}
The function \(F_c\) is square integrable and nonzero on
\([U(W)]\). Up to a constant \(B_p>0\) independent of \(c\),
its squared norm is
\[
 \|F_c\|^2 =
 B_p\,
 \frac{\displaystyle\prod_{j=1}^{p}L_{E,f}(j-e,\lambda_c)}
      {\displaystyle\prod_{j=1}^{p}L_{F,f}(r+j,\eta^j)}
       \prod_{v\in\mathcal S_c}
                 \frac{P_p(q_v)}{H_p(q_v,r/2)} .
\]
In particular this formula concerns the actual theta integral
over the definite group with the standard marked tests.
\end{theorem}
\begin{proof}
Double the factor \(U(I_p)\) in the product of the kernel
and its conjugate, and integrate on \([U(W)]\). Using the
second member of a unitary pair as the doubled member gives
the same doubled theta construction: take the product of
the rational polarizations for the pair and its conjugate.
Their diagonal graph evaluation is the Hermitian pairing
on the test, now for the \(I_p\)-action. The determinant
correction on the opposite copy makes the two actions
contragredient there, as in Theorem~\ref{thm:normalfirst828}.
On the first copy the external multiplier at the argument
\(h\) is the actual target multiplier \(\alpha(\det h)\).
The two multipliers in the integrated variable cancel.
Equivalently one can transpose the two members with the
adjoint polarization and read the graph pairing back in
the original coordinates.

Here the Siegel rank is \(p\); the space integrated in the
doubled theta integral has dimension \(n\) and Witt index
\(w\le m\) by its real signature. It satisfies
\(n-w\ge p+2m>p\). The ordinary convergent theta integral
therefore gives the holomorphic value of the Siegel
Eisenstein series at \(s_0=(n-p)/2=r/2\), with factor one
in the measures of
\href{https://msp.org/pjm/2013/264-1/pjm-v264-n1-p11-s.pdf}
{Yamana, Theorem~2}. In particular the squared kernel,
for a fixed point of \(U(I_p)\), is integrable on \([U(W)]\).
The integral of its square is bounded over a compact set
of representatives of the definite quotient, by continuity
of the doubled theta integral. Cauchy--Schwarz on
\([U(I_p)]\), followed by integration in \(W\), proves
\(F_c\in L^2\). The same bound permits Fubini in its norm
on the three quotients.

Pair the restricted doubled Eisenstein family with the
constant on the two definite quotients. By compactness
evaluation at \(s_0\) commutes with the pairing. The resulting
value is a positive common measure factor times \(\|F_c\|^2\).
In the initial half-plane of absolute convergence every
rational maximal isotropic subspace of the double of \(I_p\)
is a graph, by definiteness. Thus the ordinary graph
unfolding gives the product of the flat graph integrals for
the constant representation of \(U(I_p)\), up to the fixed
diagonal volume.

At a marked place the row test is supported on
\((y,z)\in M_{p,r}(O_v)\times K_p\) with value \(\chi_v(\det z)\).
A translate by \(h^{-1}\) in the row basis has no overlap
unless \(h\in K_p\), in which case the external character
cancels the unit factor from the substitution. Its
coefficient and hence its integral are \(P_p(q_v)\), the
additive measure of this support, for every section parameter.
At an unmarked split place a target translate with elementary
divisor valuations \(b_i\) gives the absolute lattice overlap
\(q^{-n\sum_i|b_i|/2}\), one copy for each of the \(n\) columns.
The phase on the positive uniformizer is
\(\alpha_v(\varpi)\xi^{-n}=u_v^{-1}\). This follows by reversing
the source and target lattice substitutions as in
Theorem~\ref{thm:marksections835}. Moving from the theta point
\(s_0\) to \(s\) inserts \(q^{-(s-s_0)\sum_i|b_i|}\) by the
split graph height in Section~\ref{sec:squarefactors827}.
As \(n/2-s_0=p/2\), the split integral is precisely that
of Theorem~\ref{thm:splitzeta827} in dimension \(p\)
with phase \(u_v^{-1}\). Its two inverse phases give both
components of the Hecke character \(\lambda_c\), independent
of the ordering.

At an inert finite place all spaces and the section are
unramified with hyperspecial data. The diagonal graph has
the unimodular polarization of Lemma~\ref{lem:squarespherical828}.
The even-rank space character is unramified and trivial on
the common uniformizer, since its restriction to \(F_v^\times\)
is \(\eta_v^n=1\). The unramified doubling formula in that
lemma, in Siegel rank \(p\), has standard numerator for the
constant representation at \(s+1/2\) and denominator
\(\prod_{j=1}^p L_{F_v}(2s+j,\eta_v^j)\); both \(p,n\) are
even. Its spherical section has value one at the identity,
the unmarked vacuum norm. Thus the numerator arguments at
this place are \(s+j-p/2\) over \(E_w\). The same arguments
occur at the split places by the preceding computation.

At infinity \(U(I_p)\) is compact, with graph height one.
The polynomial \(P_*\) gains \((\det h)^T\) from its \(T\)
full row determinants. Its target vacuum and multiplier
have exponents \(-m\) and \(m-T\), respectively. They cancel
the polynomial exponent. Hence the real coefficient is
the positive squared norm of the real vector, constant on
the integrated graph and independent of \(s,c\).
The finite Euler product in the unfolding region thus has
numerator \(\prod_{j=1}^p L_{E,f}(s+j-p/2,\lambda_c)\), the
denominator above, and the correction \(P_p(q_v)/H_p(q_v,s)\)
at each mark. Equality of meromorphic continuations permits
evaluation at \(s_0\); individual local comparison factors
at interior points are not evaluated as divergent integrals.

At this point the denominator arguments \(r+j\) all exceed
one. The numerator arguments are the integers \(1-e,\ldots,r+e\).
For \(1-e,\ldots,e\) the gamma factors and the nonvanishing
of these same Hecke values were verified in
Theorem~\ref{thm:thirdsplitprod836}. At the remaining positive
integers the Euler product is absolutely convergent and
nonzero. Each marked correction is positive and finite.
All changes of global measure and all real coefficients are
positive and independent of the twist. The norm identity
and the nonvanishing claim follow.
\end{proof}
\clearpage

\subsection{The norm on the smaller side and the ratio of global norms}
\label{sec:thirdsmallnorm840}
Write \(F'_c=V^{\prime0}_{Q_*}\) on \([U(W)]\), with source
\(G=U(S_0)\) of rank \(r\). Its integral in the source
variable is absolutely convergent by Lemma~\ref{lem:rankdecay837}.
To compute its square norm we shall initially cut off the
integration on \([G]\).

\begin{lemma}[Exclusion of the boundary orbits in the smaller norm]
\leavevmode\label{lem:nongraph840}
In the doubled norm construction with source \(S_0\), the
section determined by the standard marked test vanishes on
every non-graph rational orbit of the two source groups.
Its theta-point coefficient on the graph is integrable in
absolute value on \(G(\mathbb A)\).
\end{lemma}
\begin{proof}
At a chosen marked split place use the two position coordinates
\(X_1,X_2\in M_{n,r}(F_v)\) for the test and its conjugate.
Their top \(r\)-by-\(r\) blocks are units, so both matrices
are injective on their respective source spaces. In the split
position model the double acts by invertible substitutions
on the \(2r\) columns, with the usual scalar determinant
factor. The diagonal evaluation functional integrates the
test on \(X_1=X_2\), as in Section~\ref{sec:squarefactors827}.
In intrinsic terms the variable \((X_1,-X_2)\), considered
as a map on the doubled space, vanishes on the diagonal
\(r\)-plane. Thus the section at a general doubled element
integrates over maps vanishing on the transported \(r\)-plane.
One sees this directly by applying the column substitution
before diagonal evaluation; a change of Haar on the linear
fiber only multiplies by a scalar. At a rational non-graph
orbit this plane has nonzero intersection with a source
summand, as proved in Theorem~\ref{thm:carrierline831}.
This intersection remains nonzero on the chosen split
component. Vanishing on it is impossible for the corresponding
\(X_i\). Translation by the two source groups preserves that
obstruction. The flat continuation multiplies the fiber
value by a height, so the same vanishing persists.

At a mark the graph overlap is compactly supported on
\(K_r\), by the unit top block. At unmarked finite places the
vacuum overlap bound and the self-dual lattice count of
Lemma~\ref{lem:selfdualtail813} apply with source rank \(r\)
and target dimension \(n=p+r\ge3r\) at our nodes. Indeed
increasing the target dimension only increases the exponent
of the intersection index in the denominator. Thus the
absolute integral outside the full compact subgroup at an
unmarked \(v\) is \(O(q_v^{-r/2})\). The compact contribution
is one, and the sum of the errors over the fixed number
field is finite.

At a real place retain \(b=2m,\ T=\mu_\tau,\ L=T+n-2m\)
from Theorem~\ref{thm:thirdreal838}. The maximal compact
group of the source fixes \(Q_*\) including the vacuum and
splitting factors. In fact on the positive source block
the exponents from vacuum, external character and determinant
polynomial are \(m,T-m,-T\), and on the negative block they
are \(p+2m,T-m,-L\). Both sums are zero. Every monomial of
the polynomial involves only letters of positive Hom sign.
For each paired source-column boost the lowering operator
therefore kills the vector and the nonconstant creation
terms pair to zero against it, as in
Lemma~\ref{lem:oscillatorcartan710}. The vacuum contributes
one inverse hyperbolic cosine per target row. In addition
every positive source column occurs to degree \(T\) in the
positive determinant powers and every negative source column
to degree \(L\). The normalized coefficient on a Cartan
element with parameters \(t_i\ge0\) is consequently
\[
                  \prod_{i=1}^m(\cosh t_i)^{-(n+T+L)}.
\]
The compact invariance just proved makes the coefficient
positive on the entire real group. The Cartan density is
bounded by a constant times \(\exp(2(r-1)\sum_i t_i)\).
Since \(n>2r\) on the prescribed fields, the real integral
is finite and strictly positive. Together the local estimates
prove the absolute integrability on \(G(\mathbb A)\).
\end{proof}

\begin{theorem}[Global norm of the smaller standard theta integral]
\leavevmode\label{thm:thirdsmallnorm840}
The function \(F'_c\) is square integrable and nonzero.
There is a constant \(B_r>0\) independent of \(c\) on the
given field and lift index such that
\[
 \|F'_c\|^2=
 B_r\,
 \frac{\displaystyle\prod_{j=1}^{r} L_{E,f}(e+j,\lambda_c)}
      {\displaystyle\prod_{j=1}^{r} L_{F,f}(p+j,\eta^j)}
       \prod_{v\in\mathcal S_c}
                     \frac{P_r(q_v)}{H_r(q_v,p/2)} .
\]
\end{theorem}
\begin{proof}
Let \(u\) be smooth and compactly supported on \([G]\).
Insert it in the integral defining \(F'_c\) and denote the
result by \(T_c u\). The doubled Siegel rank is now \(r\).
The space integrated in the Siegel--Weil formula has
dimension \(n\) and Witt index at most \(m\), so
\(n-m>r\). The ordinary convergent formula in
\href{https://msp.org/pjm/2013/264-1/pjm-v264-n1-p11-s.pdf}
{Yamana, Theorem~2} applies at \(s_1=(n-r)/2=p/2\).
In fact \(n>2r\), so the Eisenstein series is in the
absolute convergence range already at the theta point.
As in Theorem~\ref{thm:thirdlargenorm839}, the integrated
square of the fixed-point kernel is finite and locally
bounded in its source point by the convergent formula.
Cauchy--Schwarz and the compact support of \(u\) justify
Fubini for the norm of \(T_cu\).

Unfold the doubled series against \(u\) and its conjugate
at this point. On the compact support absolute locally
uniform convergence permits termwise integration.
Lemma~\ref{lem:nongraph840} kills every non-graph term
before integrating in a unipotent variable. On the graph
the integrated diagonal variable gives the ordinary matrix
coefficient of \(u\), up to complex conjugation according
to the ordering of the copies. In particular
\[
    \|T_cu\|^2\le C\,\|u\|^2\int_{G(\mathbb A)}|K_c(a)|\,da,
\]
where \(K_c\) is the actual theta-point graph coefficient
and \(C>0\) is a Haar constant independent of \(c\). Here
we used unitarity of right translation on the rational
quotient. This proves boundedness on the dense compactly
supported space, since the last integral is finite.

Choose smooth compact cutoffs \(0\le u_j\le1\) tending to
one pointwise and in \(L^2([G])\). Such cutoffs come from an
exhaustion of this finite-volume quotient, or from compact
exhaustions at the fixed finite level. Their images converge
in \(L^2\) by the bound. At each \(y\in U(W)(\mathbb A)\)
they also converge to \(F'_c(y)\) by the absolute convergence
in Lemma~\ref{lem:rankdecay837}. An almost everywhere
convergent subsequence of the \(L^2\) sequence identifies
its limit with \(F'_c\). Moreover the source matrix
coefficients of \(u_j\) converge uniformly to the coefficient
of the constant, namely \(\operatorname{vol}([G])\).
This follows by Cauchy--Schwarz and \(L^2\) convergence on
both sides of each unitary translation. Integrating against
the \(L^1\) graph coefficient therefore gives the ordinary
product of local graph integrals for the squared norm of
\(F'_c\), up to positive fixed Haar factors. The cutoff
argument accounts for the noncompact source in this calculation.

At a mark the test is supported on
\(K_r\times M_{p,r}(O_v)\) in its top and bottom blocks.
Two translates in the source have disjoint support unless
the quotient belongs to \(K_r\). There the test is fixed
under the combined action, with norm \(P_r(q_v)\).
At an unmarked split place the source lattice coefficient
has phase \(u_v=\alpha_v(\varpi)^{-1}\xi^n\) at determinant
valuation one, and modulus \(q^{-n\ell(a)/2}\) in elementary
divisor notation. Apply Theorem~\ref{thm:splitzeta827} in
dimension \(r\), with parameter \(p/2\) since \(n=r+p\).
Its two Euler components are those of \(\lambda_c\), and
the arguments are \(p/2+j-r/2=e+j\). At inert places the
spherical unitary doubling computation in
Lemma~\ref{lem:squarespherical828} applies in rank \(r\)
for the even space character of rank \(n\), trivial on the
common uniformizer. With the constant input and section of
value one its numerator has the same \(e+j\) arguments, and
its denominator has arguments \(p+j\) and characters \(\eta^j\).
At the marks the correction displayed in the theorem gives
exactly \(P_r(q_v)\) after cancellation of the local
reciprocal denominator. The real factors were proved positive,
finite and independent of the twist in
Lemma~\ref{lem:nongraph840}. This gives the formula. Its
finite Hecke functions are taken at real arguments exceeding
one, hence are nonzero. The norm is therefore nonzero.
\end{proof}

\begin{corollary}[The third multiplier as a ratio of actual global norms]
\leavevmode\label{cor:thirdnormratio840}
For a positive \(K_3\) independent of \(c\) at the given
node and lift index one has
\[
              \frac{\|F'_c\|^2}{\|F_c\|^2}
                    =\frac{K_3}{\mathcal D_e(\sigma_c)} .
\]
In particular if the two functions are proportional, the
squared modulus of their unique proportionality factor
in Corollary~\ref{cor:thirdseed838} is the displayed value.
\end{corollary}
\begin{proof}
Divide the large norm by the small one. The arguments
\(e+1,\ldots,e+r\) occur in both numerators and cancel.
The remaining numerator values are exactly \(1-e,\ldots,e\).
The denominators of the two global products are independent
of \(c\). At a marked place the relative correction is one
since
\[
       P_p(q) H_r(q,p/2)
           = \prod_{i=1}^{p+r}(1-q^{-i})
           = P_r(q) H_p(q,r/2).
\]
Upon division by the corresponding ratio for a reference
twist, the paired integer values give
\(\mathcal D_e(\sigma_c)/\mathcal D_e(\sigma_{c_*})\)
by Theorem~\ref{thm:thirdsplitprod836}. Inverting and fixing
the positive normalization at \(c_*\) gives \(K_3\).
\end{proof}

Thus the test still needed on the common rational quotient is
whether the two nonzero standard functions span the same line.
Equivalently their mixed inner product must attain the
Cauchy--Schwarz bound for the two norms just evaluated.
\clearpage

\subsection{Holomorphy of the two standard functions on the common group}
\label{sec:thirdholomorphic841}
Use the compact type in Theorem~\ref{thm:thirdreal838} and
the two nonzero square-integrable functions \(F_c,F'_c\)
of Section~\ref{sec:thirdsmallnorm840}. Fix a real embedding.
Write \(i\) for a positive coordinate of \(W\), \(\beta\)
for a negative one, \(a\) for a positive coordinate of
\(S_0\), and \(d\) for a negative one. The holomorphic
equations in the positive-first order use the root operators
on \(U(W)\) of compact torus weight \(\epsilon_i-\epsilon_\beta\);
equivalently these are the contractions on the oscillator
with \(W\) as source and a positive definite target.

\begin{theorem}[Holomorphy after integration on the indefinite source]
\leavevmode\label{thm:thirdholomorphic841}
Both standard functions are killed by all these operators
at every real place. The assertion holds on their entire
matched compact types, including the arbitrary row choices
of Theorem~\ref{thm:thirdreal838}.
\end{theorem}
\begin{proof}
On the larger side a noncompact lowering operator is a sum
of products of derivatives in a positive and a negative
source column, one term per target row. The polynomial
\(P_\eta\) is independent of every negative column. This
proves the equation before integration.

On the smaller side retain the positive-sign Fock blocks
\(Y_{ia},D_{\beta d}\) and denote the opposite-sign blocks
by \(Z_{id},V_{\beta a}\). Under the compact torus of
the target \(W\), the four blocks \(Y,D,Z,V\) have weights
\(\epsilon_i,\epsilon_\beta,-\epsilon_i,-\epsilon_\beta\),
respectively. Under the source torus the corresponding
weights are \(-\epsilon_a,-\epsilon_d,\epsilon_d,\epsilon_a\).
Here the letter is dualized on the opposite-sign Hom
coordinate. These rules follow by diagonal substitution on
both indices; the vacuum weights are independent of the
monomial.

We record the two infinitesimal actions used below. On a
polynomial \(H(Y,D)\) involving only these two blocks, the
root operator of weight \(\epsilon_i-\epsilon_\beta\) on
\(W\) acts by a fixed nonzero scalar times
\[
                         \sum_a Y_{ia}V_{\beta a}H .
\]
On such a polynomial the root operator of weight
\(\epsilon_a-\epsilon_d\) on \(S_0\) acts by a nonzero
scalar times
\[
                         \sum_{\beta'} V_{\beta'a}D_{\beta'd}H .
\]
To verify the formulas, a matrix unit of either unitary
Lie algebra switches the two indicated coordinates,
and acts separately on each fixed index of the other
space. The two modes in each switch have opposite Hom
sign. The two off-diagonal complexified symplectic operators
on these modes are multiplication by their product and
differentiation by their mixed second derivative. This
follows directly from the Heisenberg commutators
\([\partial_x,x]=[\partial_y,y]=1\): multiplication by
\(xy\) sends each annihilation operator in its commutator
to the other creation operator up to sign, and
\(\partial_x\partial_y\) performs the opposite switch.
Their compact weights distinguish the two root directions.
There is no scalar term for a nonzero torus weight. For the
target root the other terms are derivatives in \(Z_{id}\)
and \(D_{\beta d}\), and for the source root they are
derivatives in \(Y_{ia}\) and \(Z_{id}\). They kill \(H\).
The creation coefficients are constant on each indicated
sum by permutation of like-signed coordinates, and nonzero
because the matrix unit acts nontrivially on each pair.
Changes of root-vector normalizations can be absorbed in
the scalar coefficients. The determinant splittings have
zero differential on these off-diagonal operators.

Write \(P(Y)=\prod_l\det(\eta_lY)\) and
\(Q=(\det D)^L P(Y)\), with \(L=T+n-2m\ge1\).
For the fixed \(i,\beta\), use the polynomials
\[
       H_{ad}=Y_{ia}\,(\operatorname{adj}D)_{d\beta}
                            (\det D)^{L-1}P(Y).
\]
They depend only on \(Y,D\). The identity
\(D\operatorname{adj}D=(\det D)1_m\) gives exactly
\[
       \sum_{a,d}\sum_{\beta'} V_{\beta'a}D_{\beta'd}H_{ad}
                    =\sum_a Y_{ia} V_{\beta a}Q .
\]
Consequently the unwanted target derivative of \(Q\) is
a linear combination of source derivatives of the \(H_{ad}\).
This identity holds in the polynomial algebra even on the
singular locus of \(D\).

Pass to rational theta sums. A source derivative on the
real test gives the corresponding right derivative in the
variable integrated on \([G]\). Its integral is zero: right
translation preserves quotient Haar measure, and the
difference quotient can be integrated by the uniform rapid
decay of Lemma~\ref{lem:rankdecay837} for compact families
of source translations and Schwartz derivatives. The same
lemma justifies taking the target derivative under the
integral at every target point. Thus all the operators
annihilate the smaller function as asserted. The argument
works for each \(\eta\) and for their linear combinations.
\end{proof}

\begin{corollary}[Real matrix coefficients of the two global functions]
\leavevmode\label{cor:thirdrealmatrix841}
There is one function \(M_\infty\) on \(U(W)(F_\infty)\),
normalized by \(M_\infty(1)=1\), such that for any
\(f,h\in\{F_c,F'_c\}\) at the given twist,
\[
                 \langle R(a)f,h\rangle
                        =M_\infty(a)\langle f,h\rangle
                   \qquad (a\in U(W)(F_\infty)).
\]
The normalization on the common repeated frame is as in
Corollary~\ref{cor:thirdseed838}.
\end{corollary}
\begin{proof}
The evaluations of the real polynomial vectors define compact
equivariant maps from one irreducible compact type into
\(L^2\), using Theorem~\ref{thm:thirdreal838} to identify the
two domains. Both maps are nonzero by
Theorems~\ref{thm:thirdlargenorm839} and~\ref{thm:thirdsmallnorm840}.
On this finite-dimensional type every pairing obtained from
any two of these maps is a scalar multiple of one invariant
Hermitian form, by Schur's lemma. This also allows zero for
the mixed scalar.

For passage to the real group, the polynomial vectors used
in the pairings are smooth for the Hilbert-space real action.
On the larger side any derived real polynomial has square
integrable image by the convergent Siegel--Weil argument
with the compact source in Theorem~\ref{thm:thirdlargenorm839}.
On the smaller side the estimate used for the cutoff in
Theorem~\ref{thm:thirdsmallnorm840} applies also after a
fixed archimedean differential operator. Such an operator
gives a polynomial times the Gaussian. For any fixed finite
degree the absolute Cartan coefficient on the integrated
source is \(O(\prod_i(\cosh t_i)^{-n})\). Indeed the lowering,
diagonal and raising factorization of
Lemma~\ref{lem:oscillatorcartan710} first reduces degree a
finite number of times, multiplies each monomial by the
vacuum factor and powers of \(\operatorname{sech} t_i\),
and only finitely many subsequent raising terms can pair
with the fixed degree on the other side. The remaining
coefficients involve bounded powers of \(\tanh t_i\), also
uniformly under compact rotations. This bound is integrable
since \(n>2r\). The marked rank support is unchanged by
these real operators. Thus all distributional derivatives
are in \(L^2\), giving smooth vectors by the usual Sobolev
criterion for a unitary Lie group action.

Order complexified real enveloping-algebra monomials with
holomorphic raising operators on the left, compact operators
in the middle and lowering operators on the right. As in
Lemma~\ref{lem:carriermatrix831}, a lowering factor kills
the vector and a left raising factor kills the pairing by
unitary adjunction, using Theorem~\ref{thm:thirdholomorphic841}
for both entries. All derivatives at the identity reduce
to compact derivatives, already determined up to the scalar
of the invariant pairing. The matrix coefficients are real
analytic. To check this here, take a compact-invariant
orthonormal basis of the real tangent space in the sense of
an invariant inner product, with the noncompact directions
in pairs \(X,JX\). On a holomorphic vector their sum of
squares reduces to a sum of commutators in the compact
algebra as in Lemma~\ref{lem:carriermatrix831}. This sum is
compact-invariant and acts on the irreducible compact type
by a scalar, as does the compact sum of squares. The
coefficient therefore satisfies an analytic elliptic
eigenfunction equation. The same argument applies to the
product over real places. Equality of derivatives gives
equality on the connected real group. Take \(M_\infty\)
from the normalized self-pairing of \(F_c\) to conclude.
\end{proof}

The mixed scalar in this corollary is the global inner
product at \(a=1\). Identifying the real coefficients fixes
their dependence on real translations; the absolute value
of that mixed scalar is still needed for proportionality.
\clearpage

\subsection{The local representations and the standard adelic vector}
\label{sec:thirdadelic842}
Continue with the two functions \(F_c,F'_c\) on \([U(W)]\).
At a split place, define two unitary characters of \(F_v^\times\)
by
\[
       \beta_B(z)=\nu_v(z)^{-1}\xi^{p v(z)},\qquad
       \beta_S(z)=\kappa_v(z)\xi^{-r v(z)} .
\]
Thus \(\beta_B/\beta_S\) is the unitary part \(u_0\) of
\(\omega_0\) in Section~\ref{sec:marksections835}. We use
the frame and the ordering of its two blocks fixed there.

\begin{lemma}[The marked principal series and its specified vector]
\leavevmode\label{lem:thirdmarkedrep842}
At a mark consider normalized induction to
\(\operatorname{GL}_n(F_v)\) from the upper block parabolic
of type \((r,p)\) of
\((\beta_B\circ\det_r)\boxtimes(\beta_S\circ\det_p)\).
It is irreducible. On either standard global integral, varying
the local full-rank test induces a map from this same
representation, with the same nonzero local vector for the
standard pair.
\end{lemma}
\begin{proof}
Use the model with translations on the left and parabolic
covariance on the right. For \(a_0(t)=\beta_S(\det t)
|\det t|^{-r/2}\), the two functions of the frame are
\[
   a_0(t)^{-1}J_r(\Psi;t,0),\qquad
   a_0(t)^{-1}\omega_0(\det t)^{-1}J_p(\Phi;t,0).
\]
Their left action is pure translation by
Theorem~\ref{thm:marksections835}. On right multiplication
by a parabolic element with Levi \((A,D)\), their common
covariance is
\[
 \delta_P(A,D)^{-1/2}
             \beta_B(\det A)^{-1}\beta_S(\det D)^{-1},
 \qquad \delta_P(A,D)=|\det A|^p|\det D|^{-r}.
\]
Indeed the column integral transforms by
\(\omega_0(\det A)^{-1}\) and one must also divide by the
value of \(a_0\) on the parabolic. The row integral gives
the same result by the row change of variables. On the
standard pair the two sections coincide, with value one at
the identity.

Both frame maps are onto the smooth induced space. On a
Grassmannian coordinate chart trivialize the bases using
an identity minor. A local section of compact open support
is lifted by taking its scalar function of the plane times
\({\bf1}_{K_b}(h)\omega_0(\det h)^{-1}\) in the basis
coordinate \(h\), where \(b=p\) or \(r\), incorporating
the fixed displayed scalar for the chosen frame. Its fiber
integral has the prescribed value; compactly many charts
and a disjoint finite refinement suffice. By
Proposition~\ref{prop:framecoinvariants835} the kernel of
each map consists of sums of the weighted changes of basis
minus the original vector. Each such difference pairs to
zero in the global integral by right invariance of the Haar
measure of the integrated group. The bigger integral is
over a compact quotient; the smaller integral with these
full-rank tests converges by Lemma~\ref{lem:rankdecay837}.
This proves factorization and equivariance of the global maps.

For irreducibility one may apply the segment criterion of
\href{https://www.numdam.org/article/ASENS_1980_4_13_2_165_0.pdf}
{Zelevinsky, \emph{Induced representations of reductive
\(p\)-adic groups II}, Theorem~4.2}. In his segment
subrepresentation convention each determinant character in
rank \(b\) is the segment subrepresentation of the increasing
sequence from \(\beta|\cdot|^{-(b-1)/2}\) to
\(\beta|\cdot|^{(b-1)/2}\). The constant vector embeds in
that normalized Borel induction after the determinant twist.
The two segments here are on disjoint cuspidal lines since
\(\beta_B/\beta_S\) is nontrivial on the residue units.
They cannot be linked. The criterion gives irreducibility
of their product. In particular, since the common standard
vector has nonzero global image on both sides, each image
of this local induced module is a copy of it.
\end{proof}

\begin{theorem}[A common adelic representation with fixed local vectors]
\leavevmode\label{thm:thirdadelic842}
For a fixed twist every nonzero irreducible projection of
\(F_c\) or \(F'_c\) in \(L^2([U(W)])\) has the same abstract
adelic representation \(\pi\). In each such copy the
projected vector lies on the same abstract tensor line.
At the marks use the common vector in
Lemma~\ref{lem:thirdmarkedrep842}. At an unmarked finite
place the vector is spherical. Its standard multiset is
\[
 \beta_B(\varpi)A_r(q_v)\sqcup\beta_S(\varpi)A_p(q_v)
       \ \text{at a split place},\qquad
 A_r(q_w)\sqcup A_p(q_w)
       \ \text{at an inert place}.
\]
At the real places take the common holomorphic representation
and repeated-frame vector determined in
Theorem~\ref{thm:thirdholomorphic841}.
\end{theorem}
\begin{proof}
First the two vectors lie in the discrete spectrum.
They are square integrable by the preceding norm formulas.
Their compact types are finite dimensional and annihilated
by the holomorphic lowering operators. They have a fixed
infinitesimal character: the repeated frame spans an
extremal line for the compact roots in the rectangular
representation of Theorem~\ref{thm:thirdreal838}. Choose the
compact positive roots for which this is a highest line and
adjoin the roots of the annihilating noncompact subspace.
These form a positive system, obtained by putting all
positive coordinates before all negative ones. The vector
is a highest-weight vector for that system, so the center
acts by the infinitesimal character of its highest weight.
In the Langlands spectral decomposition the continuous part
is a sum of unitary inductions integrated over imaginary
split Levi parameters for proper parabolics. On each such
parameter space the real Casimir eigenvalue is a nonconstant
quadratic polynomial in the continuous parameter, by the
infinitesimal-character formula for normalized induction.
The norm-square term is definite on the imaginary split
parameters. The locus of any specified eigenvalue has
Lebesgue measure zero. The discrete data in the decomposition
are countable. Thus no vector of this fixed infinitesimal
character has a continuous part. This applies to the present
unitary group with anisotropic center.

At infinity each nonzero irreducible projection with the
specified vector has a fixed highest weight. The simple
highest-weight module is the unique irreducible quotient
of its Verma module, so the real constituent is determined.
The highest-weight space in the simple module is a line.
The repeated frame supplies that vector in both forms;
equivalently the normalized real coefficients were identified
in Corollary~\ref{cor:thirdrealmatrix841}. The unitary
globalization is unique since inner products on the cyclic
module are determined by the norm of its highest line and
unitary adjunction. This argument works on each real group
separately.

At a marked place project the equivariant maps of
Lemma~\ref{lem:thirdmarkedrep842} onto any nonzero
automorphic summand. In the tensor decomposition of that
summand, contraction against any smooth vector on the other
places gives a local intertwiner, nonzero for at least one
contraction. Irreducibility and Schur's lemma show that the
summand has the indicated local representation and that
all such contractions of the standard vector lie on the
same specified line. This also proves the tensor-line
assertion there.

At an unmarked split place the smaller rectangular lift
is from the integrated constant in rank \(r\). The local
Grassmannian computation of Lemma~\ref{lem:secondhecke829}
works before integration and with arbitrary ranks. To
recall its phases here, the positive elementary target
coset of size \(h\) has phase \(\beta_S(\varpi)^h\) and
absolute multiplier \(q^{rh/2}\); the source coset with
\(l\) negative elementary divisors has phase \(u_0(\varpi)^{-l}\)
and absolute multiplier \(q^{nl/2}\). The lattice count
for residue rank \(a\) expands
\({n-a\brack h}_q\) using \({r-a\brack l}_q\) with coefficient
\({n-r\brack h-l}_q q^{l(n-r-h+l)}\). Every source translate
has the same integral against the constant. Dividing by the
two averaging masses gives, by the elementary symmetric
expansion in Theorem~\ref{thm:allcarriers830}, the segments
\(\beta_S A_{n-r}\) and \(\beta_S u_0 A_r\).
Their phases are as asserted.

For the larger rectangular integral transpose the matrix
and replace each group variable by its inverse transpose.
The rank \(n\) target in these coordinates has phase
\(\beta_B^{-1}\) and inverse square-root Jacobian on its
\(p\) columns. The integrated rank \(p\) source has phase
\(u_0\) and direct Jacobian, since its former action had
multiplier \(\omega_0^{-1}\). The same count now gives
\(\beta_B^{-1}A_{n-p}\) and \(\beta_B^{-1}u_0 A_p\).
Pullback by the inverse transpose inverts the multiset.
Since the segments \(A_b\) are stable under inversion and
\(\beta_B/u_0=\beta_S\), this agrees with the smaller
integral. The proof is valid for arbitrary local unitary
phases, even when some entries coincide.

At an inert place apply Theorem~\ref{thm:inert175} to the
spherical joint local quotient of the oscillator. We give
the justification of the quotient in this application where
one of the constants is on a noncompact group. With all
other components fixed, vary the Schwartz test at this
inert place. The integrated variable still acts trivially
after integration by Haar invariance. The larger-integral
map remains \(L^2\)-valued by the convergence argument of
Theorem~\ref{thm:thirdlargenorm839}. For the smaller-integral
map the other places include a split full-rank mark, so
both rapid source convergence and the exclusion of boundary
orbits survive. The local coefficient at the varying inert
place for any fixed smooth test satisfies the same lattice
overlap bound up to a test-dependent constant. Indeed in a
self-dual lattice model a finite Schwartz vector is a finite
linear combination of Heisenberg translates of the vacuum.
An overlap of any two such translates is bounded by the
same square-root intersection-index bound: the intersection
of two translated lattices is empty or a coset of their
intersection, and the extra phase is unitary. The proof of
Lemma~\ref{lem:nongraph840} therefore gives the needed local
\(L^1\) estimate and the cutoff argument of
Theorem~\ref{thm:thirdsmallnorm840} applies. Projection and
contraction in the other places now give a nonzero local
map with trivial integrated member and the specified
spherical output. Dual conventions do not change the
reciprocal multiset here. Every inert place is odd and the
spaces are unramified quasi-split with the given full
stabilizers. Their dimensions \(n,r,p\) are even. Thus a
lift from the constant in rank \(r\) adjoins \(A_p\) to
\(A_r\), and a lift from the constant in rank \(p\) adjoins
\(A_r\) to \(A_p\), with all uniformizer phases one.

At each unmarked place the multiset determines the spherical
representation and its compact-fixed line has dimension one.
A discrete irreducible automorphic representation is the
restricted tensor product of its local constituents. The
conditions just proved fix all these constituents. They
also fix the one vector line in each factor: at a mark this
comes from the common vector in the induced model, regardless
of the dimension of its fixed space at a principal level.
Consequently in any abstract copy all projections lie on
the same tensor line. This proves the theorem, allowing at
this stage several automorphic copies of the representation.
\end{proof}
\clearpage

\subsection{Multiplicity one and the third theta comparison}
\label{sec:thirdglobal843}
Fix a sufficiently large node, a lift index \(j\), and one of its
prescribed twists \(c=\chi_M^j\vartheta\). Write \(U=U(W)\),
\(r=3m\), \(p=\dim W'\), and \(n=p+r\). Both \(p\) and \(r\) are
even and \(p>r\). Theorem~\ref{thm:thirdadelic842} identifies an
irreducible adelic representation \(\pi\) of \(U\) and one line of
vectors in its abstract tensor product. To compare functions on the
automorphic quotient we must count copies of \(\pi\) there.

\begin{proposition}[Automorphic multiplicity of the common representation]
\leavevmode\label{prop:thirdmultiplicity843}
The representation \(\pi\) of Theorem~\ref{thm:thirdadelic842}
has multiplicity one in \(L^2_{\rm disc}([U])\).
\end{proposition}
\begin{proof}
Let \(\mu\) be the fixed unitary Weil splitting character of
\(\mathbb A_E^\times/E^\times\) in dimension one, with
\(\mu|_{\mathbb A_F^\times}=\eta_{E/F}\). Its restriction to
norm-one ideles is \(\upsilon_1\). Define two Hecke characters of \(E\) by
\[
 B(a)=\nu(a/\bar a)^{-1}\mu(a)^p,\qquad
 S(a)=\kappa(a/\bar a)\mu(a)^{-r}.
\]
These characters are unitary. They are conjugate self-dual because
\(\mu(a\bar a)=\eta_{E/F}(N_{E/F}a)=1\), and are trivial on
\(\mathbb A_F^\times\) because \(p,r\) are even. At a good split
place, take the idele whose coordinates above that place are
\((z,1)\). Its quotient by its conjugate has coordinates
\((z,z^{-1})\), so \(B,S\) on the chosen factor are the characters
\(\beta_B,\beta_S\) of Theorem~\ref{thm:thirdadelic842}. At a
good inert place the uniformizer can be taken in \(F_v\), so both
characters have value one on it.

For a unitary character \(\lambda\), write \(\lambda[b]\) for its
tensor product with the \(b\)-dimensional irreducible representation
of the Arthur \(\mathrm{SL}_2(\mathbb C)\). Its standard unramified
parameter contributes the \(b\) twists with exponents
\((b-1)/2,(b-3)/2,\ldots,(1-b)/2\). The preceding computation and
Theorem~\ref{thm:thirdadelic842} identify the standard parameters
of \(\pi\) at almost all places with those of
\(B[r]\boxplus S[p]\), including both places of \(E\) over a split
place by conjugate duality.

Take \(R=2n\) and \(M=2R+1\). Choose the quasi-split odd unitary
group \(H\) over \(F\) of dimension \(M\), given by \(R\) hyperbolic
planes and a line, with the opposite Hermitian sign convention for
the dual pair. Use a nontrivial adelic additive character and
splitting characters \(1,\mu\) for the spaces of dimensions \(n,M\),
respectively. Their restrictions to \(\mathbb A_F^\times\) are
\(\eta_{E/F}^{\,n}\) and \(\eta_{E/F}^{\,M}\), as required.
Write \(\sigma\) for the restricted tensor product of the local
small theta lifts of \(\pi_v\) to \(H(F_v)\). Since \(n<R\),
Li's stable-range inequality gives
\[
 1\ \le\ m_{\rm disc,U}(\pi)\ \le\ m_{\rm disc,H}(\sigma).
\]
The first inequality follows also from the nonzero discrete
vectors in Theorems~\ref{thm:thirdlargenorm839},
\ref{thm:thirdsmallnorm840}, and~\ref{thm:thirdadelic842}.
The stable-range result applies to the discrete spectrum including
its residual part. We use it in the form stated by
\href{https://arxiv.org/pdf/2103.07956}
{Chen and Zou, Theorem~3.1, after J.-S.~Li, Theorem~A}.

We determine the packet on \(H\) using its standard parameters.
In the unramified theta correspondence from dimension \(n\) to
dimension \(M\), the input parameter is twisted by the inverse
of the large splitting character times the small one, and a
string of length \(M-n\) with the small splitting character is
adjoined. This is formula~(4.2) of
\href{https://arxiv.org/pdf/2103.07956}{Chen and Zou}.
Thus the unramified data of \(\sigma\) are those of
\[
 \widetilde\psi=(B\mu^{-1})[r]\ \boxplus\
                 (S\mu^{-1})[p]\ \boxplus\ {\bf1}[M-n].
\]
This is an elliptic global parameter for \(H\), with ordinary
base change. For the first two characters the restriction to the
ideles of \(F\) is \(\eta_{E/F}\). A conjugate self-dual character
of this restriction has conjugate-symplectic sign. Tensoring it
with an even-dimensional irreducible of \(\mathrm{SL}_2\), whose
invariant form is alternating, gives conjugate-orthogonal sign.
The trivial character and the odd-dimensional last factor have
positive sign as well. This is the required parity in odd dimension.
The lengths are distinct since \(M-n=3n+1>p>r\). Thus no simple
summand is repeated.

The quasi-split discrete classification assigns one global parameter
to the near-equivalence class of \(\sigma\). To identify it, take
standard base change at all good places of \(E\). The isobaric
cuspidal data for the assigned parameter agree there with the
three strings above. Strong multiplicity one for isobaric
representations of general linear groups identifies the multisets of
cuspidal constituents with their real twists. These multisets
determine the Arthur blocks. Indeed group the constituents by
their unitary cuspidal part. The largest exponent in each group
is \((b-1)/2\) for its longest block, with multiplicity the number
of such blocks. Remove their full strings and iterate. Hence the
assigned parameter is \(\widetilde\psi\).

For this parameter we check the local multiplicities entering the
global packet. At a split place the group is general linear and
its packet has one member. At a nonsplit finite place the three
Weil factors are unitary characters and the Deligne
\(\mathrm{SL}_2\) acts trivially. Interchanging the two
\(\mathrm{SL}_2\)'s therefore produces a tempered parameter.
Its local packet has no repeated representation. The co-tempered
packet is obtained by Aubert duality on those representations,
with the component-group labels adjusted in the character identity.
Aubert duality is an involution on irreducibles up to sign, so
distinct underlying members stay distinct. This uses
\href{https://arxiv.org/pdf/2410.13504}
{Atobe, Gan, Ichino, Kaletha, M\'inguez and Shin, Theorem~5.4.1},
with its tempered character identities on the quasi-split group.
At each real place the multiplicity-one theorem for real unitary
Arthur packets applies, including singular infinitesimal characters.
See \href{https://webusers.imj-prg.fr/~colette.moeglin/MR7.pdf}
{M{\oe}glin and Renard, Theorem~1.3}.

In the global formula each restricted tensor product of packet
entries satisfying the global character condition contributes one
summand. The local facts imply that the underlying representation
\(\sigma\) can occur in at most one such tensor product, since its
local representations are determined at every place. We have
already identified its unique global parameter. This gives
\(m_{\rm disc,H}(\sigma)\le1\), proving the assertion by Li's
inequality. Here the global input is
\href{https://arxiv.org/pdf/1206.0882}
{Mok's quasi-split multiplicity formula, Theorem~2.5.2}.
For its auxiliary intertwining and trace formula inputs we use
\href{https://arxiv.org/pdf/2410.13504}
{Atobe and collaborators, Introduction and main theorems}, and
the twisted weighted fundamental lemma in the form proved in
\href{https://arxiv.org/abs/2611.26750}
{Halleck-Dub\'e's November 2026 preprint, \S1.14}.
The auxiliary group \(H\) is used only for this upper bound.
\end{proof}

\begin{theorem}[The three theta exchanges on one field and interval]
\leavevmode\label{thm:thirdglobal843}
There is a unique \(\delta_3(c)\in\mathbb C^\times\) such that
\[
                F'_c(y)=\delta_3(c)F_c(y)
                    \qquad(y\in U(W)(\mathbb A)).
\]
Put \(d=m/2\), \(e=(p-r)/2\),
\(\sigma_c=\alpha^{-1}\upsilon_1^{n/2}\), and
\(b_{j,\vartheta}=\gamma_{12}(c)\delta_3(c)/N_R\).
For every prescribed choice \(i\) and every
\(g\in U(W_j)(\mathbb A)\) one has
\[
             b_{j,\vartheta} L_{i,\vartheta}(g)=\mathcal B_i(g),
\]
with
\[
 |b_{j,\vartheta}|^2=
 \frac{K_{12}K_3}
 {N_R^2\,\mathcal D_d(c\upsilon_1^d)\,
          \mathcal D_m(c^{-1}\upsilon_1^m)\,
          \mathcal D_e(\sigma_c)} .
\]
Here \(K_{12},K_3>0\) and \(N_R\) are independent of the prescribed
twist at the fixed node and lift index.
\end{theorem}
\begin{proof}
Both standard functions have nonzero norm. They are in the discrete
spectrum and every nonzero irreducible projection is in a copy of
\(\pi\), on the specified tensor line by
Theorem~\ref{thm:thirdadelic842}. The proposition shows that this
copy occurs exactly once. Thus the two nonzero vectors span the
same line in \(L^2([U(W)])\). Their smooth representatives satisfy
the equality everywhere because a nonempty open set has positive
quotient measure. The scalar is unique. Taking squared norms and
using Corollary~\ref{cor:thirdnormratio840} gives
\(\lvert\delta_3(c)\rvert^2=K_3/\mathcal D_e(\sigma_c)\).
Corollary~\ref{cor:thirdseed838} extends the identity from the
standard functions to every prescribed \(i\) and \(g\).
Theorem~\ref{thm:secondglobal833} gives the other two factors in the
denominator. All fields, residue levels, interval endpoints, tests
and measures used in the period identity are those fixed before
the introduction of \(H\).
\end{proof}
\clearpage

\clearpage
\section{References}\label{nav:bibliography-XV}
\begingroup\def\section*#1{}\begin{thebibliography}{99}\small
\bibitem{XV-linked-1} \href{https://www.claymath.org/library/cw/arthur/pdf/61.pdf}{Arthur, \emph{An introduction to the trace formula}, Theorem~8.1}.
\bibitem{XV-linked-2} \href{https://arxiv.org/pdf/2410.13504}{Atobe, Gan, Ichino, Kaletha, M\'inguez and Shin, Theorem~5.4.1}.
\bibitem{XV-linked-3} \href{https://doi.org/10.1007/BF01445125}{Banaszczyk, \emph{New bounds in some transference theorems in the geometry of numbers}}.
\bibitem{XV-linked-4} \href{https://arxiv.org/pdf/2103.07956}{Chen and Zou, Theorem~3.1, after J.-S.~Li, Theorem~A}.
\bibitem{XV-linked-5} \href{https://msp.org/pjm/1987/129-1/pjm-v129-n1-p02-s.pdf}{Davidson, \S2, Theorem~2.7}.
\bibitem{XV-linked-6} \href{https://arxiv.org/pdf/1207.4709}{Gan, Qiu and Takeda, \S\S11.5--11.6}.
\bibitem{XV-linked-7} \href{https://arxiv.org/abs/2611.26750}{Halleck-Dub\'e's November 2026 preprint, \S1.14}.
\bibitem{XV-linked-8} \href{https://doi.org/10.1515/crll.1992.428.177}{J.-S.~Li, Theorem~3.1}.
\bibitem{XV-linked-9} \href{https://webusers.imj-prg.fr/~colette.moeglin/MR7.pdf}{M{\oe}glin and Renard, Theorem~1.3}.
\bibitem{XV-linked-10} \href{https://arxiv.org/pdf/1206.0882}{Mok's quasi-split multiplicity formula, Theorem~2.5.2}.
\bibitem{XV-linked-11} \href{https://msp.org/pjm/2013/264-1/pjm-v264-n1-p11-s.pdf}{S.~Yamana, \emph{The Siegel--Weil formula for unitary groups}, Theorem~2}.
\bibitem{XV-linked-12} \href{https://www.mathnet.ru/eng/im2220}{V.~P.~Platonov, \emph{The problem of strong approximation and the Kneser--Tits conjecture for algebraic groups}}.
\bibitem{XV-linked-13} \href{https://www.numdam.org/article/ASENS_1980_4_13_2_165_0.pdf}{Zelevinsky, \emph{Induced representations of reductive \(p\)-adic groups II}, Theorem~4.2}.
\end{thebibliography}\endgroup
\end{document}
